\chapter{Obstructions}
\label{prov:ch:obs}
\label{obs:ch}

The chapter collects the negative results: what a smaller numerical basis or a
smaller rule inventory cannot do, together with the model or invariant that
proves it. Each result names the exact class of theories or presentations that
it excludes and states what it does not exclude. The results are grouped by
method. \Cref{obs:sec:powers} uses coordinatewise powers of the standard model
and exact truth transfer for conjunctive formulas. \Cref{obs:sec:symmetry} uses
the automorphisms of such powers. \Cref{obs:sec:majorants,obs:sec:quotient} use
growth invariants of term languages and a dimension argument for term
interpretations. \Cref{obs:sec:profiles} counts definable values and reaches
arbitrary relational bi-interpretations through witness extraction.
\Cref{obs:sec:zeroprofile} classifies zero sets. \Cref{obs:sec:rules} collects
small counter-interpretations of individual rule deletions.
\Cref{obs:sec:progress-independence} builds a nonstandard model of all iterator
computations, by a tower of ultrapowers, in which the progress rule fails, and a
refinement in which preservation of fixed points and seed shift hold as well;
it ends with two covers of a nonstandard model of arithmetic which show that
neither iterator fusion nor the successor numbering law suffices alone, even
with induction against fixed numerals. The packages that
prove the obstructions also contain positive statements about the presentations
that the obstructions are measured against; these are
in \cref{obs:sec:progressive,obs:sec:cc,obs:sec:sources,obs:sec:partial}.

None of the results is a lower bound on the size of a proof checker, and none
excludes equiconsistency. \Cref{obs:tab:summary} lists the main statements.

\begin{longtable}{@{}p{0.23\linewidth}p{0.36\linewidth}p{0.33\linewidth}@{}}
\caption{The obstructions of the chapter.}\label{obs:tab:summary}\\
\toprule
Result & Excluded & Not excluded\\
\midrule
\endfirsthead
\toprule
Result & Excluded & Not excluded\\
\midrule
\endhead
\bottomrule
\endfoot
\Cref{obs:thm:direct-target} & case introduction with target $z$ in place of
$\Suc z$, with all other native rules & other readings of the weaker calculus\\
\Cref{obs:thm:boolcases} & deletion of Boolean cases from the progressive
calculus & a weaker case principle\\
\Cref{obs:thm:symmetry} & ordinary bi-interpretation with $\iPRA$ of any theory
true in an eventually constant power: equation induction, conjunctive
induction, all true conjunctive sequents, many sorts & one-way
interpretations, equiconsistency, disjunctive axioms\\
\Cref{obs:prop:monotone} & a term compiler into the zero-seeded recursor &
relational interpretations\\
\Cref{obs:thm:adaptive}, with
\cref{obs:thm:single-body,obs:thm:circuit} & terms for addition, selection and
pairing when seed and update are finite adaptive query programs on one body &
an independent input, retained registers, relational interpretations\\
\Cref{obs:thm:quotient} & finite-dimensional interpretations with a
term-defined addition, any quotient equality & graph-defined operations\\
\Cref{obs:thm:fc} & terms for an equality test, minimum and truncated
subtraction in the tied-seed language & formula-defined relations\\
\Cref{obs:thm:no-biinterp} & every finite-dimensional bi-interpretation of
$\iPRA$ with an existential-free theory true in the standard tied-seed
algebra & one-way interpretations, theories with positive existence axioms\\
\Cref{obs:thm:zeroprofile} & terms for predecessor, zero test, addition and
initialized iteration from positive iteration & graph definitions\\
\Cref{obs:prop:forced-progress,obs:prop:majors,obs:prop:root-vs-congruence,obs:thm:RL-independent,obs:prop:reversed}
& five individual rule deletions or replacements & redesigned calculi with
their own interpretation proofs\\
\Cref{obs:thm:progress-independence,obs:cor:pi-extensional} & deletion of
\rulename{QI} by progress from the core: the other seven rules do not derive
$\Q{y}(\It{y}{0}{u}{u}=0)$, nor does classical extensional computation with
coverage derive its induction instance & fusions of \rulename{QI} with another rule, weaker
induction interfaces, equivalence of uniqueness of iteration with induction\\
\Cref{obs:thm:fixed-points,obs:cor:fp-native} & replacement of equation
induction by preservation of fixed points and seed shift, with the canonical
successor iterator and cancellative addition; \rulename{QI} over the seven rules
with \rulename{FP} and \rulename{Shift} & general fusion
(\cref{obs:prop:fusion}), a derivable step over the identity scheme alone, the
uniqueness rule for the successor update\\
\Cref{obs:thm:numbering-barrier,obs:thm:fusion-barrier,obs:cor:qi-identity} &
recovery of equation induction from numeral induction with the identity law,
all numbering laws $\It{n}{\overline{k}}{u}{\Suc u}=\Suc^{k}n$ and
extensionality, or with \rulename{Fusion} without
$\It{n}{0}{u}{\Suc u}=n$; \rulename{QI} over the seven rules with the identity
scheme & \rulename{Fusion} together with $\It{n}{0}{u}{\Suc u}=n$, which gives
induction\\
\Cref{obs:cor:barrier-biinterp} & bi-interpretation with $\iPRA$ of the fusion
theory with finite parameters, and of the numbering theory with standard
parameters & nonstandard parameters for the numbering theory, one-way
interpretations, equiconsistency\\
\end{longtable}

\section{Conventions and the presentations under test}
\label{obs:sec:setup}

This section fixes the vocabulary shared by \pkg{adaptive\_seed},
\pkg{research}, \pkg{local\_scopes}, \pkg{checker\_and\_barriers},
\pkg{compact\_kernel}, \pkg{uniform\_predecessor}, \pkg{scope\_kernel},
\pkg{constructed\_computation} and \pkg{counter\_congruence}.

\subsection{Three kinds of exclusion}
\label{obs:sec:kinds}

The statements of the chapter have three strengths.

\emph{Term definability.} A reduced term language is given a standard
interpretation: every fixed term denotes a primitive recursive function on
$\mathbb{N}$, by structural recursion on the term followed by numerical
primitive recursion. No uniform evaluator of all descriptions is used. A
term-definability obstruction is a statement about these functions. It
excludes a direct compiler of the primitive recursive signature into terms with
their usual meanings. It does not exclude a definition of the missing function
by a formula.

\emph{Derivability in a specified rule system.} A rule system is shown to lose
a theorem of another one by a reading under which every rule of the first is
sound and the theorem is false. A classical structure is a one-node Kripke
model, so soundness in a classical structure is a sound test of intuitionistic
derivability. Such a result concerns the displayed rules with the identity
translation of formulas. It does not exclude a different translation.

\emph{Ordinary bi-interpretability.} This is the strongest kind.

\begin{definition}[Ordinary interpretations]
\label{obs:def:biinterp}
An \emph{ordinary interpretation} $\mathcal{I}$ of a theory $U$ in a theory $T$
has a dimension $d\ge 0$ and consists of formulas of $T$: a domain formula
$\delta(\vec x)$ in $d$ variables for each sort of $U$, an equality formula
$\varepsilon(\vec x,\vec y)$, and a graph formula for each function and
relation symbol of $U$. The formulas may contain finitely many parameters,
subject to a parameter formula. $T$ proves that $\varepsilon$ is an equivalence
on $\delta$ respected by the graphs, that the function graphs are total and
single-valued modulo $\varepsilon$, and the translation of every axiom of $U$.
The domain, the graphs and the quotient equality are arbitrary formulas.

A \emph{bi-interpretation} of $T$ and $U$ is a pair of ordinary interpretations
$\mathcal{I}$ of $U$ in $T$ and $\mathcal{J}$ of $T$ in $U$ such that $T$
proves that a formula $\eta$ defines an isomorphism of the identity
interpretation of $T$ onto $\mathcal{J}$ interpreted through $\mathcal{I}$, and
$U$ proves the corresponding statement for the other composite. The formula
$\eta$ of $T$ is the \emph{source-side return isomorphism}.
\end{definition}

The bridge of \cref{base:thm:bridge} is a bi-interpretation in this sense, of
dimension one, with identity domains, actual equality and return isomorphisms
with graph $x=y$. The obstructions of
\cref{obs:sec:symmetry,obs:sec:profiles} exclude all bi-interpretations in the
sense of \cref{obs:def:biinterp}, not only those of this special form. They
use only the source-side return isomorphism.

\subsection{The paired-introduction form of the native core}
\label{obs:sec:paired}

The packages of this chapter state the native calculus with a progress
introduction that concludes a paired formula:
\[
\infrule{\Gamma,\ E(t,z)\proves E(t,\Suc z)}{\Gamma\proves \D{x}{y}E(x,y)}
\quad\rulename{DI-progress},
\qquad E \text{ a literal equation},\ z \text{ fresh for } \Gamma,t.
\]
The core $\Nat$ of \cref{base:sec:kernel} has \rulename{QI} by progress
instead. The \emph{paired-introduction form} $\Nat_{\mathrm{p}}$ is $\Nat$ with
\rulename{QI} replaced by \rulename{DI-progress}. Several packages also supply
a computation instance as a discharged local assumption (rule \rulename{Comp},
called \rulename{Supply} in some packages): from $\Gamma,H\proves A$ with $H$
an instance of a computation equation infer $\Gamma\proves A$.

\begin{lemma}
\label{obs:lem:paired}
$\Nat$ and $\Nat_{\mathrm{p}}$ derive the same sequents. \rulename{Comp} is
admissible in both, and the computation equations are derivable from
\rulename{Comp} and \rulename{Hyp}.
\end{lemma}

\begin{proof}
Let $\Gamma,E(t,z)\proves E(t,\Suc z)$ with $z$ fresh. In $\Nat$,
\rulename{QI} for the matrix $y\mapsto E(t,y)$ gives
$\Gamma\proves\Q{y}E(t,y)$. Weaken by $E(t,0)$. \rulename{Hyp} gives $E(t,0)$,
and \rulename{QE} at the target $\Suc z'$ gives
$\Gamma,E(t,0)\proves E(t,\Suc z')$. \rulename{DI} by cases with witness $t$
gives $\Gamma\proves\D{x}{y}E(x,y)$. Conversely, let
$\Gamma,E(z)\proves E(\Suc z)$. In $\Nat_{\mathrm{p}}$, \rulename{DI-progress}
with a witness coordinate $u$ that does not occur in $E$ gives
$\Gamma\proves\D{u}{y}E(y)$. \rulename{DE} supplies the clause $\Q{y}E(y)$,
which does not contain the fresh witness, and \rulename{Hyp} concludes
$\Q{y}E(y)$. \rulename{Comp} is a cut with a computation equation
(\cref{base:sec:structural}); conversely \rulename{Comp} followed by
\rulename{Hyp} derives the equation.
\end{proof}

In the paired-introduction form the formula $\Q{y}B$ is interderivable with
$\D{u}{y}B$ for an unused $u$: \rulename{DE} followed by \rulename{Hyp} gives
one direction, and \rulename{DI} by cases at the witness $0$ with \rulename{QE}
gives the other. The rule-level statements of \cref{obs:sec:rules} refer to
$\Nat_{\mathrm{p}}$ where the difference matters.

\subsection{Reduced term languages}
\label{obs:sec:languages}

\begin{definition}[Seedless recursors]
\label{obs:def:languages}
Each of the following languages has variables, $0$ and one recursion
constructor. Only the update body binds the displayed bound variables; its
other variables are passive parameters, fixed through the recursion. Bodies are
fixed finite earlier terms.
\begin{enumerate}[label=(\alph*)]
\item \emph{Zero-seeded recursion} has, besides $0$ and $\Suc$, the
constructor $\mathsf{L}(n;u,v.h)$ with
\[
\mathsf{L}(0;u,v.h)=0,\qquad
\mathsf{L}(\Suc n;u,v.h)=h\subst{n}{u}\subst{\mathsf{L}(n;u,v.h)}{v}.
\]
\item \emph{Tied-seed recursion} has the constructor $\mathsf{B}(n;u,v.h)$ and
the abbreviation $\Suc a\defeq\mathsf{B}(0;u,v.a)$ for $u,v$ not in $a$, with
\begin{align}
\mathsf{B}(0;u,v.h)&=\Suc\bigl(h\subst{0}{u}\subst{0}{v}\bigr),\label{obs:eq:B0}\\
\mathsf{B}(\Suc n;u,v.h)&=h\subst{n}{u}\subst{\mathsf{B}(n;u,v.h)}{v}.\label{obs:eq:BS}
\end{align}
The two substitutions are simultaneous. The initial state is tied to the
update; there is no independent initializer. The standard interpretation is
\[
\mathsf{B}_h(0,\vec p)=1+h(0,0,\vec p),\qquad
\mathsf{B}_h(m+1,\vec p)=h(m,\mathsf{B}_h(m,\vec p),\vec p),
\]
so that $\Suc$ denotes the successor. The \emph{unary-update} variant
$\mathsf{F}(n;v.h)$ is the case in which $u$ does not occur in $h$.
\item \emph{Positive iteration} has, besides $0$ and $\Suc$, the constructor
$\mathsf{R}^{+}(n,a;v.h)$, standing for $h^{n+1}(a)$:
\[
\mathsf{R}^{+}(0,a;v.h)=h\subst{a}{v},\qquad
\mathsf{R}^{+}(\Suc n,a;v.h)=h\subst{\mathsf{R}^{+}(n,a;v.h)}{v}.
\]
\end{enumerate}
For reference, the \emph{counter-aware initialized recursor} is
$\mathsf{R}(n,a;u,v.h)$ with $\mathsf{R}(0,a;u,v.h)=a$ and
$\mathsf{R}(\Suc n,a;u,v.h)=h\subst{n}{u}\subst{\mathsf{R}(n,a;u,v.h)}{v}$. It
is ordinary primitive recursion and compiles every primitive recursive
description directly.
\end{definition}

The tied-seed language is the most capable of the three. The next two lemmas
record what it does define; they are the positive boundary of the obstructions
of \cref{obs:sec:majorants,obs:sec:profiles}.

\begin{lemma}[Terms of the tied-seed language]
\label{obs:lem:Bterms}
Put
\[
W(n)\defeq\mathsf{B}(n;u,v.u),\quad
Z(n)\defeq\mathsf{B}(n;u,v.0),\quad
H(n,a)\defeq\mathsf{B}(n;u,v.a),\quad
P(n)\defeq W(H(n,n)).
\]
The computation equations \eqref{obs:eq:B0}, \eqref{obs:eq:BS} and equality
reasoning prove
\[
W(0)=1,\ W(\Suc n)=n,\quad Z(0)=1,\ Z(\Suc n)=0,\quad
H(0,a)=\Suc a,\ H(\Suc n,a)=a,\quad P(0)=0,\ P(\Suc n)=n .
\]
With induction for literal equations, the two gates
\[
G(n,a)\defeq P\bigl(P(\mathsf{B}(a;u,v.H(Z(n),v)))\bigr),\qquad
G'(n,a)\defeq P\bigl(P(\mathsf{B}(a;u,v.H(n,v)))\bigr)
\]
satisfy $G(0,a)=0$, $G(\Suc q,a)=a$ and $G'(0,a)=a$, $G'(\Suc q,a)=0$.
\end{lemma}

\begin{proof}
The equations for $W$, $Z$ and $H$ are instances of \eqref{obs:eq:B0} and
\eqref{obs:eq:BS}. Then $P(0)=W(H(0,0))=W(\Suc 0)=0$ and
$P(\Suc n)=W(H(\Suc n,\Suc n))=W(\Suc n)=n$. For $G$ at the control $0$, the
update is $H(Z(0),v)=H(1,v)=v$. Induction on the literal equation
$\mathsf{B}(a;u,v.H(Z(0),v))=1$ has the base $\Suc(H(1,0))=\Suc 0$ and a
successor step that repeats the previous state. Two predecessors give
$G(0,a)=0$. At the control $\Suc q$ the update is $H(0,v)=\Suc v$. Induction on
$\mathsf{B}(a;u,v.H(Z(\Suc q),v))=\Suc\Suc a$ has the base
$\Suc(H(0,0))=\Suc\Suc 0$ and a successor step that applies $\Suc$. Two
predecessors give $G(\Suc q,a)=a$. The argument for $G'$ exchanges the two
cases. The inductions keep the actual update bodies and the control parameter
fixed; no extensionality principle for recursion descriptions and no induction
on a compound formula is used.
\end{proof}

\begin{lemma}[Zero-start counter recursion is a macro]
\label{obs:lem:zerostart}
For a body $h$ with bound variables $u,v$ and fresh $r,s$, put
\[
L(n;u,v.h)\defeq
\mathsf{B}\bigl(\Suc n;\,r,s.\,G(r,\,h\subst{P r}{u}\subst{s}{v})\bigr).
\]
Then $L(0;u,v.h)=0$ and
$L(\Suc n;u,v.h)=h\subst{n}{u}\subst{L(n;u,v.h)}{v}$ follow from
\cref{obs:lem:Bterms} without further induction.
\end{lemma}

\begin{proof}
Let $K(r,s)=G(r,h\subst{P r}{u}\subst{s}{v})$. The gate equations and
$P(\Suc r)=r$ give $K(0,s)=0$ and $K(\Suc r,s)=h\subst{r}{u}\subst{s}{v}$.
Hence $\mathsf{B}_K(0)=\Suc(K(0,0))=1$ and
$\mathsf{B}_K(1)=K(0,\mathsf{B}_K(0))=0$, which is $L(0)=0$. At a successor,
$L(\Suc n)=\mathsf{B}_K(\Suc\Suc n)=K(\Suc n,\mathsf{B}_K(\Suc n))
=h\subst{n}{u}\subst{L(n)}{v}$.
\end{proof}

\begin{remark}
\Cref{obs:lem:Bterms} is in \pkg{counter\_congruence} (with the gate $G'$),
\pkg{research} (with $P(n)=W(H(Z(n),W(n)))$, which computes the same function)
and \pkg{local\_scopes} (with the gate $G$ and the internal proof).
\Cref{obs:lem:zerostart} is in \pkg{local\_scopes}. Predecessor, a zero test, a
one-sided gate and zero-start recursion are therefore not what the tied-seed
language lacks. The obstruction is to combining two unbounded inputs.
\end{remark}

\subsection{The progressive calculus}
\label{obs:sec:Hrules}

\begin{definition}[Progressive binder calculus]
\label{obs:def:progressive}
Numerical terms are variables, $0$, $\Suc t$ and the counter-aware recursor
$\mathsf{R}(n,a;u,v.h)$ with its two computation equations. Object formulas are
\[
A::=(s=t)\mid\mathsf{H}_{x,y}A .
\]
Internal clauses are $\mathsf{P}_{y}A(y)$ with $A$ an object formula. Contexts
are finite reusable lists of object formulas and clauses. Write
$\Gamma\mathrel{\triangleright}M$ when $M$ is a member of $\Gamma$ or a proved
object formula; clauses are not goals. The ordinary readings are
\[
(\mathsf{P}_{y}B)^{\circ}=\forall y\,\bigl(B^{\circ}(y)\to B^{\circ}(\Suc y)\bigr),
\qquad
(\mathsf{H}_{x,y}B)^{\circ}=\exists x\,\forall y\,\bigl(B^{\circ}(x,y)\to B^{\circ}(x,\Suc y)\bigr).
\]
Besides the computation equations, the rules are
\begin{gather*}
\infrule{\Gamma\mathrel{\triangleright}s=t\quad\Gamma\proves A(s)}{\Gamma\proves A(t)}\ \rulename{Tr}
\qquad
\infrule{\Gamma,\ B(t,z)\proves B(t,\Suc z)}{\Gamma\proves\mathsf{H}_{x,y}B(x,y)}\ \rulename{HI}
\\[4pt]
\infrule{\Gamma\mathrel{\triangleright}\mathsf{H}_{x,y}B(x,y)\quad\Gamma,\ \mathsf{P}_{y}B(a,y)\proves A}{\Gamma\proves A}\ \rulename{HE}
\qquad
\infrule{\Gamma\mathrel{\triangleright}\mathsf{P}_{y}B(y)\quad\Gamma\proves B(t)}{\Gamma\proves B(\Suc t)}\ \rulename{Step}
\\[4pt]
\infrule{\Gamma\mathrel{\triangleright}\mathsf{P}_{y}E(y)\quad\Gamma\proves E(0)}{\Gamma\proves E(t)}\ \rulename{Iter}
\qquad
\infrule{\Gamma,\ \operatorname{sg}(t)=0\proves A\quad\Gamma,\ \operatorname{sg}(t)=1\proves A}{\Gamma\proves A}\ \rulename{BoolCases}
\end{gather*}
In \rulename{Iter} the matrix $E$ is a literal equation. In \rulename{HI} the
variable $z$ is fresh for $\Gamma$, $t$ and the passive parameters; in
\rulename{HE} the variable $a$ is fresh for $\Gamma$, $A$ and the passive
parameters. Here $1=\Suc 0$ and $\operatorname{sg}(n)=\mathsf{R}(n,0;u,v.1)$,
so $\operatorname{sg}(0)=0$ and $\operatorname{sg}(\Suc n)=1$. There is no
primitive hypothesis rule and no other induction rule. The \emph{case-free
progressive calculus} omits \rulename{BoolCases}.
\end{definition}

\rulename{HI} is an unrestricted logical introduction; the only induction is
\rulename{Iter}, at literal equations, in an elimination. The calculus with
\rulename{BoolCases} is a complete presentation of $\iPRA$
(\cref{obs:thm:progressive}); \cref{obs:thm:boolcases} shows that
\rulename{BoolCases} cannot be deleted.

\section{Coordinatewise powers and truth transfer}
\label{obs:sec:powers}

This section draws on \pkg{research}, \pkg{compact\_kernel},
\pkg{scope\_kernel} and \pkg{checker\_and\_barriers}. One family of structures
serves four results: equation induction does not imply numerical cases
(\cref{obs:thm:direct-target}), Boolean cases are not admissible in the
progressive calculus (\cref{obs:thm:boolcases}), and, in
\cref{obs:sec:symmetry}, positive axioms and conjunctive induction do not give
a theory bi-interpretable with $\iPRA$.

\subsection{The powers}

\begin{definition}[Coordinatewise powers]
\label{obs:def:power}
Let $\mathcal{N}$ be a finitary expansion of the standard structure
$(\mathbb{N},0,\Suc)$ by total functions and relations. The signature may be
infinite; in particular it may contain a symbol for every primitive recursive
description. $\mathcal{N}$ may be many-sorted, with nonempty sorts and with
$(\mathbb{N},0,\Suc)$ as a distinguished numerical sort. For a nonempty index
set $J$, the structure $\mathcal{N}^{(J)}$ has, in each sort, the families
$a=(a(i))_{i\in J}$ of elements of that sort that are constant outside a finite
subset of $J$. Functions are interpreted coordinatewise. A relation holds of
families exactly when it holds at every coordinate. Equality is actual equality
of families, and constants are constant families.

For $J=\{1,2\}$ the structure is the full product
$\mathcal{A}=\mathcal{N}\times\mathcal{N}$. For $J=\mathbb{N}$ it is the
structure $\mathcal{E}$ of eventually constant sequences. For $i\in J$, the
family $e_i$ of the numerical sort has the value $1$ at $i$ and $0$ elsewhere.
\end{definition}

Finitely many families are simultaneously constant outside the union of their
exceptional sets, so coordinatewise functions preserve the domain. Each element
of $\mathcal{E}$ has a finite description, so $\mathcal{E}$ is countable when
the sorts of $\mathcal{N}$ are.

\begin{definition}[Conjunctive formulas]
\label{obs:def:conjunctive}
The class $\mathcal{C}$ of \emph{conjunctive formulas} is generated from atomic
formulas and $\top$ by $\wedge$, $\exists$ and $\forall$, with unrestricted
quantifier alternation. A conjunctive formula without $\forall$ is
\emph{primitive positive}. No disjunction, implication or negation occurs
inside a conjunctive formula.
\end{definition}

\begin{theorem}[Conjunctive transfer]
\label{obs:thm:transfer}
For every $\varphi(\vec x)\in\mathcal{C}$, every nonempty $J$ and every tuple
$\vec a$ of $\mathcal{N}^{(J)}$,
\[
\mathcal{N}^{(J)}\models\varphi(\vec a)
\quad\Longleftrightarrow\quad
\mathcal{N}\models\varphi(\vec a(i))\ \text{ for every } i\in J .
\]
\end{theorem}

\begin{proof}
Induction on the formula. Terms evaluate coordinatewise, so atomic formulas
transfer by the definition of the relations and of equality. Conjunction is
immediate.

\emph{Existential quantifier.} If $\mathcal{N}^{(J)}\models\theta(b,\vec a)$,
the induction hypothesis gives $\mathcal{N}\models\theta(b(i),\vec a(i))$ at
every $i$. Conversely, suppose
$\mathcal{N}\models\exists y\,\theta(y,\vec a(i))$ at every $i$. There is a
finite set $F\subseteq J$ such that $\vec a(i)$ is one fixed tuple
$\vec a_\infty$ for all $i\notin F$. Choose a witness $b_i$ for each of the
finitely many $i\in F$ and one witness $b_\infty$ for $\vec a_\infty$ if
$J\ne F$. The family with the values $b_i$ on $F$ and $b_\infty$ elsewhere
belongs to $\mathcal{N}^{(J)}$, and the induction hypothesis gives
$\mathcal{N}^{(J)}\models\theta(b,\vec a)$. Only finitely many witnesses are
chosen.

\emph{Universal quantifier.} Suppose
$\mathcal{N}^{(J)}\models\forall y\,\theta(y,\vec a)$. Fix $i$ and an element
$b$ of the relevant sort of $\mathcal{N}$. The constant family $\overline{b}$
is an element of $\mathcal{N}^{(J)}$; instantiate at $\overline{b}$ and apply
the induction hypothesis at $i$ to get $\mathcal{N}\models\theta(b,\vec a(i))$.
Conversely, if every coordinate satisfies $\forall y\,\theta(y,\vec a(i))$,
then for every family $b$ every coordinate satisfies
$\theta(b(i),\vec a(i))$, and the induction hypothesis gives
$\mathcal{N}^{(J)}\models\theta(b,\vec a)$.
\end{proof}

\begin{corollary}[Regular axioms and conjunctive induction]
\label{obs:cor:power-axioms}
Let $J$ be nonempty.
\begin{enumerate}[label=(\roman*)]
\item $\mathcal{N}^{(J)}$ satisfies every sentence
$\forall\vec x\,(\varphi\to\psi)$ and every sentence
$\forall\vec x\,(\varphi\to\bot)$ with $\varphi,\psi\in\mathcal{C}$ that is
true in $\mathcal{N}$.
\item For every $A(n,\vec p)\in\mathcal{C}$, with $n$ of the numerical sort,
$\mathcal{N}^{(J)}$ satisfies
\[
\forall\vec p\,\Bigl[\bigl(A(0,\vec p)\wedge\forall n\,(A(n,\vec p)\to A(\Suc n,\vec p))\bigr)\to\forall n\,A(n,\vec p)\Bigr].
\]
\end{enumerate}
Consequently the contextual induction rule for matrices in $\mathcal{C}$ is
sound in $\mathcal{N}^{(J)}$ under arbitrary retained assumptions.
\end{corollary}

\begin{proof}
(i) If $\varphi(\vec a)$ holds in $\mathcal{N}^{(J)}$, then
$\varphi(\vec a(i))$ holds in $\mathcal{N}$ at every $i$ by
\cref{obs:thm:transfer}; truth of the sentence in $\mathcal{N}$ gives
$\psi(\vec a(i))$ at every $i$, and the converse direction of the theorem gives
$\psi(\vec a)$. For a negative constraint, a true antecedent in
$\mathcal{N}^{(J)}$ gives a true antecedent in $\mathcal{N}$ at any one
coordinate, which contradicts the constraint.

(ii) Fix $\vec p$ and suppose the base and the progress hypothesis hold in
$\mathcal{N}^{(J)}$. Fix $i\in J$ and put
$U_i=\{k:\mathcal{N}\models A(k,\vec p(i))\}$. Transfer of the base gives
$0\in U_j$ for every $j$. Let $k\in U_i$ and let $a$ be the family with the
value $k$ at $i$ and $0$ elsewhere. At $i$ the formula $A(a(i),\vec p(i))$
holds by the choice of $k$, and at every other coordinate it is the transferred
base; hence $\mathcal{N}^{(J)}\models A(a,\vec p)$. The progress hypothesis
gives $A(\Suc a,\vec p)$, and transfer at $i$ gives $k+1\in U_i$. Ordinary
induction in $\mathbb{N}$ gives $U_i=\mathbb{N}$ for every $i$, and the
converse direction of the theorem gives $A(a,\vec p)$ for every $a$.

For the contextual rule, fix an assignment that makes the retained assumptions
true. The induction variable is fresh, so changing its value does not change
their truth, and the argument applies under that assignment. The retained
assumptions may be arbitrary formulas. Witnesses for $A$ are not assumed to be
terms.
\end{proof}

In the primitive recursive signature the corollary covers every defining
equation, the injectivity of $\Suc$, the separation $0=\Suc x\to\bot$, every
universal Horn sentence
$\forall\vec x\,(E_1\wedge\dots\wedge E_r\to F)$ with equations $E_k$ and an
equation or $\bot$ as $F$ that is true in $\mathbb{N}$, and induction for
literal equations with parameters.

\begin{definition}
\label{obs:def:Talg}
$T_{\mathrm{alg}}$ is intuitionistic first-order logic with equality in the
signature with a symbol for every primitive recursive description, with all
defining equations, the separation axioms and induction for every literal
equation. Coverage is absent. For a structure $\mathcal{N}$ as in
\cref{obs:def:power}, $T_{\mathcal{C}}(\mathcal{N})$ is intuitionistic logic
with equality together with all sentences of
\cref{obs:cor:power-axioms}(i) that are true in $\mathcal{N}$ and all
induction sentences of \cref{obs:cor:power-axioms}(ii);
$T_{\mathrm{pp}}(\mathcal{N})$ is its restriction to primitive positive
$\varphi,\psi,A$.
\end{definition}

By \cref{obs:cor:power-axioms}, every $\mathcal{N}^{(J)}$ is a model of
$T_{\mathcal{C}}(\mathcal{N})$, hence of $T_{\mathrm{pp}}(\mathcal{N})$, and,
for the primitive recursive signature, of $T_{\mathrm{alg}}$ together with all
universal Horn sentences true in $\mathbb{N}$.

\begin{proposition}[What fails in a power]
\label{obs:prop:coverage-fails}
Let $J$ have at least two elements and let the signature contain
$\operatorname{sg}$. In $\mathcal{N}^{(J)}$:
\begin{enumerate}[label=(\roman*)]
\item $e_i$ is neither $0$ nor a successor, so coverage
$\forall x\,(x=0\vee\exists y\,x=\Suc y)$ fails;
\item the quantifier-free induction instance for
$\varphi(n)\defeq(\operatorname{sg}(n)=0\vee\operatorname{sg}(n)=1)$ fails:
$\varphi(0)$ and $\varphi(\Suc a)$ hold for every $a$, and $\varphi(e_i)$
fails;
\item no formula with finitely many parameters defines a linear order of the
numerical sort when $J$ is infinite (\cref{obs:lem:involution}).
\end{enumerate}
\end{proposition}

\begin{proof}
Every successor family is positive at every coordinate, and $e_i$ is $0$ at
some coordinate and $1$ at $i$. The family $\operatorname{sg}(\Suc a)$ is the
constant $1$, while $\operatorname{sg}(e_i)=e_i$ is neither the constant $0$
nor the constant $1$. For (iii), an order defined over a finite tuple of
parameters is preserved by the involution of \cref{obs:lem:involution} that
fixes the tuple, and an order-preserving involution of a linear order is the
identity (\cref{obs:lem:rigid}).
\end{proof}

The same holds for the relation $x\le y$ of the primitive recursive signature:
for $i\ne j$ neither $e_i\le e_j$ nor $e_j\le e_i$ holds, since each fails at
one coordinate.

\begin{remark}
\label{obs:rem:boolean}
The missing feature is a disjunctive range assertion, not numerical
computation. A primitive recursive characteristic of equality still reflects
equality in a power: its value is the zero family exactly when its arguments
agree at every coordinate. What fails is the assertion that the value is
globally $0$ or globally $1$. In the same way, equation induction proves the
algebraic identities $t(\sel(n,a,b))=\sel(n,t(a),t(b))$ for each fixed term
context $t$, and $\sel(b,0,1)=b$ for the Boolean value $b=\sel(n,0,1)$; these
hold for the mixed values $(0,1)$ and $(1,0)$ of $\mathcal{A}$ and do not
entail $b=0\vee b=1$. Coordinatewise predecessor satisfies $P(0)=0$ and
$P(\Suc x)=x$ in every power, so adding valid predecessor equations does not
restore coverage.
\end{remark}

\begin{remark}
The product $\mathcal{A}$ with literal-equation induction is in
\pkg{scope\_kernel} and \pkg{checker\_and\_barriers}; \pkg{research}
attributes the product model of non-coverage to an earlier draft. The structure
$\mathcal{E}$ with induction for equations and for finite conjunctions of
atoms, universal Horn sentences and arbitrary finitary expansions is in
\pkg{research}. Primitive positive transfer, the universal quantifier and the
many-sorted form are in \pkg{compact\_kernel}. The statement for an arbitrary
index set is the same proof. % ed.: the packages treat J = {1,2} and J = N separately; the proof of the existential case only needs finitely many exceptional coordinates.
\end{remark}

\subsection{Equation induction does not replace numerical cases}

Consider the \emph{direct-target} replacement of case introduction:
\[
\infrule{\Gamma,\ B(t,0)\proves B(t,z)}{\Gamma\proves\D{x}{y}B(x,y)}
\qquad z \text{ fresh for } \Gamma,t \text{ and the parameters of } B .
\]
It is sound for the ordinary reading of $\D{x}{y}$ in every structure, with no
coverage axiom. Let $\Nat^{\mathrm{dt}}$ be $\Nat$ (or $\Nat_{\mathrm{p}}$)
with \rulename{DI} by cases replaced by this rule.

\begin{theorem}[\pkg{scope\_kernel}]
\label{obs:thm:direct-target}
Every theorem of $\Nat^{\mathrm{dt}}$ is valid in every power
$\mathcal{N}^{(J)}$ of the standard model of the iterator language, under the
ordinary reading. Write $1=\Suc 0$ and define
\begin{align*}
E(n,y,p)&\defeq\bigl(\sel(p,\,\sel(n,y,1),\,\sel(n,1,0))=0\bigr), &
K(n,p)&\defeq\Q{y}E(n,y,p),\\
M(p)&\defeq\Q{n}K(n,p), &
F&\defeq\Q{p}M(p).
\end{align*}
The closed formula $F$ is derivable in $\Nat$ without the progress rule and is
false in $\mathcal{A}$. Hence $\Nat^{\mathrm{dt}}$ does not derive $F$.
\end{theorem}

\begin{proof}
\emph{Soundness.} Interpret $\Q{y}B$ and $\D{x}{y}B$ by their ordinary
readings. \rulename{Hyp}, transport, \rulename{QE} and \rulename{DE} are sound
in every structure with actual equality. The computation equations hold
coordinatewise. The progress rule is sound by
\cref{obs:cor:power-axioms}(ii). The direct-target rule is sound because its
fresh variable already ranges over the whole domain.

\emph{Derivation of $F$ in $\Nat$.} A closure $\Q{y}B(y)$ follows from a
derivation of $B(\Suc z)$ with $z$ fresh: \rulename{DI} by cases with an unused
witness gives $\D{u}{y}B(y)$, and \rulename{DE} with \rulename{Hyp} gives
$\Q{y}B(y)$. At $n=\Suc z$ the computation equations give
$\sel(\Suc z,y,1)=1$ and $\sel(\Suc z,1,0)=0$, so $E(\Suc z,y,p)$ is
interderivable with $\sel(p,1,0)=0$, independently of $y$. Hence
$E(\Suc z,0,p)\proves E(\Suc z,\Suc w,p)$ by transport, and $K(\Suc z,p)$ is
derivable. This is the successor instance of $M(p)$, so $M(p)$ is derivable
for a free $p$; its instance $M(\Suc w)$ is the successor instance of $F$.
Only case introduction, \rulename{DE}, \rulename{QE}, the computations and
transport are used.

\emph{$F$ is false in $\mathcal{A}$.} Let $p_*=(0,1)$. Coordinatewise
computation gives
\[
\mathcal{A}\models E(n,y,p_*)
\quad\Longleftrightarrow\quad n_1=0,\ n_2>0,\ y_1=0 .
\]
Hence $K(0,p_*)$ holds, since its antecedent $E(0,0,p_*)$ is false, and
$K((0,1),p_*)$ fails: its antecedent holds at $y=0$ and its consequent fails at
$y=(1,0)$. So $M(p_*)$ is false. At $p=0$ the formula $E(0,y,0)$ is $y=0$, so
$K(0,0)$, which reads $0=0\to\forall y\,(y=0)$, is false and $M(0)$ is true.
Therefore $F$, which reads $M(0)\to\forall p\,M(p)$, is false.
\end{proof}

The theorem excludes this replacement with the identity translation of
formulas. It does not show that every interpretation of the weaker calculus in
$\iPRA$ fails, and it does not show that no restricted form of numerical cases
suffices.

\subsection{Boolean cases are not admissible in the progressive calculus}

\begin{theorem}[\pkg{checker\_and\_barriers}]
\label{obs:thm:boolcases}
Let $C$ and $Z$ be the terms $C(n,a,b)=\mathsf{R}(n,a;u,v.b)$ and
$Z(n)=\mathsf{R}(n,1;u,v.0)$, so that $C(0,a,b)=a$, $C(\Suc n,a,b)=b$,
$Z(0)=1$ and $Z(\Suc n)=0$. For a numerical term $t$ put
\[
F(t)\defeq\mathsf{H}_{u,y}\,\mathsf{H}_{v,z}\bigl(C(t,Z(y),z)=0\bigr),
\]
with $u,v$ unused. The case-free progressive calculus of
\cref{obs:def:progressive} derives $F(0)$ and $F(1)$ and does not derive
$F(\operatorname{sg}(t))$. Hence \rulename{BoolCases}, with an arbitrary object
formula as conclusion, is not admissible in the case-free calculus, in the
presence of all primitive recursive symbols, their defining equations and
contextual induction for literal equations.
\end{theorem}

\begin{proof}
\emph{The endpoints.} The derivations use the derived identity
$\Gamma,A\proves A$ of \cref{obs:lem:Hidentity}, which needs no cases. At
$t=0$ the inner matrix is interderivable with $Z(y)=0$, independently of $z$.
Its progress step is therefore an identity, and \rulename{HI} gives
$\mathsf{H}_{v,z}(C(0,Z(y),z)=0)$ for arbitrary $y$. The outer \rulename{HI}
ignores its local assumption. At $t=1$ the inner matrix is interderivable with
$z=0$, independently of $y$. For the outer \rulename{HI}, open the assumed
inner formula at $y$ by \rulename{HE}, introduce the inner formula at $\Suc y$
by \rulename{HI} with the same witness, convert its local assumption to the
matrix at $y$ by the computation equations, apply \rulename{Step}, and convert
back.

\emph{The model.} Let $\mathcal{A}$ be the product of the standard model of the
primitive recursive signature. Under the readings of
\cref{obs:def:progressive}, \rulename{Tr}, \rulename{HI}, \rulename{HE} and
\rulename{Step} are logical rules, valid in every structure with actual
equality, and \rulename{Iter} is induction for a literal equation, valid by
\cref{obs:cor:power-axioms}(ii). The computation equations hold. So every
theorem of the case-free calculus is valid in $\mathcal{A}$.

Let $e=(0,1)$, so $\operatorname{sg}(e)=e$. For $y=(y_1,y_2)$ and
$z=(z_1,z_2)$ the inner equation $C(e,Z(y),z)=0$ says
$(Z(y_1),z_2)=(0,0)$. If $y_1=0$ it is false for every $z$, so the inner
progress statement is vacuously true. If $y_1>0$ it says $z_2=0$; it holds at
$z=(0,0)$ and fails at $\Suc z=(1,1)$, so the inner statement is false. Thus
the inner formula is true at $y=(0,0)$ and false at $\Suc y=(1,1)$, which
refutes the outer progress statement for every $u$. Hence
$\mathcal{A}\not\models F(e)$, and $F(\operatorname{sg}(t))$ is not derivable.

\emph{Non-admissibility.} From the assumption $\operatorname{sg}(t)=0$,
transport of $F(0)$ gives $F(\operatorname{sg}(t))$; from
$\operatorname{sg}(t)=1$, transport of $F(1)$ gives the same. An admissible
\rulename{BoolCases} would derive $F(\operatorname{sg}(t))$.
\end{proof}

\begin{proposition}[The binder depth is sharp for this separation]
\label{obs:prop:depth}
Let $G(t)$ be an equation or a formula $\mathsf{H}_{x,y}E(x,y,t)$ with $E$ a
literal equation. If $G(0)$ and $G(1)$ hold in the standard model, then
$G((0,1))$ holds in $\mathcal{A}$.
\end{proposition}

\begin{proof}
For an equation this is coordinatewise evaluation. Otherwise let $a_0$ and
$a_1$ be standard witnesses for $t=0$ and $t=1$. If
$E((a_0,a_1),(y_1,y_2),(0,1))$ holds in $\mathcal{A}$, then $E(a_0,y_1,0)$ and
$E(a_1,y_2,1)$ hold in $\mathbb{N}$, so $E(a_0,y_1+1,0)$ and $E(a_1,y_2+1,1)$
hold, which is $E$ at $\Suc y$ in $\mathcal{A}$.
\end{proof}

The first failure arises when an implication tests a compound premise: one
coordinate makes an inner implication vacuous while the other would fail it.
\Cref{obs:thm:boolcases} rejects the deletion of \rulename{BoolCases}. It is
not a statement about other logical bases, other primitive operations or other
interpretations; the stronger statement about bi-interpretations of the
case-free calculus is \cref{obs:cor:casefree}.

\section{Symmetry}
\label{obs:sec:symmetry}

This section draws on \pkg{research} and \pkg{compact\_kernel};
\pkg{checker\_and\_barriers} applies the result to the case-free progressive
calculus. \pkg{research} attributes the method, a finite-parameter involution
against a definable linear order, to earlier drafts on unary arithmetic. The
obstruction of \cref{obs:prop:coverage-fails} concerns the original domain. An
interpretation on a different definable domain is not touched by it; the
symmetry argument handles that freedom.

\begin{lemma}[Finite parameters leave an involution]
\label{obs:lem:involution}
Let $J$ be infinite. For every finite tuple $\vec c$ of
$\mathcal{N}^{(J)}$, the expansion $(\mathcal{N}^{(J)},\vec c)$ has an
automorphism of order two that is not the identity on the numerical sort.
\end{lemma}

\begin{proof}
Outside a finite set $F\subseteq J$ every component of $\vec c$ is constant.
Choose distinct $i,j\notin F$ and exchange the coordinates $i$ and $j$ in every
family of every sort. The exchange preserves the domain, commutes with every
coordinatewise function and preserves every coordinatewise relation, whatever
the size of the signature. It fixes $\vec c$, has order two, and moves $e_i$ to
$e_j$.
\end{proof}

\begin{lemma}[Order rigidity]
\label{obs:lem:rigid}
Every classical model of $\iPRA$ has a parameter-free definable strict linear
order, for example
$x<y\iffdef\Suc x\mathbin{\dot{-}}y=0$ with truncated subtraction. An
automorphism of finite order of a linearly ordered structure that preserves the
order is the identity.
\end{lemma}

\begin{proof}
The order laws are theorems of primitive recursive arithmetic with
quantifier-free induction. If $a<\sigma(a)$, repeated application gives
$a<\sigma(a)<\dots<\sigma^{r}(a)=a$ for the order $r$ of $\sigma$; the case
$\sigma(a)<a$ is symmetric. Linearity gives $\sigma(a)=a$.
\end{proof}

\begin{lemma}[Automorphisms under a bi-interpretation]
\label{obs:lem:aut}
Let a structure $M$ interpret a structure $N$ by an ordinary interpretation
with parameters among $\vec c$. Every automorphism of $M$ that fixes $\vec c$
acts coordinatewise on representing tuples and induces an automorphism of $N$;
this is a group homomorphism
$\operatorname{Aut}(M/\vec c)\to\operatorname{Aut}(N)$. If in addition $N$
interprets a copy of $M$, and an isomorphism $\eta$ of $M$ onto that copy is
definable in $M$, and $\vec c$ contains the parameters of $\eta$ and
representatives in $M$ of the parameters of the second interpretation, then the
homomorphism is injective.
\end{lemma}

\begin{proof}
The domain, the equality and the graphs of the interpretation are definable
over $\vec c$, so they are preserved by $\sigma\in\operatorname{Aut}(M/\vec c)$;
this gives the induced automorphism $\widehat\sigma$ and the homomorphism. The
automorphism $\widehat\sigma$ fixes the parameters of the second
interpretation, since $\sigma$ fixes their representatives, and so induces an
automorphism $\widehat{\widehat\sigma}$ of the copy of $M$. Since $\eta$ is
definable over parameters fixed by $\sigma$,
$\eta(\sigma(a))=\widehat{\widehat\sigma}(\eta(a))$. If $\widehat\sigma$ is the
identity, so is $\widehat{\widehat\sigma}$, hence $\eta(\sigma(a))=\eta(a)$,
and injectivity of $\eta$ gives $\sigma(a)=a$. Only this one composite is used.
The interpretations may have any finite dimension, quotient equality and
domains different from the original ones.
\end{proof}

\begin{theorem}[Symmetry obstruction]
\label{obs:thm:symmetry}
Let $J$ be infinite and let $T$ be a theory, in the language of a structure
$\mathcal{N}$ as in \cref{obs:def:power}, that has $\mathcal{N}^{(J)}$ as a
model. Then $T$ and $\iPRA$ are not bi-interpretable in the sense of
\cref{obs:def:biinterp}. Arbitrary formula-defined domains, graphs and quotient
equality are allowed, in any finite dimension and over finite products of the
sorts. Finite parameter tuples are allowed, including parameters of the return
isomorphism, provided the interpretations and the isomorphism exist in
$\mathcal{N}^{(J)}$; this covers a parameter condition that $T$ proves to be
inhabited and under which $T$ proves all the required facts.
\end{theorem}

\begin{proof}
Apply a hypothetical bi-interpretation in $\mathcal{N}^{(J)}$, which yields a
classical model $\mathcal{M}$ of $\iPRA$. Let $\vec c$ collect the parameters
of the first interpretation, of the return isomorphism, and representatives in
$\mathcal{N}^{(J)}$ of the parameters of the second interpretation.
\Cref{obs:lem:involution} gives a non-identity involution $\sigma$ of
$(\mathcal{N}^{(J)},\vec c)$. Its image in $\operatorname{Aut}(\mathcal{M})$
has order dividing two and preserves the definable order, so it is the identity
by \cref{obs:lem:rigid}. By \cref{obs:lem:aut} the homomorphism is injective,
so $\sigma$ is the identity on every sort, which contradicts its choice on the
numerical sort. Faithfulness of the action on the other sorts is not needed.
\end{proof}

\begin{corollary}
\label{obs:cor:symmetry-cases}
None of the following theories is bi-interpretable with $\iPRA$ in the sense of
\cref{obs:def:biinterp}:
\begin{enumerate}[label=(\roman*)]
\item $T_{\mathrm{alg}}$, also with every set of universal Horn sentences true
in $\mathbb{N}$ adjoined;
\item for every finitary expansion $\mathcal{N}$ of $(\mathbb{N},0,\Suc)$, the
theory of all universal Horn sentences true in $\mathcal{N}$ with induction for
atomic formulas or for finite conjunctions of atomic formulas;
\item $T_{\mathrm{pp}}(\mathcal{N})$ and $T_{\mathcal{C}}(\mathcal{N})$, and
every subtheory of them, recursively axiomatized or not;
\item the many-sorted theories $T_{\mathcal{C}}(\mathcal{A}')$ for a
many-sorted $\mathcal{A}'$ with numerical sort $(\mathbb{N},0,\Suc)$, with
induction on the numerical sort for conjunctive matrices whose parameters and
witnesses may range over the other sorts.
\end{enumerate}
\end{corollary}

\begin{proof}
Each theory is true in $\mathcal{E}$ by \cref{obs:cor:power-axioms}. In the
many-sorted case the family varied in the induction argument belongs to the
numerical sort and the parameters of the other sorts stay fixed.
\end{proof}

\begin{remark}[Function objects]
The many-sorted case applies when a sort consists of numerical functions and
the signature has a total application operation. Extensionality,
$\bigl(\forall n\,\operatorname{app}(f,n)=\operatorname{app}(g,n)\bigr)\to f=g$,
has a conjunctive antecedent and is preserved. Adding function objects,
application, extensionality and equational combinator laws does not by itself
remove the symmetry.
\end{remark}

\begin{corollary}[The case-free progressive calculus]
\label{obs:cor:casefree}
Every rule of the case-free progressive calculus is sound in $T_{\mathrm{alg}}$
under the readings of \cref{obs:def:progressive}; \rulename{Iter} is induction
for a literal equation. Hence $\mathcal{E}$ is a model of its ordinary
first-order realization, and that realization is not bi-interpretable with
$\iPRA$ in the sense of \cref{obs:def:biinterp}.
\end{corollary}

\begin{remark}[Scope]
\label{obs:rem:symmetry-scope}
\Cref{obs:thm:symmetry} is about ordinary model-theoretic bi-interpretations,
or about a native proposal whose claimed interpretations are proved to induce
such interpretations under the displayed readings. Soundness of a native
calculus in an ordinary theory does not show that the calculus derives the
logical rules of that theory, and a proof translation that changes the meaning
of the connectives need not give an ordinary interpretation; the theorem does
not address such translations. It does not exclude equiconsistency, one-way
interpretations, disjunctive axioms such as coverage or a Boolean range
assertion, non-regular rules, or expansions that are not interpreted
coordinatewise. A parameter condition that excludes $\mathcal{E}$ and is not
proved inhabited there is outside the claim. Changing a term-defined operation
to a positive existential graph does not escape the theorem, since existential
witnesses and conjunctive induction of arbitrary alternation are allowed.
\end{remark}

\begin{remark}
\pkg{research} proves \cref{obs:thm:symmetry} for $T_{\mathrm{alg}}$, universal
Horn sentences and conjunctions of atoms. \pkg{compact\_kernel} proves it for
$T_{\mathrm{pp}}$, for $T_{\mathcal{C}}$ and for many sorts, with the order
$x\le y\iffdef\exists z\,(x+z=y)$ in place of the one in \cref{obs:lem:rigid}.
\pkg{checker\_and\_barriers} restates the argument for
\cref{obs:cor:casefree}.
\end{remark}

\section{Majorants: almost unary term languages}
\label{obs:sec:majorants}

This section draws on \pkg{counter\_congruence}, \pkg{research},
\pkg{local\_scopes}, \pkg{adaptive\_seed}, \pkg{uniform\_predecessor} and
\pkg{validated\_leaves}. All
statements concern the standard interpretation of fixed finite terms. The
bounding functions belong to the metatheory; they are not asserted to be terms
of the reduced languages.

\subsection{Zero-seeded recursion is monotone}

\begin{proposition}[\pkg{counter\_congruence}]
\label{obs:prop:monotone}
Every function denoted by a term of the zero-seeded language of
\cref{obs:def:languages}(a) is nondecreasing in each coordinate. Hence no term
denotes the zero test $z(0)=1$, $z(\Suc n)=0$, and the terms do not denote all
primitive recursive functions.
\end{proposition}

\begin{proof}
Induction on terms. Variables, $0$ and $\Suc$ are monotone. For
$\mathsf{L}(n;u,v.h)$ the update $h(u,v,\vec p)$ is coordinatewise
nondecreasing by the induction hypothesis. Fix $\vec p$ and write $r_0=0$,
$r_{k+1}=h(k,r_k,\vec p)$. Then $r_0\le r_1$, and $r_k\le r_{k+1}$ gives
$r_{k+1}=h(k,r_k,\vec p)\le h(k+1,r_{k+1},\vec p)=r_{k+2}$. So the value is
nondecreasing in the counter. Induction on $k$ gives
$r_k(\vec p)\le r_k(\vec p\,')$ for $\vec p\le\vec p\,'$ coordinatewise.
Substitution of the counter term and of parameter terms preserves monotonicity
by the induction hypothesis. The zero test is not monotone at $0<1$.
\end{proof}

The statement is formalized as \Lean{ZTm.mono} and \Lean{ZTm.no\_zero\_test}
in the module \Lean{ZeroSeedObstruction}; \Lean{Tm.val\_zf} records that
$\sel(x,1,0)$ of the iterator language is the zero test. The proposition
excludes the direct term compiler only. The tied-seed language of
\cref{obs:def:languages}(b) escapes it, since $Z$ of \cref{obs:lem:Bterms} is
the zero test.

\subsection{Bounding coordinates}

\begin{definition}
\label{obs:def:bc}
A function $f\colon\mathbb{N}^r\to\mathbb{N}$ has the \emph{bounding
coordinate} $i$ if there is a nondecreasing unary primitive recursive function
$F$ with $f(x_1,\dots,x_r)\le F(x_i)$ for all arguments. A function is
\emph{almost unary} if it is globally bounded or has a bounding coordinate.
\end{definition}

The name is the one used in clone theory for functions bounded by a unary
function of one coordinate, whose closure under composition is standard
(Pinsker, \emph{Clones containing all almost unary functions}, Algebra
Universalis 51 (2004)). The primitive recursive majorants are the additional
content here.

\begin{lemma}[Envelopes and composition]
\label{obs:lem:envelope}
\begin{enumerate}[label=(\roman*)]
\item For a primitive recursive map $g\colon\mathbb{N}^r\to\mathbb{N}^s$ the
function $\widehat g(z)=\max\{g(\vec z)_k: z_1,\dots,z_r\le z,\ k\le s\}$ is a
nondecreasing unary primitive recursive function that bounds every output
coordinate of $g$ whenever all inputs are at most $z$. In particular a unary
bound $F$ may be replaced by its nondecreasing envelope
$\widehat F(z)=\max_{m\le z}F(m)$.
\item Constants, projections, unary primitive recursive functions and globally
bounded functions are almost unary.
\item If $f$ and $g_1,\dots,g_r$ are almost unary, so is
$f(g_1(\vec x),\dots,g_r(\vec x))$.
\end{enumerate}
\end{lemma}

\begin{proof}
(i) and (ii) are immediate. For (iii), if $f$ is globally bounded so is the
composite. If $f\le F(x_j)$, only $g_j$ matters: a constant bound $c$ of $g_j$
gives the constant bound $F(c)$, and $g_j\le N(x_i)$ gives $F(N(x_i))$, by
monotonicity of $F$.
\end{proof}

\subsection{Adaptive query programs with several registers}
\label{obs:sec:programs}

This subsection is from \pkg{validated\_leaves}. It gives one barrier for
constructors whose seed and update both query one supplied body adaptively, with
several registers and with query locations that depend on the target and the
elapsed counter; \cref{obs:thm:single-body,obs:thm:circuit} below are special
cases. The restriction that makes the barrier work is on the data interface:
numerical outputs are functions of oracle answers only.

\begin{definition}[\pkg{validated\_leaves}]
\label{obs:def:programs}
Fix integers $q,d\ge 1$ and $c\ge 0$. A \emph{query program} with $c$ control
inputs, $q$ query coordinates and output dimension $d$ is defined together with
the number $j$ of answers available at a node: a program with $j$ answers is
\begin{enumerate}[label=(\alph*)]
\item a \emph{leaf} $F$, where $F\colon\mathbb{N}^j\to\mathbb{N}^d$ is a fixed
primitive recursive map; or
\item a \emph{query node} $(Q,P')$, where $Q\colon\mathbb{N}^{c+j}\to\mathbb{N}^q$
is a fixed primitive recursive map and $P'$ is a program with $j+1$ answers; or
\item a \emph{test node} $(T,P',P'')$, where $T\subseteq\mathbb{N}^{c+j}$ is a
fixed primitive recursive predicate and $P',P''$ are programs with $j$ answers.
\end{enumerate}
A query program is a program with $0$ answers; it is a finite tree. For an oracle
$O\colon\mathbb{N}^q\to\mathbb{N}$ and controls $\vec c\in\mathbb{N}^c$, the
value $P^{O}(\vec c)\in\mathbb{N}^d$ is computed from the root with the empty
answer list: a query node $(Q,P')$ reached with answers $a_1,\dots,a_j$ sets
$a_{j+1}=O\bigl(Q(\vec c,a_1,\dots,a_j)\bigr)$ and continues with $P'$; a test
node continues with $P'$ if $T(\vec c,a_1,\dots,a_j)$ holds and with $P''$
otherwise; a leaf returns $F(a_1,\dots,a_j)$.

An \emph{adaptive-program scheme} is given by $q,d\ge 1$, a seed program $P_0$
with one control input, an update program $P_1$ with $d+2$ control inputs, both
with $q$ query coordinates and output dimension $d$, and a fixed primitive
recursive $G\colon\mathbb{N}^d\to\mathbb{N}$. Its supplied subdescriptions are a
counter term $t(\vec p)$ and one body $h(\vec u,\vec p)$ with $q$ local
variables $\vec u$ and passive variables $\vec p$. With $n=t(\vec p)$ and
$h_{\vec p}(\vec a)=h(\vec a,\vec p)$, the constructor denotes $G(\vec r_n)$,
where
\begin{equation}
\label{obs:eq:program-recursor}
\vec r_0=P_0^{h_{\vec p}}(n),\qquad
\vec r_{m+1}=P_1^{h_{\vec p}}(n,m,\vec r_m).
\end{equation}
The \emph{adaptive-program term language} has variables, any primitive
operations that are almost unary, composition, and the constructors of all
adaptive-program schemes, arbitrarily nested; the body and the counter term of a
constructor are terms of the same language.
\end{definition}
% ed.: the package describes a program as queries Q_j(c,a_1..a_{j-1}) followed by a finite decision tree over output leaves, and its proof works path by path; the tree above, in which tests and queries interleave, makes that format explicit. A test that chooses the next query is also expressible by definition by cases inside Q, so nothing is gained or lost.

The programs reach $\vec p$ only through the oracle $h_{\vec p}$. The controls
$n$, $m$ and the old registers $\vec r$ may determine query locations and tests,
but no leaf reads them: a leaf is a function of the answers alone, so a test on
$m$ may choose between two fixed output maps but cannot copy $m$ into a
register. The number of queries on every path is bounded by the depth of the
tree, independently of the controls and of the answers; the only unbounded
iteration is the displayed recursion on $m$. A program with no test whose query
maps ignore the earlier answers is \emph{non-adaptive}.

\begin{lemma}[Bounds for a query program]
\label{obs:lem:programs}
Let $P$ be a query program with $c$ control inputs and $q$ query coordinates.
\begin{enumerate}[label=(\roman*)]
\item There is a nondecreasing unary primitive recursive $B_P$ such that every
output coordinate of $P^{O}(\vec c)$ is at most $B_P(M)$, for all $\vec c$,
whenever $O(\vec u)\le M$ for all $\vec u$.
\item Let $H$ be a nondecreasing unary primitive recursive function and
$\ell\le q$. There is a nondecreasing unary primitive recursive $K_{P,H,\ell}$
such that every output coordinate of $P^{O}(\vec c)$ is at most
$K_{P,H,\ell}(s)$ whenever all coordinates of $\vec c$ are at most $s$ and
$O(\vec u)\le H(u_\ell)$ for all $\vec u$.
\end{enumerate}
\end{lemma}

\begin{proof}
Write $\widehat g$ for the envelope of \cref{obs:lem:envelope}(i); for a map
without inputs it is the largest output coordinate.

(i) Every answer on the path taken is at most $M$. A leaf $F$ reached with $j$
answers returns coordinates at most $\widehat F(M)$. Put
$B_P(M)=\max_F\widehat F(M)$ over the finitely many leaves. Tests select the
leaf and query locations select the answers, but neither enters the bound; a
query location may be arbitrarily large while its answer is at most $M$.

(ii) Let $b$ be the largest number of query nodes on a path of $P$, and let
$L(s)=\max_Q\widehat Q(s)$ over the finitely many query nodes. Put
\[
A_0(s)=s,\qquad A_{j+1}(s)=\max\bigl(A_j(s),\,H(L(A_j(s)))\bigr)\quad(j<b).
\]
Each $A_j$ is nondecreasing and primitive recursive, and $A_j\le A_{j+1}$. By
induction along the path taken, at a node reached with $j$ answers all controls
and all answers are at most $A_j(s)$. For $j=0$ this is the hypothesis on
$\vec c$; a test changes nothing; at a query node $(Q,P')$ every coordinate of
$Q(\vec c,a_1,\dots,a_j)$ is at most $\widehat Q(A_j(s))\le L(A_j(s))$, so the
new answer is at most $H(L(A_j(s)))\le A_{j+1}(s)$, since $H$ is
nondecreasing. A leaf $F$ reached with $j\le b$ answers returns coordinates at
most $\widehat F(A_j(s))\le\widehat F(A_b(s))$. Put
$K_{P,H,\ell}(s)=\max_F\widehat F(A_b(s))$ over the leaves.
\end{proof}

\begin{theorem}[Adaptive iteration with several registers]
\label{obs:thm:adaptive}
Every term of the adaptive-program term language denotes an almost unary
function.
\end{theorem}

\begin{proof}
Induction on descriptions; the local and passive variables of a body are
independent inputs for the induction hypothesis. Variables are projections, the
primitive cases hold by assumption and composition is
\cref{obs:lem:envelope}(iii). Consider a constructor with programs $P_0,P_1$ and
output map $G$, and let $\widehat G$ be the envelope of $G$.

\emph{$h$ is bounded by a constant $M$.} By \cref{obs:lem:programs}(i) every
coordinate of $\vec r_0$ is at most $B_{P_0}(M)$ and every coordinate of
$\vec r_{m+1}$ is at most $B_{P_1}(M)$, so the result is at most the constant
$\widehat G\bigl(\max(B_{P_0}(M),B_{P_1}(M))\bigr)$.

\emph{$h(\vec u,\vec p)\le H(p_i)$ for a passive coordinate.} For fixed
$\vec p$ the oracle $h_{\vec p}$ is bounded by $M=H(p_i)$, whatever its local
arguments. The same argument bounds the result by
$\widehat G\bigl(\max(B_{P_0}(H(p_i)),B_{P_1}(H(p_i)))\bigr)$, a nondecreasing
primitive recursive function of $p_i$ alone; the counter is not used.

\emph{$h(\vec u,\vec p)\le H(u_\ell)$ for a local coordinate.} Let
$K_0=K_{P_0,H,\ell}$ and $K_1=K_{P_1,H,\ell}$ be the bounds of
\cref{obs:lem:programs}(ii), and define by primitive recursion
\[
W(0,n)=K_0(n),\qquad
W(m+1,n)=\max\bigl(W(m,n),\,K_1(n+m+W(m,n))\bigr).
\]
Induction on $m$ shows that every coordinate of $\vec r_m$ is at most $W(m,n)$:
the seed has the control $n$; at the step the number $n+m+W(m,n)$ bounds the
target counter, the elapsed counter and every old register, which are the
controls of $P_1$. Hence the result is at most $\widehat G(W(n,n))$, and
$K(n)=\max_{z\le n}\widehat G(W(z,z))$ is a nondecreasing unary primitive
recursive bound, independent of every passive parameter.

In the last case apply the induction hypothesis to the counter term: a constant
bound $c$ of $t$ gives the bound $K(c)$, and $t(\vec p)\le N(p_i)$ gives
$K(N(p_i))$. The cases exhaust the induction hypotheses for $h$.
\end{proof}

If the primitive operations are primitive recursive, so is every term: a fixed
query program is a finite composition with definition by cases, primitive
recursive relative to its body, and the outer loop is primitive recursion on
$m$.

\begin{remark}[Effectiveness]
Given bounding-coordinate or global-bound witnesses for the primitive
operations of a description, the proof constructs a coordinate and a majorant
effectively. A bound of an arbitrary globally bounded primitive recursive
function cannot be computed from its bare description: for a machine $e$, the
function $b_e(x,y)$ equal to the first halting time of $e$ if it is at most
$\max(x,y)$ and $0$ otherwise is globally bounded, and a uniform extractor of
bounds would decide halting. Unary primitives have computable envelopes and
Boolean predicates have the bound $1$, so these need no witness. The existence
statement of the theorem needs no witnesses.
\end{remark}

\subsection{Single-body constructors}

\begin{definition}[\pkg{uniform\_predecessor}]
\label{obs:def:single-body}
A \emph{single-body constructor scheme} is given by integers $a,b,d,q\ge 1$ and
fixed primitive recursive maps
\[
A_i\colon\mathbb{N}\to\mathbb{N}^q\ (i\le a),\quad
U_j\colon\mathbb{N}^{d+2}\to\mathbb{N}^q\ (j\le b),\quad
F_0\colon\mathbb{N}^a\to\mathbb{N}^d,\quad
F_1\colon\mathbb{N}^b\to\mathbb{N}^d,\quad
G\colon\mathbb{N}^d\to\mathbb{N}.
\]
Its supplied subdescriptions are a counter term $t$ and one body
$h(\vec u,\vec p)$ with $q$ local variables $\vec u$ and passive variables
$\vec p$. With $n=t(\vec p)$, the constructor denotes $G(\vec r_n)$, where
\begin{align}
\vec r_0&=F_0\bigl(h(A_1(n),\vec p),\dots,h(A_a(n),\vec p)\bigr),\label{obs:eq:seed}\\
\vec r_{m+1}&=F_1\bigl(h(U_1(n,m,\vec r_m),\vec p),\dots,h(U_b(n,m,\vec r_m),\vec p)\bigr).\label{obs:eq:step}
\end{align}
The maps are fixed before the body is supplied. No passive parameter and no
retained register enters $F_0$, $F_1$ or $G$ directly: passive data and
retained state reach the seed and update combiners only through evaluations of
the one body $h$. The target counter $n$ and the elapsed counter $m$ may enter
the probe locations. The \emph{single-body term language} has variables, $0$,
any primitive operations that are almost unary, composition, and the
constructors of all schemes, arbitrarily nested.
\end{definition}

The tied-seed recursor is the scheme $d=1$, $q=2$, $a=b=1$, $A_1(n)=(0,0)$,
$U_1(n,m,r)=(m,r)$, $F_0(z)=z+1$, $F_1(z)=z$, $G(r)=r$.

\begin{theorem}[Single-body information barrier]
\label{obs:thm:single-body}
Every term of the single-body term language denotes an almost unary function.
\end{theorem}

\begin{proof}
A single-body constructor is an adaptive-program constructor with the same
counter term, body, $d$, $q$ and $G$. The seed program has no test, queries
$h_{\vec p}$ at $A_1(n),\dots,A_a(n)$ and has the leaf $F_0$; the update program
has no test, queries $h_{\vec p}$ at $U_1(n,m,\vec r),\dots,U_b(n,m,\vec r)$
and has the leaf $F_1$. Both are non-adaptive, and
\cref{obs:eq:seed,obs:eq:step} are \cref{obs:eq:program-recursor} for them. By
induction on descriptions every term of the single-body term language is
therefore a term of the adaptive-program term language with the same
denotation, and \cref{obs:thm:adaptive} applies.
\end{proof}

\subsection{Adaptive query circuits}

\begin{definition}[\pkg{adaptive\_seed}]
\label{obs:def:circuit}
A \emph{query circuit} $P$ is a finite acyclic circuit with no input nodes. Its
nodes are numerals, applications of fixed primitive recursive functions to
earlier nodes, and queries $\mathit{Q}(a,b)$ whose two arguments are earlier
nodes; one node is the output. For an oracle
$\mathit{Q}\colon\mathbb{N}^2\to\mathbb{N}$ the output is $P^{\mathit{Q}}$.
Later queries may depend on earlier answers. For an update term
$h(u,v,\vec p)$ put $h_{\vec p}(a,b)=h(a,b,\vec p)$; the \emph{tied-seed
recursor with circuit $P$} is
\[
\mathsf{R}_P(0;u,v.h)=P^{h_{\vec p}},\qquad
\mathsf{R}_P(n+1;u,v.h)=h(n,\mathsf{R}_P(n;u,v.h),\vec p).
\]
The circuit reaches $\vec p$ only through queries to $h_{\vec p}$; it receives
no separate initializer, no counter and no passive parameter. The
\emph{circuit term language} has constants, projections, any primitive
operations that are almost unary, substitution, and the constructors
$\mathsf{R}_P$ for all circuits $P$, arbitrarily nested. The combining
operations inside a circuit are arbitrary primitive recursive functions.
\end{definition}

The tied-seed recursor $\mathsf{B}$ is the circuit with one query
$\mathit{Q}(0,0)$ followed by successor. A circuit may, for example, compute
$q=h(0,0,\vec p)$ and return $q+h(q,q+1,\vec p)^2$.

\begin{lemma}
\label{obs:lem:circuit}
Every query circuit $P$ is a query program without control inputs, with $q=2$
and $d=1$, in the sense of \cref{obs:def:programs}: there are fixed primitive
recursive maps $V_1,\dots,V_k$ and $E$ such that $P^{\mathit{Q}}$ is computed by
the queries $a_j=\mathit{Q}(V_j(a_1,\dots,a_{j-1}))$ for $j\le k$ followed by
the leaf $E(a_1,\dots,a_k)$, for every oracle $\mathit{Q}$. Consequently
\begin{enumerate}[label=(\roman*)]
\item there is a unary primitive recursive $C_P$ such that
$P^{\mathit{Q}}\le C_P(M)$ whenever $\mathit{Q}(a,b)\le M$ for all $a,b$;
\item if $\mathit{Q}(a,b)\le F(a)$ for all $a,b$, with $F$ a fixed unary
primitive recursive function, then $P^{\mathit{Q}}$ is bounded by a constant
that depends on $P$ and $F$ only. The same holds with $F(b)$.
\end{enumerate}
\end{lemma}

\begin{proof}
List the query nodes in a topological order of the circuit. By induction along
that order, the value of every node is a fixed primitive recursive function of
the answers of the query nodes before it: a numeral is a constant function, an
operation node composes a fixed primitive recursive function with the functions
of its inputs, and a query node is its own answer. Applied to the two arguments
of the $j$th query node this gives $V_j$, and applied to the output node it gives
$E$. Then (i) is \cref{obs:lem:programs}(i), and (ii) is
\cref{obs:lem:programs}(ii) with $\ell=1$ or $\ell=2$, $H=\widehat F$ and
$s=0$, since there are no controls.
\end{proof}

\begin{theorem}[Tied-seed circuit obstruction]
\label{obs:thm:circuit}
Every term of the circuit term language denotes an almost unary function.
\end{theorem}

\begin{proof}
$\mathsf{R}_P(n;u,v.h)$ is the adaptive-program constructor with the counter
term $n$, the body $h$, $q=2$, $d=1$, $G(r)=r$, the seed program given for $P$
by \cref{obs:lem:circuit} (its control $n$ unused), and the update program with
controls $(n,m,r)$, the single query location $(m,r)$ and the leaf $a_1$.
Indeed \cref{obs:eq:program-recursor} then reads $r_0=P^{h_{\vec p}}$ and
$r_{m+1}=h(m,r_m,\vec p)$, which is the definition of $\mathsf{R}_P$ at $0$ and
$m+1$. Constants, projections and substitution are variables, primitive
operations and composition. So every term of the circuit term language is a
term of the adaptive-program term language with the same denotation, and
\cref{obs:thm:adaptive} applies.
\end{proof}

\begin{remark}
The direct proof of \pkg{adaptive\_seed} gives sharper bounds in the two local
cases. If $h(u,v,\vec p)\le F(u)$ or $h(u,v,\vec p)\le F(v)$, then $r_0$ is at
most a constant $b_0$ by \cref{obs:lem:circuit}(ii). In the first case
$r_{j+1}\le F(j)$, so $\max(b_0,\widehat F(n))$ bounds $r_n$. In the second case
$K(z)=\max(z,\widehat F(z))$ is nondecreasing with $K\ge\operatorname{id}$, and
induction gives $r_n\le K^{n}(b_0)$. If $h\le c$ or $h\le F(p_i)$, writing $M$
for the bound, $\max(C_P(M),M)$ bounds the result.
\end{remark}

\begin{corollary}[The tied-seed language]
\label{obs:cor:tied-seed}
Every term of the tied-seed language of \cref{obs:def:languages}(b), and of its
unary-update variant, denotes an almost unary function. This remains true after
adjoining all unary primitive recursive functions and any operations that are
globally bounded or have a bounding coordinate, for example a primitive
characteristic of equality with values in $\{0,1\}$, predecessor, zero tests
and minimum.
\end{corollary}

\begin{proof}
The language is a sublanguage of both the single-body and the circuit term
language. Directly: if $h\le c$ the values are at most $c+1$; if
$h\le F(p_i)$ they are at most $1+F(p_i)$; if $h\le F(u)$ the value at the
counter $m$ is at most $1+F(m)$; if $h\le F(v)$ it is at most $\kappa(m)$ with
$\kappa(0)=1+F(0)$ and $\kappa(m+1)=\max(\kappa(m),F(\kappa(m)))$.
\end{proof}

\begin{corollary}[What almost unary terms cannot denote]
\label{obs:cor:no-addition}
No term of the adaptive-program term language, and hence none of the
single-body or of the circuit term language, denotes addition, maximum, the
selector with $\sel(0,a,b)=a$ and $\sel(\Suc n,a,b)=b$, or an injection
$\mathbb{N}^2\to\mathbb{N}$. In particular the tied-seed language has no such
term and cannot compile the initialized iterator, since
$\It{b}{a}{v}{\Suc v}=a+b$.
\end{corollary}

\begin{proof}
By \cref{obs:thm:adaptive} it suffices that none of these functions is almost
unary. For addition or maximum, fix the supposed bounding coordinate and increase the
other argument; a global bound fails likewise. For the selector, a bound in the
control fails at the control $0$ as $a$ grows; a bound in $a$ fails at $a=0$
and control $1$ as $b$ grows; a bound in $b$ fails at $b=0$ and control $0$ as
$a$ grows. Additional dummy coordinates can be fixed. For an injection, fixing
the bounding coordinate at $0$ leaves infinitely many pairs whose values lie in
one finite interval.
\end{proof}

\begin{remark}[The boundary of the hypotheses]
\label{obs:rem:boundary}
The data interface is what the theorems restrict. (i) If the initializer
circuit receives one further passive input $y$, the update $h(u,v,x,y)=x$ and
the initializer $y+h(0,0,x,y)$ give $x+y$ at the counter $0$. (ii) A circuit
with two independent update oracles queries the constant updates $x$ and $y$
and adds the answers. (iii) If a passive initializer enters a register
directly, $r_0=a$, $r_{m+1}=\Suc r_m$ computes $a+n$. (iv) Direct register
retention escapes \cref{obs:thm:single-body,obs:thm:adaptive} even with a tied
seed:
$r_0=h(0,\vec p)$, $r_{m+1}=\Suc r_m$ with the constant body $h(u,a)=a$
computes $a+n$; here the old register reaches the update combiner without
passing through the body. (v) Finiteness of the circuit matters: for the oracle
that is $1$ exactly at $(u,v)=(x,y)$, every answer is at most $1$, and an
unbounded search over all pairs finds $(x,y)$ and returns $x+y$; such a search
is not a fixed finite circuit and is not total for all oracles. (vi) In
\cref{obs:thm:adaptive} a leaf may not read a control numerically. If the seed
and the update program both return $n+h(0,\vec p)$, with the body $h(u,a)=a$ and
the counter term $x$, then $\vec r_m=x+a$ at every stage and the output is
$x+a$; no register is retained, and finite tests on $n$ do not allow this use of
$n$. (vii) The query budget must be fixed in advance. If both programs make $n$
queries of the body $h(u,a)=a$ and return the sum of the answers, with the
counter term $x$, the output is $xa$: every answer is at most $a$, but $xa$ is
bounded neither through $a$ nor through $x$. Saying only that every execution
makes finitely many queries is not enough; (v) is the same failure for a budget
that depends on the answers. These examples mark the boundary. None of them is claimed to give a complete basis or a
smaller checker.
\end{remark}

\begin{remark}[Scope and sources]
The theorems are term-definability obstructions. They do not exclude
quantified graph definitions, a relational interpretation, a different
interpretation domain or additional axioms; an independent unbounded
initializer is one escape and is not asserted to be the only one. The
one-coordinate bound for $\mathsf{B}$ and the exclusion of addition and
selection are in \pkg{counter\_congruence}; \pkg{research} rederives them for
$\mathsf{F}$ and $\mathsf{B}$; \pkg{local\_scopes} adds the bounded-valued
primitives. \Cref{obs:thm:single-body} is from \pkg{uniform\_predecessor} and
\cref{obs:thm:circuit} from \pkg{adaptive\_seed}, each with a direct majorant
induction. The two classes are incomparable: the first has register vectors,
several probes in the update and counter-dependent probe locations, but a
non-adaptive seed; the second has an adaptive seed, but a scalar state and the
fixed update interface. \Cref{obs:thm:adaptive}, from \pkg{validated\_leaves},
is a common extension: later queries may depend on earlier answers in the seed
and in every update, with several registers and counter-dependent locations.
Its proof is the same majorant induction, and the article derives the two
earlier theorems from it. Globally bounded primitives such as an equality
characteristic are allowed in it, so it asserts none of the finite-profile
invariants of \cref{obs:sec:profiles}. \Cref{obs:cor:no-term-interp} applies to
its standard algebra as well.
\end{remark}

\section{Finite-dimensional term interpretations}
\label{obs:sec:quotient}

This section draws on \pkg{adaptive\_seed} and \pkg{local\_scopes}, which prove
the same theorem independently. It turns the bound of
\cref{obs:sec:majorants} into an obstruction to interpretations that use finite
tuples and a quotient equality, as long as the interpreted addition is given by
a tuple of terms. The theorem is about functions on $\mathbb{N}$; the domain
and the equivalence need not be definable.

\begin{theorem}[Finite quotient from coordinate bounds]
\label{obs:thm:quotient}
Let $d\ge 1$, $D\subseteq\mathbb{N}^d$, let $\sim$ be an equivalence relation
on $D$ and $e\in D$. Let $f\colon D\times D\to D$ satisfy
\[
f(e,z)\sim z\sim f(z,e)\qquad(z\in D).
\]
Suppose that for each output coordinate $i$ there are an input coordinate
$j_i\le d$ and a nondecreasing unary $g_i$ such that, uniformly on
$D\times D$, either $f(x,y)_i\le g_i(x_{j_i})$ or
$f(x,y)_i\le g_i(y_{j_i})$, or that $f(x,y)_i$ has a constant bound. Then
$D/{\sim}$ is finite. More precisely, with $L(z)=f(e,z)$ and $R(z)=f(z,e)$, a
composition of at most $\ell_d=\lceil\log_2(d+1)\rceil$ maps chosen from
$L,R$ sends every element of $D$ to a $\sim$-equivalent element of one finite
box $\prod_{i\le d}\{0,\dots,C_i\}$, and $D/{\sim}$ has at most
$\prod_i(C_i+1)$ elements.
\end{theorem}

\begin{proof}
A coordinate bounded through an $x$-input has the constant bound
$g_i(e_{j_i})$ under $L$ and the bound $g_i(z_{j_i})$ under $R$; for a
coordinate bounded through a $y$-input the roles are exchanged; a constantly
bounded coordinate is bounded under both. Regard $\{1,\dots,d\}$ as a set of
states that track dependencies of bounds. A letter sends a state $i$ to the
state $j_i$ if it leaves the dependence, and to an absorbing terminal state if
it gives a constant bound. Each state is terminated by at least one of the two
letters.

Choose the letters from the outside of the composition inward. Let $S$ be the
current set of live states, initially all $d$ states. The number of states of
$S$ terminated by $L$ plus the number terminated by $R$ is at least $|S|$, so
one letter terminates at least $\lceil|S|/2\rceil$ of them. Choose it and
replace $S$ by the set of images of the survivors; each survivor has one
image, so the new set has at most $\lfloor|S|/2\rfloor$ elements. After at most
$\ell_d$ choices $S$ is empty, since $2^{\ell_d}>d$.

Let the chosen letters be $V_1,\dots,V_k$ and $W=V_1\circ\dots\circ V_k$, so
that $V_k$ is applied first. By induction on the number of letters, each
coordinate of the partial composition $V_1\circ\dots\circ V_m$ is either
uniformly bounded or bounded by a composition of the functions $g_i$ along one
live dependency path, applied to a live coordinate of the argument; appending
the next letter on the inside replaces that coordinate by its bound, which is
legitimate because the $g_i$ are nondecreasing. When the path terminates the
bound is a constant. The empty final set gives constants $C_1,\dots,C_d$ with
$W(D)\subseteq\prod_i\{0,\dots,C_i\}$.

Both $L$ and $R$ map $D$ into $D$ and send every element to an equivalent one,
so $W(z)\sim z$ by transitivity. Every class meets the finite set
$D\cap\prod_i\{0,\dots,C_i\}$.
\end{proof}

Associativity, commutativity and compatibility of $f$ with $\sim$ are not
used. A unary bound may be replaced by its nondecreasing envelope, so
monotonicity is no extra assumption for term majorants. For $d=0$ the domain
has at most one element.

\begin{proposition}[The word length is sharp]
\label{obs:prop:sharp}
For every $d\ge 1$ there are $D=\mathbb{N}^d$, the universal equivalence, $e$
the zero tuple and an $f$ whose coordinates are projections or constants, such
that the shortest composition of $L,R$ with finite image has length exactly
$\lceil\log_2(d+1)\rceil$.
\end{proposition}

\begin{proof}
Let $k=\lfloor\log_2 d\rfloor$. If $k=0$ then $d=1$; take $f(x,y)=x$, so $R$ is
the identity and $L$ is constant, and the minimum is $1$. Let $k\ge 1$. Label
$2^k$ of the coordinates by words $w\in\{L,R\}^k$ and let $\sigma(w)$ be $w$
with its first letter removed and $R$ appended. Put
$f(x,y)_w=y_{\sigma(w)}$ if $w$ starts with $L$ and
$f(x,y)_w=x_{\sigma(w)}$ if $w$ starts with $R$; the remaining coordinates are
$0$. Under a letter, the state $w$ survives exactly when the letter is its
first letter, and then moves to $\sigma(w)$. For a word
$\alpha\in\{L,R\}^k$, read from the outside inward, the state $\alpha$
survives all $k$ letters, and the coordinate $\alpha$ of the composition is
the projection to the coordinate $R^k$ of the argument. So no word of length
$k$ has finite image, and neither has a shorter word, since extending a word
with finite image on the inside keeps the image finite. The word $L^{k+1}$ has
finite image: a state with an $R$ among its $k$ letters is terminated when that
letter is reached, and the state $L^k$ has become $R^k$ after $k$ letters and
is terminated by the last $L$. The minimum is $k+1=\lceil\log_2(d+1)\rceil$.
\end{proof}

\begin{corollary}[No finite-dimensional term interpretation of $\iPRA$]
\label{obs:cor:no-term-interp}
Let $T$ be a theory in a numerical language whose standard model on
$\mathbb{N}$ satisfies $T$ and in which every term denotes an almost unary
function; for example a theory true in the standard tied-seed algebra, in the
single-body algebra, in the circuit algebra or in the adaptive-program algebra
of \cref{obs:thm:adaptive}. There is no interpretation of
$\iPRA$ in $T$ such that the interpreted domain is a nonempty subset of a
finite power of the source domain, the interpreted equality is an arbitrary
equivalence relation on it, and the interpreted addition has one fixed tuple of
terms as realizers. Fixed numerical parameters are allowed, and successor may
be interpreted in any way consistent with the axioms. In particular there is no
bi-interpretation with these properties.
\end{corollary}

\begin{proof}
Apply the interpretation in the standard model. Let $e$ represent the
interpreted zero. The two zero laws of addition are theorems of $\iPRA$, so the
addition tuple satisfies the unit hypothesis of \cref{obs:thm:quotient}; each
of its coordinate terms is almost unary in the $2d$ input variables, and a
majorant in a fixed parameter is a constant. So the interpreted quotient is
finite. The interpreted successor is injective on the quotient and its image
omits the class of zero, which is impossible on a finite nonempty set.
\end{proof}

\begin{remark}[Limits]
\label{obs:rem:quotient-limits}
(i) Both unit laws are needed: on $\mathbb{N}$ with equality, $f(x,y)=x$
satisfies $f(z,e)=z$ and the coordinate bound, and the quotient is infinite.
(ii) The realizers must be terms. With decidable equality to zero the
quantifier-free graph
$(x=0\wedge z=y)\vee(x\ne 0\wedge z=x)$ is total and single-valued and defines
an associative operation on $\mathbb{N}$ with the two-sided unit $0$, which
selects the first nonzero argument; it is not almost unary. Likewise
$\max(x,y)$ has the unit $0$ and is defined by cases that choose different
controlling coordinates. Neither example is an interpretation of $\iPRA$; they
show that the relational route is not closed by \cref{obs:thm:quotient}. (iii)
The qualifier about the standard model is needed: a theory whose axioms exclude
the standard algebra is not covered. (iv) The logarithmic bound counts layers
of unit multiplication in a metatheoretic construction; it is not a bound on
term size, evaluation time or checker size.
\end{remark}

\begin{remark}
\pkg{adaptive\_seed} proves \cref{obs:thm:quotient} by the halving argument
given here and proves \cref{obs:prop:sharp}. \pkg{local\_scopes} proves the
theorem by counting: from each starting coordinate at most one binary word of
length $m$ keeps a dependence alive, so at most $d$ of the $2^m$ words are
admitted, and a word admitted from no coordinate exists when $2^m>d$.
\end{remark}

\section{Finite profiles and relational bi-interpretations}
\label{obs:sec:profiles}

This section draws on \pkg{local\_scopes}. The coordinate bound leaves two
questions open: whether the tied-seed language computes a bounded comparison of
two unbounded inputs, and whether an interpretation that defines its operations
by formulas escapes. A stronger invariant of tied-seed terms answers the first,
and, combined with witness extraction from proofs, answers the second for a
specified class of theories.

\subsection{Finite control}

\begin{definition}
\label{obs:def:fc}
A function $f\colon\mathbb{N}^r\to\mathbb{N}$, $r\ge 1$, has \emph{finite
control} if there are an active coordinate $i$, finite nonempty sets $A_j$ and
unary maps $\alpha_j\colon\mathbb{N}\to A_j$ for $j\ne i$, and unary functions
$F_a$ indexed by $a\in\prod_{j\ne i}A_j$, such that
\[
f(x_1,\dots,x_r)=F_{(\alpha_j(x_j))_{j\ne i}}(x_i).
\]
The value may depend without bound on one coordinate; every other coordinate
supplies only a member of a fixed finite set. The functions $F_a$ are the
\emph{profiles}. For primitive recursive $f$ all displayed maps are taken
primitive recursive, with finite encodings of the $A_j$; specified finite
codomains replace exact ranges, and no decision procedure for attained ranges
is assumed.
\end{definition}

A constant function has finite control with any active coordinate, and the
definition imposes no restriction at arity one.

\begin{lemma}
\label{obs:lem:fc-comp}
\begin{enumerate}[label=(\roman*)]
\item If $f$ has finite control and $\rho$ is unary with finite range, then
$\rho\circ f$ factors through finite-valued unary maps of all input
coordinates.
\item The functions with finite control contain the constants and projections
and are closed under composition.
\end{enumerate}
Both preserve primitive recursiveness of the displayed maps.
\end{lemma}

\begin{proof}
(i) Attach to the active coordinate the finite feature
$\beta_i(x)=(\rho(F_a(x)))_a$. Together with the $\alpha_j(x_j)$ it determines
$\rho(f(\vec x))$.

(ii) Consider $f(g_1(\vec x),\dots,g_r(\vec x))$ where $f$ has the active
input $i$ and $g_i$ has the active input $k$. Every other $g_j$ enters $f$
through the finite-valued $\alpha_j\circ g_j$, which by (i) is determined by
finite unary features of all coordinates of $\vec x$. Combine the features of
the coordinates other than $k$ with the finite controls of $g_i$. With these
combined controls fixed, the expression is one unary function of $x_k$; the
features that depend on $x_k$ are evaluated inside it. There are finitely many
combined controls.
\end{proof}

\begin{theorem}[Finite control of tied-seed terms]
\label{obs:thm:fc}
Every function of positive arity denoted by a term of the tied-seed language of
\cref{obs:def:languages}(b) has finite control, with primitive recursive
feature maps and profiles obtained by recursion on the term. A closed term is a
constant.
\end{theorem}

\begin{proof}
Variables and $0$ are covered by \cref{obs:lem:fc-comp}. It suffices to show
that $\mathsf{B}_h(n,\vec p)$ has finite control when the update
$h(u,v,\vec p)$ has; the actual counter term is then substituted by
\cref{obs:lem:fc-comp}(ii). Write $\gamma(\vec p)$ for a tuple of finite unary
features, one for each passive coordinate.

\emph{The active input of $h$ is the previous state:}
$h(u,v,\vec p)=F_{\alpha(u),\gamma(\vec p)}(v)$. For each of the finitely many
values $c$ of $\gamma$ define
$R_c(0)=1+F_{\alpha(0),c}(0)$ and $R_c(m+1)=F_{\alpha(m),c}(R_c(m))$. Then
$\mathsf{B}_h(n,\vec p)=R_{\gamma(\vec p)}(n)$, with active coordinate $n$.

\emph{The active input of $h$ is the old counter:}
$h(u,v,\vec p)=F_{\beta(v),\gamma(\vec p)}(u)$. Define
$R_c(0)=1+F_{\beta(0),c}(0)$ and $R_c(m+1)=F_{\beta(R_c(m)),c}(m)$. Again
$\mathsf{B}_h(n,\vec p)=R_{\gamma(\vec p)}(n)$.

\emph{The active input of $h$ is a passive parameter $p$.} Write $\vec q$ for
the other passive parameters and
$h(u,v,p,\vec q)=F_{\alpha(u),\beta(v),\gamma(\vec q)}(p)$. Form a finite
catalogue $(V_j)_{j\in I}$ of unary functions that contains every profile
$F_{a,b,c}$ and every initial profile
$V_{0,c}(p)=1+F_{\alpha(0),\beta(0),c}(p)$, with separate indices for separate
occurrences. Every state of the recursion is the value at $p$ of a catalogue
function. The only observation of a state used by the update is its
$\beta$-value; record all of them in the finite feature
\[
\kappa(p)=\bigl(\beta(V_j(p))\bigr)_{j\in I}.
\]
For each control value $c$ and each possible vector $k$, define a finite-state
process: its initial state is the index of $V_{0,c}$, and at the counter $m$
the state $j$ is followed by the index of $F_{\alpha(m),k_j,c}$. Let
$\theta_{c,k}(n)\in I$ be its state after $n$ transitions and pack the finitely
many processes into the finite-valued feature
$\eta(n)=(\theta_{c,k}(n))_{c,k}$. Induction on the counter gives
\[
\mathsf{B}_h(n,p,\vec q)=V_{\eta(n)_{\gamma(\vec q),\kappa(p)}}(p):
\]
the initial case selects $V_{0,c}$, and the transition uses exactly the
recorded observation $\beta(V_j(p))=\kappa(p)_j$. With
$H_{e,c}(p)=V_{e_{c,\kappa(p)}}(p)$ this reads
$\mathsf{B}_h(n,p,\vec q)=H_{\eta(n),\gamma(\vec q)}(p)$, with the active
coordinate $p$ and finite controls on $n$ and on each remaining parameter. The
processes are defined for all vectors $k$, including unrealized ones; only the
realized ones are used.
\end{proof}

The bound of \cref{obs:cor:tied-seed} follows from the theorem by taking the
maximum of the finitely many profiles and its running maximum. The new content
is that the other inputs supply only finitely many distinct profiles.

\begin{corollary}[No equality test, minimum or truncated subtraction]
\label{obs:cor:fc-excluded}
No term of the tied-seed language denotes the characteristic function of
equality or of inequality, $\min(a,b)$, or
$a\mathbin{\dot{-}}b=\max(a-b,0)$.
\end{corollary}

\begin{proof}
A binary function with finite control and active first coordinate has finitely
many distinct sections $a\mapsto f(a,b)$; with active second coordinate it has
finitely many sections $b\mapsto f(a,b)$. Each displayed function has
infinitely many distinct sections in both directions: for equality the section
at $b$ is the characteristic function of $\{b\}$; the functions
$x\mapsto\min(x,b)$ differ at an $x$ above both parameters; for truncated
subtraction, distinct second arguments are separated at a large first argument,
and distinct first arguments at the second argument $0$.
\end{proof}

\begin{remark}
The theorem does not decide whether every unary primitive recursive function is
denoted by a tied-seed term, since finite control is no restriction at arity
one. It does not concern formula-defined relations: the source has logical
equality $a=b$, which is not a term that returns its characteristic value.
Bounded-valued primitives preserve the coordinate bound
(\cref{obs:cor:tied-seed}) and need not preserve finite control; an adjoined
equality characteristic violates it.
\end{remark}

\subsection{An existential-free presentation}

Let $T_{\mathsf{B}}$ be intuitionistic first-order logic with equality in the
staged signature of the tied-seed language, with the computation equations
\eqref{obs:eq:B0}, \eqref{obs:eq:BS}, coverage
$\forall x\,(x=0\vee\exists y\,x=\Suc y)$, separation $0=\Suc 0\to\bot$, and
the universal closures of
\[
\bigl(E(0,\vec p)\wedge\forall n\,(E(n,\vec p)\to E(\Suc n,\vec p))\bigr)\to\forall n\,E(n,\vec p)
\]
for literal equations $E$. By \cref{obs:lem:Bterms} the term $P$ satisfies
$P(0)=0$ and $P(\Suc y)=y$ by computation, so coverage is equivalent, without
induction, to the existential-free sentence
$\forall x\,(x=0\vee x=\Suc(Px))$. Equality is treated as a binary predicate
with universal axioms of reflexivity, symmetry, transitivity and congruence for
every function symbol; transport through formulas is derived. Thus
$T_{\mathsf{B}}$ has an existential-free axiomatization in its own language.
An infinite axiom scheme causes no difficulty, since a proof uses finitely many
axioms.

\begin{lemma}[Finite witness extraction]
\label{obs:lem:witness}
Let the language contain the closed term $0$ and let $\Gamma$ be a finite set
of existential-free formulas. If intuitionistic first-order logic proves
$\Gamma\proves\exists z_1\cdots\exists z_d\,R(\vec x,\vec z)$, then there are
finitely many tuples of terms $\vec t_1(\vec x),\dots,\vec t_M(\vec x)$ with
\[
\Gamma\proves\bigvee_{\ell\le M}R(\vec x,\vec t_\ell(\vec x)).
\]
The formula $R$ is arbitrary; it may contain existential quantifiers. No
uniqueness of witnesses is assumed.
\end{lemma}

\begin{proof}
Take a cut-free intuitionistic sequent proof; cut elimination for
intuitionistic first-order logic is standard (for instance Avigad,
\emph{Algebraic proofs of cut elimination}, J.~Log.\ Algebr.\ Program.\ 49
(2001), Theorem 7.6). Equality is handled by the existential-free axioms, so no
equality extension of the theorem is needed.

Starting at the end sequent, trace the branches whose succedent carries the
displayed leading existential quantifiers, and keep the prefix of witness terms
already selected. At a left implication rule keep the proof of the first
premise unchanged and follow the second premise; at a left disjunction rule
follow both premises; the other left and structural rules preserve the tracked
succedent. Every antecedent formula on a tracked branch is a subformula
instance of a formula of $\Gamma$ and is existential-free, so no left
existential rule occurs there. While a leading existential quantifier remains,
an initial sequent cannot close the branch, because the succedent is not an
antecedent formula; the only applicable right rule is existential
introduction. Record its witness term and continue. When all $d$ witnesses are
introduced, stop and keep the remaining proof of the matrix instance
unchanged; the matrix is not analyzed. If a tracked branch closes by falsity
elimination, fill the unselected witnesses with $0$; the matrix instance
follows from the same falsity. There are finitely many branches, so finitely
many complete tuples $\vec t_\ell$ are recorded.

Replace each stopped branch by its kept proof of $R(\vec x,\vec t_\ell)$
followed by introductions into the common disjunction, and reassemble the
tracked left and structural inferences. An existential introduction that
supplied a recorded component is omitted, since the disjunction is already the
succedent of its transformed premise. Left implication inferences are
reconstructed from the unchanged first premise and the transformed second
premise.

No numerical parameter is bound on a tracked branch: left existential rules are
absent, and right universal introduction cannot have the tracked succedent.
Extraneous free variables of the proof are replaced by $0$ after renaming
eigenvariables, so the witness terms contain only $\vec x$ and the retained
parameters.
\end{proof}

The lemma does not give a single uniform witness term; a finite disjunction is
what the next theorem needs.

\subsection{No compression of scalars}

\begin{theorem}[No provable compression from $k$ scalars to $d$ coordinates]
\label{obs:thm:no-compression}
Let $T$ be an intuitionistic theory in the tied-seed language with an
existential-free axiomatization and with the standard tied-seed algebra as a
model. Let $k>d\ge 1$ and let $R(\vec x,\vec z)$ be any formula with $k$
variables $\vec x$ and $d$ variables $\vec z$. If
$T\proves\forall\vec x\,\exists\vec z\,R(\vec x,\vec z)$, then in the standard
model some output tuple $\vec z$ occurs with two different input tuples. In
particular $T$ does not also prove
\[
R(\vec x,\vec z)\wedge R(\vec x\,',\vec z)\to\bigwedge_{i\le k}x_i=x_i' .
\]
The relation may have many outputs for each input.
\end{theorem}

\begin{proof}
Instantiate at free $\vec x$ and collect the finitely many axioms used.
\Cref{obs:lem:witness} gives $d$-tuples of terms
$\vec t_1,\dots,\vec t_M$ such that for every input in the standard model some
$\vec t_\ell(\vec x)$ is an $R$-output. By \cref{obs:thm:fc} each coordinate
term $t_{\ell j}$ has finite control; let $q_{\ell j}$ be the number of its
profiles. On the grid $X_N=\{0,\dots,N-1\}^k$ it takes at most $q_{\ell j}N$
values, since its active argument takes at most $N$ values; constant terms are
included by adding an unused active argument. So $\vec t_\ell$ takes at most
$N^d\prod_j q_{\ell j}$ values on $X_N$, and the union of the images has at
most $CN^d$ elements with $C=\sum_\ell\prod_j q_{\ell j}$. If every output
belonged to at most one input, the $N^k$ inputs of $X_N$ would need $N^k$
distinct outputs in this union, which fails for $N>C$. Soundness in the
standard model turns a proof of the displayed implication into exactly that
condition.
\end{proof}

\begin{corollary}[Scalar pairing is not provably total]
\label{obs:cor:no-pairing}
In $T_{\mathsf{B}}$ no formula is provably a total coding of pairs of numbers
by single numbers with unique decoding. This includes a pairing relation
$\pi_1(c)=a\wedge\pi_2(c)=b$ with unary terms as projections, and arbitrary
formulas as projection relations, provided each code determines at most one
pair.
\end{corollary}

\begin{proof}
A provably total coding relation $R(a,b,c)$ with provably unique decoding is
the case $k=2$, $d=1$ of \cref{obs:thm:no-compression}.
\end{proof}

Even a proof that every unary primitive recursive function is a tied-seed term
would therefore not give pairing inside $T_{\mathsf{B}}$.

\subsection{No finite-dimensional bi-interpretation}

\begin{theorem}
\label{obs:thm:no-biinterp}
Let $T$ be an intuitionistic theory in the tied-seed language with an
existential-free axiomatization and with the standard tied-seed algebra as a
model; for example $T_{\mathsf{B}}$, or $T_{\mathsf{B}}$ with further
existential-free axioms true in the standard algebra. Then $T$ and $\iPRA$ are
not bi-interpretable in the sense of \cref{obs:def:biinterp}. The interpreted
domains, the operation graphs and the quotient equality are arbitrary
formulas. The statement is parameter-free, and it holds when fixed parameters
are retained uniformly in the interpretation proofs without assumptions that
are not existential-free.
\end{theorem}

\begin{proof}
Let $\mathcal{I}$ interpret $\iPRA$ in $T$ with dimension $d$, domain
$D\subseteq\mathbb{N}^d$ and interpreted equality $E$, and let $\mathcal{J}$
interpret $T$ in $\iPRA$ with dimension $e$. A zero-dimensional domain has
finitely many tuples and cannot take part in the source-side isomorphism with
the infinite source domain, so $d,e\ge 1$.

The composite represents source numbers by $e$-tuples of interpreted
arithmetic elements, that is, by $de$ source coordinates. The source-side
return isomorphism provides a formula $\eta(a,\vec u)$ of $T$, with
$\vec u\in D^e$ and the domain conditions included, such that
\[
T\proves\forall a\,\exists\vec u\,\eta(a,\vec u),
\qquad
T\proves\eta(a,\vec u)\wedge\eta(a',\vec u\,')\wedge E^{\mathcal{I}}_{\mathcal{J}}(\vec u,\vec u\,')\to a=a',
\]
where $E^{\mathcal{I}}_{\mathcal{J}}$ is the translation under $\mathcal{I}$
of the equality of $\mathcal{J}$. Only this part of the bi-interpretation is
used.

Fix $k\ge 1$. $\iPRA$ has a primitive recursive injective tupling of $ke$
numbers with primitive recursive projections, and proves its decoding
equations. Its translation under $\mathcal{I}$ is a provably total coding
relation from $(D/E)^{ke}$ to $D/E$ that is injective on input tuples. Define
$R_k(a_1,\dots,a_k,\vec c)$ in $T$: there are $\eta$-representatives
$\vec u_i\in D^e$ of all $a_i$, and the interpreted tupling codes the $ke$
arithmetic components of $(\vec u_1,\dots,\vec u_k)$ by $\vec c\in D$. From the
totality of $\eta$ and of the tupling, finite existential reasoning gives
\[
T\proves\forall a_1\cdots\forall a_k\,\exists c_1\cdots\exists c_d\,R_k(\vec a,\vec c).
\]
The inputs $a_i$ range over the whole source domain; no choice is used.

Suppose $R_k(\vec a,\vec c)$ and $R_k(\vec a\,',\vec c)$. The same tuple
$\vec c$ represents the same arithmetic element, by reflexivity of $E$ on $D$.
Injectivity of the interpreted tupling gives componentwise $E$-equality of the
arithmetic components of $\vec u_i$ and $\vec u_i'$, hence
$E^{\mathcal{I}}_{\mathcal{J}}(\vec u_i,\vec u_i')$ for each $i$, since
interpretations respect equality. The second property of $\eta$ gives
$a_i=a_i'$. So $T$ proves that the output fibers of $R_k$ are disjoint.
Surjectivity of $\eta$ and uniqueness of representatives are not needed.

With $k=d+1$ this contradicts \cref{obs:thm:no-compression}.
\end{proof}

\begin{corollary}[Provably inhabited parameter families]
\label{obs:cor:parameters}
The exclusion holds for a bi-interpretation with a finite tuple $\vec p$ of
$r$ source parameters subject to a condition $\Theta(\vec p)$, when
$T\proves\exists\vec p\,\Theta(\vec p)$ and $T$ proves the interpretation and
return-isomorphism facts uniformly under $\Theta(\vec p)$.
\end{corollary}

\begin{proof}
The construction with the parameters retained proves totality and disjoint
fibers of $R_{k,\vec p}$ under $\Theta(\vec p)$. Put
$\widehat R_k(\vec a,\vec p,\vec c)\defeq\Theta(\vec p)\wedge R_{k,\vec p}(\vec a,\vec c)$.
Then $T\proves\forall\vec a\,\exists\vec p\,\exists\vec c\,\widehat R_k$, and an
identical output $(\vec p,\vec c)$ uses identical parameters, so the fibers are
disjoint. There are $d+r$ output coordinates; take $k=d+r+1$. The formula
$\Theta$ is inside the output matrix, so no restriction on its quantifiers is
needed.
\end{proof}

\begin{remark}[Scope]
\label{obs:rem:biinterp-scope}
Positive existential domain formulas do not rescue a bi-interpretation: they
occur inside the matrix $R_k$, which witness extraction does not analyze, and
the return isomorphism removes every domain condition from the unrestricted
source inputs. The theorem does not exclude equiconsistency or one-way
interpretations on complicated domains; the return isomorphism is essential.
It does not cover a parameter condition that $T$ does not prove inhabited, a
theory with positive existence axioms that cannot be eliminated in the
tied-seed language, or every axiomatization in that language. It does not
extend automatically to additional bounded numerical primitives, which preserve
the coordinate bound and need not preserve finite control.
\end{remark}

\subsection{The selector boundary}

The obstruction has an exact positive counterpart: a bounded primitive does not
help (\cref{obs:cor:tied-seed}), and one unbounded selector restores primitive
recursion.

\begin{proposition}[\pkg{local\_scopes}]
\label{obs:prop:selector-repair}
Adjoin to the tied-seed language an operation $C$ with $C(0,a,b)=a$ and
$C(\Suc n,a,b)=b$, allowed inside update bodies. For terms $g(\vec p)$ and
$h(n,z,\vec p)$ and fresh $u,v$ put
\[
K(u,v,\vec p)\defeq C\bigl(u,\,g(\vec p),\,h(W(u),v,\vec p)\bigr),\qquad
J(n,\vec p)\defeq\mathsf{B}\bigl(\Suc n;u,v.K(u,v,\vec p)\bigr),
\]
with $W$ as in \cref{obs:lem:Bterms}. Equality reasoning and the computation
equations prove $J(0,\vec p)=g(\vec p)$ and
$J(\Suc n,\vec p)=h(n,J(n,\vec p),\vec p)$, with no induction, predecessor
search, equality characteristic, pairing or independent initializer.
\end{proposition}

\begin{proof}
Suppress $\vec p$. The first update discards the tied initial state:
$J(0)=\mathsf{B}(\Suc 0;u,v.K)=K(0,\mathsf{B}(0;u,v.K))=C(0,g,\dots)=g$. At a
successor,
$J(\Suc n)=\mathsf{B}(\Suc\Suc n;u,v.K)=K(\Suc n,\mathsf{B}(\Suc n;u,v.K))
=C(\Suc n,g,h(W(\Suc n),J(n)))=h(n,J(n))$. These hold for arbitrary parameters.
\end{proof}

Zero, successor and projections are terms and composition is substitution, so
every primitive recursive description compiles into the enlarged language, and
conversely each description of the enlarged language is primitive recursive.
This is numerical completeness only; the logical and interpretation claims for
such a basis need their own proofs and their own size accounting.

A formula definition of $C$ does not supply this compiler. With coverage and
$0\ne\Suc 0$, the term $Z$ decides $z=0$, and
$(z=0\wedge w=a)\vee(z\ne 0\wedge w=b)$ is provably total and single-valued, so
a symbol with this graph is a definitional extension. Two things do not follow.
The tied-seed signature has recursion symbols for its original update terms
only; a recursion whose update contains the new symbol is a new graph whose
totality and uniqueness need proof. And literal-equation induction for terms
with the new symbol translates into induction for formulas that are not literal
equations of the old language.

\section{Zero profiles of positive iteration}
\label{obs:sec:zeroprofile}

This section draws on \pkg{checker\_and\_barriers}. Positive iteration
$\mathsf{R}^{+}(n,a;v.h)=h^{n+1}(a)$ of \cref{obs:def:languages}(c) is not
almost unary: $\mathsf{R}^{+}(n,a;v.\Suc v)=n+a+1$ grows in both arguments. A
different invariant applies.

\begin{theorem}[Zero-profile theorem]
\label{obs:thm:zeroprofile}
For every term $t(x_1,\dots,x_r)$ built from variables, $0$, $\Suc$ and
$\mathsf{R}^{+}$, the zero set of $t$ in the standard interpretation is
$\varnothing$, $\mathbb{N}^r$, or $\{\vec x:x_j=0\}$ for a single $j$. The
profile and the coordinate are obtained by recursion on the term.
\end{theorem}

\begin{proof}
Induction on terms, with bound update variables fresh. $0$ is always zero, a
successor never, and $x_j$ exactly when $x_j=0$. Consider
$\mathsf{R}^{+}(n,a;v.h)$ with update $h(v,\vec x)$. By the induction
hypothesis the zero set of $h$ has one of four forms. If $h$ is never zero, the
iterator is never zero. If $h$ is always zero, the iterator is always zero. If
$h$ is zero exactly when a passive variable $x_j$ is zero, the iterator has the
same profile. These three cases use that at least one update is applied. If
$h$ is zero exactly when its state $v$ is zero, every update preserves the zero
status of the state, so the final state has the zero profile of the seed $a$,
which is classified by the induction hypothesis. The zero profile of the
counter is not used: the number of updates is positive. A variable that occurs
both in the counter and in the seed or the update keeps its role in the
respective profile.
\end{proof}

\begin{corollary}
\label{obs:cor:zeroprofile}
No term of this language denotes the predecessor (zero set $\{0,1\}$), the zero
test that is $1$ at $0$ and $0$ elsewhere (zero set $\mathbb{N}\setminus\{0\}$),
addition (zero set $\{(0,0)\}$), or the initialized iterator, whose
constant-update case $\It{n}{a}{v}{b}$ is the selector and gives the zero test
at $a=1$, $b=0$.
\end{corollary}

The definable shifted addition $n+a+1$ is never zero and does not contradict
the corollary.

\begin{remark}[Scope]
\label{obs:rem:zeroprofile-scope}
This is a standard-model term-definability obstruction. In intuitionistic
logic with coverage, injectivity of $\Suc$ and separation $0\ne\Suc m$, the
formula
\[
(n=0\wedge z=a)\ \vee\ \exists m\,\bigl(n=\Suc m\wedge z=\mathsf{R}^{+}(m,a;v.h)\bigr)
\]
defines a total single-valued graph of the initialized iteration of each
update $h$ of the old language. Separation is needed to exclude overlap of the
two branches. The observation does not give a bi-interpretation with $\iPRA$:
induction for the graph-defined expanded language and closure under nesting of
newly defined updates remain to be proved, and in the native calculus the
connectives of the graph formula are themselves reconstructed from a selector,
so assuming them at that point is circular. The invariant does not survive a
zero-test primitive: $\mathsf{R}^{+}(n,0;v.Z v)$ alternates between $1$ and
$0$, so zero status then depends on the iteration count.
\end{remark}

The motivation is a raw encoding in which a two-field numerical array $[a,h]$
has the value $a$ and a three-field array $[n,a,h]$ the value
$\mathsf{R}^{+}(n,a;v.h)$, which would allow a uniform predecessor expansion in
place of a zero/successor conditional. The theorem shows that this encoding
cannot inherit the term compiler or the selector-based logical bootstrap.

\section{Rule-level obstructions}
\label{obs:sec:rules}

This section draws on \pkg{research}, \pkg{adaptive\_seed},
\pkg{local\_scopes}, \pkg{counter\_congruence} and
\pkg{constructed\_computation}. Each result rejects one tempting deletion or
replacement of a rule by a reading that validates the remaining rules. Each
concerns the stated rule system with the identity translation of formulas.

\subsection{Atomic case introduction cannot be made automatic}

In $\Nat_{\mathrm{p}}$ a certificate chooses between the two introductions of
$\D{x}{y}$. The \emph{forced-progress} variant removes the choice: it uses
\rulename{DI-progress} when the matrix is a literal equation and \rulename{DI}
by cases when the matrix is compound. Equivalently, it deletes \rulename{DI}
by cases at literal-equation matrices.

\begin{proposition}
\label{obs:prop:forced-progress}
$\Nat_{\mathrm{p}}$ derives $\D{x}{y}(y=1)$, with $x$ unused. The
forced-progress variant does not.
\end{proposition}

\begin{proof}
In $\Nat_{\mathrm{p}}$, \rulename{DI} by cases needs
$0=1\proves\Suc z=1$. From $0=1$, transport through the selector gives
$s=\sel(0,s,t)=\sel(1,s,t)=t$ for all terms $s,t$, in particular $\Suc z=1$. No
induction on a compound formula is used.

For the variant, evaluate terms in $\mathbb{N}$, read equations as equality
and $\Q{y}B$ as $B(0)\to\forall y\,B(y)$, and read $\D{x}{y}$ by
\[
(\D{x}{y}B)^{\dagger}=
\begin{cases}
\exists x\,\forall y\,\bigl(B^{\dagger}(x,y)\to B^{\dagger}(x,\Suc y)\bigr), & B \text{ a literal equation},\\
\exists x\,\bigl(B^{\dagger}(x,0)\to\forall y\,B^{\dagger}(x,y)\bigr), & B \text{ compound}.
\end{cases}
\]
The reading is defined on each fixed formula. Forced progress introduction is
sound by its direct meaning. Case introduction at a compound matrix is sound by
the zero-or-successor property of $\mathbb{N}$. Paired elimination at an atomic
matrix is sound because induction in $\mathbb{N}$ turns the supplied
progressiveness into the conditional closure $\Q{y}B(a,y)$; at a compound
matrix it is sound directly. \rulename{QE}, \rulename{Hyp} and the computation
equations are sound. Transport is sound because numerical substitution does not
change whether a matrix is a literal equation. Under $\dagger$ the formula
$\D{x}{y}(y=1)$ says $\exists x\,\forall y\,(y=1\to\Suc y=1)$, which is false
at $y=1$.
\end{proof}

The padding equivalence between a formula and a gated paired closure does not
repair the counterexample: padding moves a proof obligation into a compound
matrix, and the elimination returns only the conditional closure, not the
progressiveness needed to introduce a raw atomic paired formula. A conditional
closure can hold because its zero antecedent is false, without its matrix being
progressive. A public language in which every atomic matrix is gated evades the
obstruction at the cost of new grammar and decoder obligations.

\begin{remark}
\pkg{research} rederives the proposition from an earlier draft on flat
computation certificates; \pkg{compact\_kernel} and \pkg{uniform\_predecessor}
restate the conclusion. For the root-determined introduction with unrestricted
transport the corresponding strictness is formalized as
\Lean{rdprf\_not\_cases} and \Lean{rdprf\_strictly\_weaker}, with the
progressive reading over the natural numbers as the sound reading.
\end{remark}

\subsection{Elimination majors must be proved, not only retained}

\begin{proposition}[\pkg{adaptive\_seed}]
\label{obs:prop:majors}
$\Nat_{\mathrm{p}}$ derives $\Q{y}(0=0)$. Restrict $\Nat_{\mathrm{p}}$ so that
the major premise of \rulename{QE} and of \rulename{DE} and the equation of a
transport must be members of the retained context, keep \rulename{Comp}, the
two introductions and, optionally, direct induction at equation goals, and add
no cut rule. The restricted calculus does not derive $\Q{y}(0=0)$.
\end{proposition}

\begin{proof}
In $\Nat_{\mathrm{p}}$, \rulename{DI} by cases with the witness $0$ derives
$\D{x}{y}(0=0)$ from $0=0\proves 0=0$; \rulename{DE} applied to this derived
major supplies $\Q{y}(0=0)$, and \rulename{Hyp} concludes.

In the restricted calculus consider a goal that is a closure $\Q{y}B$ in a context
of equations only, initially empty. \rulename{Hyp} does not close it. An
introduction concludes a paired formula. \rulename{Comp} adds an equation and
keeps the goal; induction applies to equation goals only. Transport along a
retained equation leaves a minor that is again a closure, in the same context.
\rulename{QE} and \rulename{DE} need a retained closure or paired formula, and
there is none. Induction on the height of a derivation shows that no derivation
exists.
\end{proof}

The proposition concerns this deletion with closures as public formulas and the
paired-introduction form; in $\Nat$ the formula is immediate by \rulename{QI}.
It does not show that the closure binder must remain an object connective, and an
unrestricted cut repairs the displayed theorem. The package verifies the
example by a mutation of a checker in which the recursive check of a major is
replaced by a membership test.

\subsection{Root replacement does not give congruence}

A further merger encodes an equation $s=t$ as a clause that sends $r=s$ to
$r=t$, so that equality transport and closure use become one clause
elimination. This yields replacement at the root of an equation only.

\begin{proposition}[\pkg{local\_scopes}]
\label{obs:prop:root-vs-congruence}
The computation equations, zero/successor cases, induction, reflexivity and
root replacement do not derive congruence of $\Suc$ for a primitive equation
predicate.
\end{proposition}

\begin{proof}
Take $\mathbb{N}$ with the standard functions and interpret the equation
predicate by the equivalence relation $\approx$ whose only non-singleton class
is $\{2,3\}$. Every computation equation holds because both sides have the same
value. Reflexivity and root replacement are sound for an equivalence relation.
Cases and induction are sound on $\mathbb{N}$ for every interpreted formula.
But $2\approx 3$ and $\Suc 2\not\approx\Suc 3$.
\end{proof}

The counter-model is not the equality of the complete theory, which is actual
equality. \Cref{obs:thm:cc-congruence} shows what must be added: congruence in
one fixed position, the counter of a recursor, suffices. A Leibniz clause
$P(s)\to P(t)$ schematic in a formula context also merges the two
eliminations, and moves the work into context application and higher-order
substitution; no smaller checker for that route is established.

\subsection{The witness order of the paired binder}

A reversed binder
$\mathsf{E}_{x,y}B\defeq\forall y\,\exists x\,(B(x,0)\to B(x,y))$ would merge
the two eliminations and remove the reusable clause. Consider the naive
introduction that fixes a witness term $t(z)$ in the instruction and checks
$B(t(z),0)\proves B(t(z),\Suc z)$, its progress variant with a witness
independent of $z$ at a literal equation, and the ordinary instance
elimination.

\begin{proposition}[\pkg{counter\_congruence}]
\label{obs:prop:reversed}
There is a literal equation matrix $B$ such that the ordinary reading
$\forall y\,\exists x\,(B(x,0)\to B(x,y))$ is true in $\mathbb{N}$ and
$\mathsf{E}_{x,y}B$ is false under a reading that validates the naive rules.
Hence the naive rules do not derive the equivalence of every $\mathsf{E}$
formula with its ordinary reading.
\end{proposition}

\begin{proof}
Read $\mathsf{E}_{x,y}B$, with free parameters fixed, as: there is a unary
primitive recursive $f$ with $\forall y\,(B(f(y),0)\to B(f(y),y))$. The naive
case rule is sound with $f(0)=0$ and $f(\Suc z)=t(z)$; the progress rule is
sound with a constant $f$; instance elimination and transport are sound. The
gated macros for the connectives keep their two-valued meanings: an existential
macro needs a constant selector, and the universal macro has an unused
existential coordinate.

Let $F$ be a total computable function that is not primitive recursive and let
$T(n,c)$ be the primitive recursive predicate saying that $c$ is a halting
computation trace of a fixed program for $F$ on the input $n$. Choose a literal
equation $B$ with $B(x,0)\equiv\top$ and $B(x,\Suc n)\equiv T(n,x)$, by a
primitive recursive characteristic and a case selector. The ordinary reading is
true. The selector reading is false: a primitive recursive $f$ would give the
trace $f(\Suc n)$, and the primitive recursive output decoder would compute
$F(n)$.
\end{proof}

The source witness obtained at a fresh universal parameter may have to be
chosen inside the proof scope of that parameter, after an existential
elimination; a witness fixed before the subproof does not allow this
interleaving. The proposition excludes this rule package with this reading. A
repair with delayed witness choice or an intermediate introduction judgment is
not excluded.

\subsection{Two side conditions that cannot be dropped}

\begin{example}[Freshness of the step variable for the witness]
\label{obs:ex:fresh}
The formula $\D{x}{y}\bigl(\It{y}{0}{v}{\Suc x}=y\bigr)$ is false in
$\mathbb{N}$: its zero instance is true, and its reading requires every
positive number to equal one fixed $\Suc x$. If an introduction could take the
step variable $z$ as its witness, the step
$\It{\Suc z}{0}{v}{\Suc z}=\Suc z$ would be a computation instance. The
condition that $z$ is fresh for the witness is therefore necessary
(\pkg{local\_scopes}).
\end{example}

\begin{example}[Local assumptions do not escape a major]
\label{obs:ex:leak}
Case introduction derives $\D{x}{y}(0=1)$ while locally assuming $0=1$. A
checker that lets this assumption leak from the subproof of a major into the
context of the caller accepts $0=1$ (\pkg{research}).
\end{example}

\pkg{constructed\_computation} proves, by a finite-range invariant, that
restricting the witnesses of \rulename{DI} to $0$ or variables and at the same
time the targets of \rulename{QE} to $1$, to finitely many numerals, or to
variables loses the theorem $\D{x}{y}(\sel(y,N+1,x)=N+1)$ for a suitable $N$;
the statement and its proof are given with the specialization rules in another
chapter.

\section{Independence of the progress rule}
\label{obs:sec:progress-independence}

This section draws on \pkg{progress\_independence} and, in
\cref{obs:sec:fixed-points}, on \pkg{ipra\_fixed\_points\_and\_checker}, and in
\cref{obs:sec:barriers} on \pkg{iteration\_barriers}. The
first package has two independent write-ups of the same result; the section follows the first, takes
the explicit derivation of the positive side and the extensional corollary
from the second, and compares the two constructions in
\cref{obs:rem:pi-saturation}. It shows that
\rulename{QI} by progress is not admissible over the other seven rules of the
core. Write $\Nat^{-}$ for $\Nat$ without \rulename{QI}.

The README of the repository records three facts that shape a countermodel.
Each of the other seven rules is independent of the rest, by readings of the
language over the natural numbers (\Lean{NativeModels}). Such a reading cannot
work for \rulename{QI}: over the natural numbers the soundness conditions for
\rulename{DI} and \rulename{DE} already imply the condition for \rulename{QI},
whatever an equation means (\Lean{Sem.Sound.qIInd\_of}); so the numerical
domain must be nonstandard. And $\Nat^{-}$ derives $0=1$ from every successor
cycle of odd length (\Lean{NoInd.no\_odd\_cycle}), so the model has no such
cycle; \cref{obs:thm:no-cycle} extends this to every length. The structures
used earlier in the chapter do not serve. Every power $\mathcal{N}^{(J)}$
satisfies induction for literal equations (\cref{obs:cor:power-axioms}), so
\rulename{QI} is sound in it; what fails there is coverage, hence \rulename{DI}
by cases (\cref{obs:prop:coverage-fails}). An ultrapower of the standard model
with its internal iterator satisfies induction as well, by the theorem of
\L o\'s. The countermodel has to do the opposite of a power: keep coverage and
both computation equations at every count, nonstandard counts included, and
lose induction for one equation. Coverage and acyclicity force the nonstandard
part of the successor graph to consist of chains infinite in both directions,
so every step function that is definable with parameters needs an orbit
infinite in both directions along each chain.

All arguments of this section are carried out in classical metamathematics with
the axiom of choice. This adds no classical axiom to the object calculus: a
classical structure is a sound test of intuitionistic derivability
(\cref{obs:sec:kinds}).

\subsection{Coherent iterator structures}

\begin{definition}
\label{obs:def:pi-coherent}
An \emph{iterator language} is a first-order language $L$ with $0$, $\Suc$ and
iterator symbols $I_h(n,a,\vec p)$, one for each member $h(u,\vec p)$ of a set
of \emph{descriptors}: terms of $L$ with a designated state variable $u$ and a
finite tuple of passive variables $\vec p$. The variable $u$ need not occur in
$h$. The language is built in stages, so that the body of an iterator symbol is
a term of an earlier stage. An $L$-structure is \emph{coherent} if it satisfies
the following first-order sentences.
\begin{enumerate}[label=(\alph*)]
\item \emph{Successor conditions.} $\Suc$ is injective, $0$ is not a successor,
every element is $0$ or a successor, and $\Suc^{k}x\ne x$ for each $k\ge 1$.
\item \emph{Computation.} For every iterator symbol,
$I_h(0,a,\vec p)=a$ and
$I_h(\Suc n,a,\vec p)=h(I_h(n,a,\vec p),\vec p)$.
\item \emph{Cross-extensionality.} For every two iterator symbols $I_h$, $I_g$
of $L$, the case $h=g$ included,
\[
\forall\vec p\,\forall\vec q\,\Bigl[\forall u\,\bigl(h(u,\vec p)=g(u,\vec q)\bigr)\to
\forall n\,\forall a\,\bigl(I_h(n,a,\vec p)=I_g(n,a,\vec q)\bigr)\Bigr].
\tag{Ext}
\]
\end{enumerate}
\end{definition}

By (Ext) an iterator depends on the unary function $u\mapsto h(u,\vec p)$ and
on its two numerical inputs, not on the presentation of the function;
repetition, omission and reordering of parameters are included. An independent
choice of an orbit for each written body would not have this property:
substitution of terms for the passive parameters changes the written body
without changing the function, and the substitution lemma for the native
syntax would fail (\cref{obs:lem:pi-subst}).

\begin{lemma}[Orbits in one ultrapower]
\label{obs:lem:pi-orbit}
Let $A$ be a nonempty $L$-structure, $\mathcal{U}$ a nonprincipal ultrafilter
on $\mathbb{N}$ and $B=A^{\mathbb{N}}/\mathcal{U}$. For every term
$h(u,\vec p)$ of $L$ and every tuple $\vec p$ of $B$, the function
$f(v)=h^{B}(v,\vec p)$ has a two-sided orbit: a family $(x_j)_{j\in\mathbb{Z}}$
of elements of $B$ with $f(x_j)=x_{j+1}$ for all $j$.
\end{lemma}

\begin{proof}
Let $\vec p=[\vec p_i]$, put $f_i(v)=h^{A}(v,\vec p_i)$ and choose any
$a_i\in A$. For $j\in\mathbb{Z}$ let
$x_j=\bigl[f_i^{\max(0,i+j)}(a_i)\bigr]$, with finite external iteration in
each coordinate. Term functions of $B$ are computed coordinatewise, and
$\max(0,i+j)+1=\max(0,i+j+1)$ for all $i\ge -j$, a cofinite set of
coordinates. Hence $f(x_j)=x_{j+1}$.
\end{proof}

The lemma concerns functions defined by terms with parameters. No induction
axiom of $A$ is used, and the orbit need not pass through a prescribed
element. In the structure $\mathfrak{M}_0$ of \cref{obs:lem:pi-seed} the
function $u\mapsto\Suc(I_{\mathrm{id}}(u,u))$ maps everything into the standard
ray and has no two-sided orbit; this is why the ultrapower precedes every
adjunction of iterators.

\begin{lemma}[Chains]
\label{obs:lem:pi-chains}
In a structure with the successor conditions, the elements $\Suc^{k}0$,
$k\in\mathbb{N}$, are pairwise distinct and form one component of the successor
graph, the \emph{standard ray}; every other component is a copy of
$(\mathbb{Z},+1)$.
\end{lemma}

\begin{proof}
Every nonzero element has exactly one predecessor and $0$ has none. An element
outside the standard ray does not reach $0$ after finitely many predecessor
steps, so its iterated predecessors and successors are all defined; they are
pairwise distinct by injectivity and the absence of cycles. The same two
properties make the numerals distinct.
\end{proof}

\begin{lemma}[Extension of a coherent structure]
\label{obs:lem:pi-extension}
Let $A$ be a coherent $L$-structure, $B=A^{\mathbb{N}}/\mathcal{U}$ as in
\cref{obs:lem:pi-orbit}, and let $\mathcal{D}$ be a set of terms of $L$ with
designated state variables that have no iterator symbol in $L$. Then $B$ is a
coherent elementary extension of $A$, and $B$ has an expansion to
$L\cup\{I_h:h\in\mathcal{D}\}$ that is coherent for the enlarged language. The
expansion changes no element and no function of $B$.
\end{lemma}

\begin{proof}
The conditions of \cref{obs:def:pi-coherent} are first-order, so they hold in
$B$ by the theorem of \L o\'s, also for the parameter tuples that are new in
$B$, and the diagonal embedding is elementary.

By \cref{obs:lem:pi-chains} the successor graph of $B$ is the standard ray
together with copies of $\mathbb{Z}$; choose an origin $b$ in each copy. Call a
function $f\colon B\to B$ \emph{presented} if $f=h^{B}(-,\vec p)$ for a term
$h$ of $L$ and a tuple $\vec p$ of $B$, and \emph{old} if $h$ can be taken with
an iterator symbol in $L$. Assign to every presented function $f$ one map
$\mathit{It}_f\colon B\times B\to B$ as follows.

If $f$ is old, say $f=g^{B}(-,\vec q)$ with $I_g$ in $L$, put
$\mathit{It}_f(n,a)=I_g^{B}(n,a,\vec q)$. This does not depend on the choice of
$g$ and $\vec q$, by (Ext) in $B$.

If $f$ is not old, put $\mathit{It}_f(\Suc^{k}0,a)=f^{k}(a)$ on the standard
ray. For each copy of $\mathbb{Z}$ with origin $b$ and each seed $a$ choose a
two-sided orbit $(x_j)_{j\in\mathbb{Z}}$ of $f$ by \cref{obs:lem:pi-orbit} and
put $\mathit{It}_f(\Suc^{j}b,a)=x_j$, where negative powers are iterated
predecessors. The axiom of choice provides the choices.

In both cases $\mathit{It}_f(0,a)=a$ and
$\mathit{It}_f(\Suc n,a)=f(\mathit{It}_f(n,a))$: for old $f$ by the computation
axioms of $I_g$, on the standard ray by the definition, and on a copy of
$\mathbb{Z}$ because $x_{j+1}=f(x_j)$. No condition links the values on a copy
of $\mathbb{Z}$ to the seed, and none is required.

For $h\in\mathcal{D}$ define
$I_h(n,a,\vec p)=\mathit{It}_f(n,a)$ with $f=h^{B}(-,\vec p)$. The computation
axioms of the new symbols hold. (Ext) holds for every pair of symbols of the
enlarged language: for two old symbols it holds in $B$ already, and in the
remaining cases two instances with the same function $f$ are both equal to
$\mathit{It}_f$, an old instance by the first clause of the definition.
\end{proof}

\begin{lemma}[The seed]
\label{obs:lem:pi-seed}
Let $L_0=\{0,\Suc,I_{\mathrm{id}}\}$, where $I_{\mathrm{id}}$ is the iterator
symbol of the descriptor $u$ without parameters. Let
$\mathfrak{M}_0=\mathbb{N}\sqcup\mathbb{Z}$, with $0$ and $1=\Suc 0$ in the
first summand, the ordinary successor on each summand, and
\[
I_{\mathrm{id}}(n,a)=
\begin{cases}
a, & n\in\mathbb{N},\\
1, & n\in\mathbb{Z}.
\end{cases}
\]
Then $\mathfrak{M}_0$ is coherent, and
$I_{\mathrm{id}}(c,0)=1\ne 0$ for every $c$ in the second summand.
\end{lemma}

\begin{proof}
The successor conditions are evident. $I_{\mathrm{id}}(0,a)=a$, and
$I_{\mathrm{id}}(\Suc n,a)=I_{\mathrm{id}}(n,a)$ because $\Suc$ preserves the
summands. (Ext) for the one symbol without parameters is trivial. No predicate
that distinguishes the summands is in the language.
\end{proof}

\begin{proposition}[The tower]
\label{obs:prop:pi-tower}
Let $L_{i+1}$ be $L_i$ together with an iterator symbol for every term of
$L_i$ with a designated state variable that has no symbol in $L_i$, and let
$L_\omega=\bigcup_i L_i$. There is a coherent $L_\omega$-structure
$\mathfrak{M}_\omega$ that extends $\mathfrak{M}_0$, in which every term of
$L_\omega$ with a designated state variable has an iterator symbol, and in
which $I_{\mathrm{id}}(c,0)=1\ne 0$ for the elements $c$ of the second summand
of $\mathfrak{M}_0$.
\end{proposition}

\begin{proof}
Given the coherent $L_i$-structure $\mathfrak{M}_i$, let $\mathfrak{M}_{i+1}$
be the expansion to $L_{i+1}$ of its ultrapower given by
\cref{obs:lem:pi-extension}; the ultrapower is taken in the whole language
$L_i$, and only descriptors over $L_i$ receive symbols at this stage, so no
body contains a symbol that is being introduced. Replace each ultrapower by an
isomorphic copy that contains the preceding structure. Then for every $j$ the
reducts $(\mathfrak{M}_i|L_j)_{i\ge j}$ form an elementary chain, and by the
elementary chain theorem the reduct to $L_j$ of the union
$\mathfrak{M}_\omega$ is an elementary extension of each of them. Every
successor condition, every computation axiom and every instance of (Ext)
mentions finitely many symbols, hence is a sentence of some $L_j$ that holds
in $\mathfrak{M}_j$, and so holds in $\mathfrak{M}_\omega$. The equation
$I_{\mathrm{id}}(c,0)=1$ and the inequation $1\ne 0$ hold in $\mathfrak{M}_0$
and are preserved. A term of $L_\omega$ is a term of some $L_i$ and has a
symbol in $L_{i+1}$.
\end{proof}

\subsection{Native terms in the tower}

Translate native terms into $L_\omega$ by $x^{*}=x$, $0^{*}=0$,
$(\Suc t)^{*}=\Suc t^{*}$ and
\[
\It{n}{a}{u}{h}^{*}=I_{h^{*}}(n^{*},a^{*},\vec p),
\]
where $\vec p$ lists the free variables of $h$ other than $u$ in a fixed order.
Inner iterators are translated first, so $h^{*}$ is a term of some $L_i$ and
its symbol belongs to $L_{i+1}$; the number of stages is the nesting depth of
iterators. Parameters are numerical variables; there is no sort of functions
and no internal evaluator. The value $[\![t]\!]_\rho$ of a native term under an
assignment $\rho$ in $\mathfrak{M}_\omega$ is the value of $t^{*}$. The body
$u$ is the descriptor of $I_{\mathrm{id}}$, so
$[\![\It{y}{0}{u}{u}]\!]_\rho=I_{\mathrm{id}}(\rho(y),0)$. % ed.: the report interprets a native body by any staged descriptor that presents its function and appeals to (Ext); for the body u this is the symbol of the seed.

\begin{lemma}[Substitution]
\label{obs:lem:pi-subst}
For native terms $t,s$, capture-avoiding substitution and every assignment
$\rho$ in $\mathfrak{M}_\omega$,
\[
[\![t\subst{s}{x}]\!]_\rho=[\![t]\!]_{\rho[x\mapsto[\![s]\!]_\rho]},
\]
and the value of a term does not change under renaming of bound variables.
The same holds for native formulas under the ordinary readings of
$\Q{y}$ and $\D{x}{y}$ with actual equality.
\end{lemma}

\begin{proof}
Induction on $t$. The cases of variables, $0$ and $\Suc$ are immediate. Let
$t=\It{n}{a}{u}{h}$ with $u$ renamed away from $x$ and $s$. The translation of
$t\subst{s}{x}$ is $I_{g}(\dots,\vec q)$ for the descriptor
$g=(h\subst{s}{x})^{*}$, whose passive variables $\vec q$ are those of $h$
other than $x$ together with those of $s$; in general $g$ is a different
descriptor from $h^{*}$. By the induction hypothesis the counter and the seed
have the required values, and for every element $v$
\[
[\![h\subst{s}{x}]\!]_{\rho[u\mapsto v]}=[\![h]\!]_{\rho[x\mapsto[\![s]\!]_\rho][u\mapsto v]},
\]
so the unary function of $g$ at the parameters $\rho(\vec q)$ equals the unary
function of $h^{*}$ at the parameters given by
$\rho[x\mapsto[\![s]\!]_\rho]$. (Ext) for the pair $I_g$, $I_{h^{*}}$ gives
equal iterator values. Renaming of the bound variable is the same argument
with equal functions. The statement for formulas follows by induction on
formulas.
\end{proof}

This is the point at which independent choices for written bodies would be
insufficient: the two descriptors are syntactically unrelated, and only (Ext),
imposed for all parameter tuples and preserved through all stages, identifies
their iterators.

\subsection{The theorem}

\begin{theorem}[Progress is not admissible]
\label{obs:thm:progress-independence}
Let $E(y)\defeq\bigl(\It{y}{0}{u}{u}=0\bigr)$.
\begin{enumerate}[label=(\roman*)]
\item $\Nat^{-}$ derives $E(0)$ and the sequent $E(z)\proves E(\Suc z)$.
\item Every sequent derivable in $\Nat^{-}$ is valid in $\mathfrak{M}_\omega$
under the ordinary readings, with actual equality.
\item $\Q{y}E(y)$ is false in $\mathfrak{M}_\omega$, and so is
$\D{x}{y}E(y)$ with $x$ unused.
\end{enumerate}
Hence $\Nat^{-}$ does not derive $\Q{y}E(y)$, which $\Nat$ derives by one
application of \rulename{QI}: the rule \rulename{QI} is not admissible over the
other seven rules, and \rulename{DI-progress} is not admissible over the other
rules of the paired-introduction form. The countermodel is a classical
structure with actual equality in which the successor is injective, coverage
and separation hold, there is no successor cycle, and both computation
equations hold for every finite iterator term at every count.
\end{theorem}

\begin{proof}
(i) $E(0)$ is the zero computation. For the step write
$a\defeq\It{z}{0}{u}{u}$ and $b\defeq\It{\Suc z}{0}{u}{u}$, and work in the
context $E(z)$, that is $a=0$. The derivation uses only primitive rules:
\begin{enumerate}[label=(\arabic*)]
\item \rulename{Hyp}: $a=0$.
\item[(2), (3)] The successor computation, twice: $b=a$, since the body $u$
with $a$ substituted for $u$ is $a$.
\setcounter{enumi}{3}
\item \rulename{Tr} with $s=b$, $t=a$, the constant template $h(v)=b$ and
$r=a$; both premises, $s=t$ and $h(s)=r$, are the equation $b=a$ of (2) and
(3), and the conclusion $r=h(t)$ is $a=b$.
\item \rulename{Tr} with $s=a$, $t=0$, the identity template $h(v)=v$ and
$r=b$; the premises are $a=0$ from (1) and $a=b$ from (4), and the conclusion
$r=h(t)$ is $b=0$, that is $E(\Suc z)$.
\end{enumerate}
One application of \rulename{QI} to (5) gives $\Q{y}E(y)$ in $\Nat$. As a
proof tree this has six nodes: one \rulename{Hyp} leaf, two successor
computation leaves, two \rulename{Tr} nodes and one \rulename{QI} node. No
derived equality rule and no induction lemma is used; derived reflexivity,
symmetry and transitivity (\cref{base:sec:structural}) are not needed.

(ii) A sequent is valid if every assignment that satisfies its context
satisfies its conclusion. \rulename{Hyp} is immediate. The zero computation is
the computation axiom for $I_{h^{*}}$. For the successor computation, the left
side has the value $h^{*}(I_{h^{*}}(n,a,\vec p),\vec p)$ by the computation
axiom, and the right side $h\subst{\It{n}{a}{u}{h}}{u}$ has the same value by
\cref{obs:lem:pi-subst}. Transport, from $s=t$ and $h(s)=r$ to $r=h(t)$, is
sound because equality is actual and
$[\![h(s)]\!]_\rho=[\![h(t)]\!]_\rho$ by \cref{obs:lem:pi-subst}.
\rulename{QE} is modus ponens and instantiation, with
\cref{obs:lem:pi-subst} for the instance $B(t)$. \rulename{DE} is existential
elimination: the value of the fresh variable is the witness, and freshness
keeps the context and the conclusion unchanged. For \rulename{DI} by cases let
$\rho$ satisfy $\Gamma$ and take the value of $t$ as witness. Suppose $B(t,0)$
holds under $\rho$. Every element is $0$ or $\Suc d$ for some $d$; in the
second case $\rho[z\mapsto d]$ still satisfies $\Gamma$ and $B(t,0)$ and gives
$t$ the same value, since $z$ is fresh, so the premise gives $B(t,\Suc z)$
there. Hence $B(t,y)$ holds for every $y$. All seven arguments are first-order
soundness arguments in one structure; coverage is used for \rulename{DI} only.

(iii) By the translation, $E(y)$ says $I_{\mathrm{id}}(y,0)=0$. It holds at
$y=0$ and fails at $y=c$ by \cref{obs:prop:pi-tower}. So
$E(0)\to\forall y\,E(y)$ is false, and with an unused witness coordinate so is
$\exists x\,(E(0)\to\forall y\,E(y))$.
\end{proof}

With the seven independence results recorded in the README
(\Lean{NativeModels}), each of the eight rule families of $\Nat$ is independent
of the other seven.

The model also satisfies extensionality in the update, stated in the native
term language. This strengthens the list of laws at the end of
\cref{obs:thm:progress-independence} from the computation equations to the
whole first-order theory of extensional computation.

\begin{corollary}[Extensional computation does not give induction]
\label{obs:cor:pi-extensional}
Let $T_{\mathrm{ext}}$ be the classical first-order theory with equality over
the native terms whose axioms are: $\Suc x=\Suc y\to x=y$, $\Suc x\ne 0$,
$x=0\vee\exists y\,x=\Suc y$, $\Suc^{k}x\ne x$ for each $k\ge 1$, the two
computation equations for every iterator term, and, for every two native bodies
$h$ and $g$ with the free variables $\vec p$ of $h$ and $\vec q$ of $g$ other
than $u$,
\[
\forall\vec p\,\forall\vec q\,\Bigl[\forall u\,(h=g)\to
\forall n\,\forall a\,\bigl(\It{n}{a}{u}{h}=\It{n}{a}{u}{g}\bigr)\Bigr].
\]
Then $T_{\mathrm{ext}}$ does not prove
$E(0)\wedge\forall y\,\bigl(E(y)\to E(\Suc y)\bigr)\to\forall y\,E(y)$ for
$E(y)\defeq\bigl(\It{y}{0}{u}{u}=0\bigr)$. A fortiori no intuitionistic
subtheory of $T_{\mathrm{ext}}$ proves it.
\end{corollary}

\begin{proof}
All axioms hold in $\mathfrak{M}_\omega$ under the translation of native terms.
The successor laws and coverage are successor conditions of
\cref{obs:def:pi-coherent}, and the computation equations were verified in
part (ii) of the proof of \cref{obs:thm:progress-independence}. For the
extensionality axiom fix values of $\vec p$ and $\vec q$ under an assignment
$\rho$. The value of $h$ at $\rho[u\mapsto v]$ is
$(h^{*})^{\mathfrak{M}_\omega}(v,\rho(\vec p))$, and similarly for $g$, so the
antecedent says that the unary functions of $h^{*}$ at $\rho(\vec p)$ and of
$g^{*}$ at $\rho(\vec q)$ coincide. The translations of the two iterator terms
are $I_{h^{*}}(n,a,\rho(\vec p))$ and $I_{g^{*}}(n,a,\rho(\vec q))$, and (Ext)
for the pair $I_{h^{*}}$, $I_{g^{*}}$ gives their equality. The displayed
induction implication is false in $\mathfrak{M}_\omega$: $E(0)$ holds, every
$E(y)\to E(\Suc y)$ holds because
$I_{\mathrm{id}}(\Suc y,0)=I_{\mathrm{id}}(y,0)$, and $E(c)$ fails by
\cref{obs:prop:pi-tower}.
\end{proof}

Equality in $\mathfrak{M}_\omega$ is actual, so transport of equals in
arbitrary formulas is also sound there; strengthening \rulename{Tr} to full
equality substitution leaves the countermodel intact.

\begin{remark}[What the construction needs]
\label{obs:rem:pi-needs}
Four requirements carry the proof, and dropping any of them leaves a gap. The
iterator of a function that is not old is chosen once for the function, not for
a written body. (Ext) is imposed for all pairs of symbols and all parameter
tuples, so that it transfers to the parameters of later ultrapowers. Each
ultrapower is taken in the whole current language, so that the earlier
iterators and their coherence are carried along. And a stage adds symbols only
for descriptors over the preceding language, so that the orbit lemma applies to
them. The construction makes no assertion about iteration of arbitrary
functions; \cref{obs:lem:pi-orbit} is false for arbitrary maps.
\end{remark}

\begin{remark}[A countable model]
The languages $L_i$ and $L_\omega$ are countable. By the downward
L\"owenheim--Skolem theorem $\mathfrak{M}_\omega$ has a countable elementary
substructure that contains $c$; it satisfies the same sentences, hence is
coherent and refutes $E(c)$, and the soundness argument applies to it. This
does not contradict the obstruction to readings over the natural numbers: the
successor of the countable carrier has nonstandard chains and is not the map
$n\mapsto n+1$ under any enumeration that would make it a reading in that
sense. No computable presentation of the model is claimed. The second write-up
of the package proves countability by the same argument, and so does not add to
the statement here.
\end{remark}

\begin{remark}[The second construction: saturation in place of ultrapowers]
\label{obs:rem:pi-saturation}
The second write-up of \pkg{progress\_independence} builds the same kind of
model with three differences, none of which changes the theorem.
\begin{enumerate}[label=(\alph*)]
\item The seed is $\mathbb{N}\sqcup\mathbb{Z}$ with the identity iterator
$G(n,a)=a$ on the first summand and $G(n,a)=\Suc a$ on the second, and with a
constant $c$ for an element of the second summand. The computation equations
hold because $\Suc$ preserves the summands, and $G(c,0)=\Suc 0\ne 0$.
\item Two-sided orbits come from $\omega$-saturation instead of
\cref{obs:lem:pi-orbit}. Let $M$ be $\omega$-saturated and $f$ a unary function
defined by a term with finitely many parameters. The type
$\{\exists x\,f^{m}(x)=y:m\ge 1\}$ is finitely satisfiable, by $y=f^{N}(w)$ for
any $w$, so some $b_0$ realizes it. If $b_j$ realizes it, the type
$\{f(x)=b_j\}\cup\{\exists w\,f^{m}(w)=x:m\ge 1\}$ over the parameters of $f$
and $b_j$ is finitely satisfiable as well: for the requirements up to $N$
choose $w$ with $f^{N+1}(w)=b_j$ and put $x=f^{N}(w)$, so that
$x=f^{m}(f^{N-m}(w))$ for $m\le N$. A realization is $b_{j-1}$, again with
predecessors of every finite depth, and dependent choice gives all negative
indices; the positive indices are $b_{j+1}=f(b_j)$ for $j\ge 0$. The orbit
need not be injective.
\item Each stage takes an $\omega$-saturated elementary extension of the
preceding expansion, and the extension step groups the presented functions by
equality: a class that contains a function with an old iterator keeps it, which
is unambiguous by (Ext), and every other class receives one iterator
$J_f(\Suc^{m}0,a)=f^{m}(a)$ on the standard ray and $J_f(\Suc^{i}d_C,a)=b_i$ on
each copy $C$ of $\mathbb{Z}$ with origin $d_C$, for one orbit of $f$ chosen
independently of the seed. All templates without a symbol receive symbols
interpreted by these iterators, so (Ext) holds between old and new symbols and
among new ones, as in \cref{obs:lem:pi-extension}.
\end{enumerate}
The write-up states the outcome in the generality of
\cref{obs:lem:pi-extension} and \cref{obs:prop:pi-tower}: a structure in a
countable language with the successor conditions and a countable coherent
family of iterators has an elementary extension with an expansion by coherent
iterators for all finite-depth templates, preserving the old iterators. The
difference shows that two-sided orbits of term-defined functions are a
compactness phenomenon: every $\omega$-saturated structure has them, and the
ultrapower of \cref{obs:lem:pi-orbit}, which is $\aleph_1$-saturated for a
countable language, is one instance. The proof of the second write-up is
complete; each of its steps is checked above or coincides with a step of
\cref{obs:lem:pi-extension,obs:prop:pi-tower,obs:lem:pi-subst} and of the
soundness argument in the proof of \cref{obs:thm:progress-independence}.
\end{remark}

\begin{proposition}[A seed with cancellative addition]
\label{obs:prop:pi-cancellative}
Let
$M_0=\{(0,n):n\in\mathbb{N}\}\cup\{(k,z):k\ge 1,\ z\in\mathbb{Z}\}$ with
componentwise addition, $0=(0,0)$, $1=(0,1)$ and $\Suc(k,z)=(k,z+1)$. Let
$I_{\mathrm{id}}(n,a)$ be $a$ if the first coordinate of $n$ is $0$ and $1$
otherwise, and let $K(n,a)=a+n$ be the iterator symbol of the descriptor
$\Suc u$. Then $M_0$ is coherent for the language
$\{0,\Suc,I_{\mathrm{id}},K\}$, $K(n,0)=n$ holds, addition $a+n=K(n,a)$ is
associative, commutative and cancellative with the unit $0$, and
$I_{\mathrm{id}}((1,0),0)=1$. The tower over $M_0$ gives a countermodel as in
\cref{obs:thm:progress-independence} that satisfies these laws in addition.
\end{proposition}

\begin{proof}
$M_0$ is a submonoid of the group $\mathbb{Z}^2$, so addition is associative,
commutative and cancellative; it is conical, and $\Suc x=x+1$. The successor is
injective, omits $0$, has no cycle, and every nonzero element has a
predecessor in $M_0$. The computation axioms of $I_{\mathrm{id}}$ hold because
$\Suc$ preserves the first coordinate, and those of $K$ are $a+0=a$ and
$a+\Suc n=\Suc(a+n)$. (Ext) between the two symbols has a false antecedent,
since the identity and the successor are different functions. The displayed
laws are first-order sentences of the seed language and survive the tower.
\end{proof}

The stronger seed shows that the countermodel does not depend on a defect of
addition. It does not separate uniqueness of iteration from induction: in
$M_0$ the term $f(n)=n+I_{\mathrm{id}}(n,0)$ satisfies $f(0)=0$ and
$f(\Suc n)=\Suc f(n)$, and $f(c)=\Suc c\ne c=K(c,0)$ at $c=(1,0)$, so the
instance of the rule \rulename{UR} of \cref{obs:sec:partial} for the update
$\Suc u$ fails there. The instance of \rulename{UR} for the identity update
and the constant template $f(z)=0$ already excludes the witness of the theorem,
since its conclusion is $0=\It{t}{0}{u}{u}$.

\subsection{Fixed points and seed shift do not give induction}
\label{obs:sec:fixed-points}

This subsection draws on the mathematical part of
\pkg{ipra\_fixed\_points\_and\_checker}. The models of
\cref{obs:thm:progress-independence,obs:prop:pi-cancellative} refute the
identity-iterator scheme $\It{n}{a}{u}{u}=a$. The subsection refines the
extension step so that two laws of iteration survive the tower, both theorems
of $\iPRA$ proved by induction on the counter, and starts the tower from a seed
in which an induction instance fails whose base and step follow from these laws.
For a descriptor $h(u,\vec p)$ write $I_h(n,a,\vec p)$ as in
\cref{obs:def:pi-coherent}. The two schemes are
\[
\forall\vec p\,\forall a\,\Bigl[h(a,\vec p)=a\to
\forall n\,\bigl(I_h(n,a,\vec p)=a\bigr)\Bigr],
\tag{FP}
\]
\[
\forall\vec p\,\forall n\,\forall a\,\Bigl[
I_h\bigl(n,h(a,\vec p),\vec p\bigr)=h\bigl(I_h(n,a,\vec p),\vec p\bigr)\Bigr].
\tag{Shift}
\]
(FP) says that iteration preserves every fixed seed of its update; for the
body $u$ it is the scheme $\It{n}{a}{u}{u}=a$. (Shift) moves one application of
the update from the seed to the output. Both compare iterators of one update;
fusion, which compares iterators of different updates, is treated after the
theorem. A coherent structure is \emph{fixed-point coherent} if it satisfies
every instance of (FP) and (Shift) for its iterator symbols.

\begin{lemma}[Extension by orbit maps]
\label{obs:lem:fp-extension}
Let $A$ be a fixed-point coherent $L$-structure, and let
$B=A^{\mathbb{N}}/\mathcal{U}$ and $\mathcal{D}$ be as in
\cref{obs:lem:pi-extension}. Then $B$ is fixed-point coherent, and $B$ has an
expansion to $L\cup\{I_h:h\in\mathcal{D}\}$ that is fixed-point coherent for
the enlarged language. The expansion changes no element and no function of $B$.
\end{lemma}

\begin{proof}
(FP) and (Shift) are schemes of first-order sentences with all parameters
bound, so, like the conditions of \cref{obs:def:pi-coherent}, they hold in $B$
by the theorem of \L o\'s, also at the parameter tuples that are new in $B$.

\emph{Orbit maps.} Let $f=h^{B}(-,\vec p)$ be a presented function, with $h$ a
term of $L$ and $\vec p=[\vec p_i]$, and put $f_i=h^{A}(-,\vec p_i)$. For
$j\in\mathbb{Z}$ define $H_j\colon B\to B$ by
\[
H_j\bigl([a_i]\bigr)=\bigl[f_i^{\max(0,i+j)}(a_i)\bigr],
\]
with finite external iteration in each coordinate. If $a_i=a'_i$ on a member
of $\mathcal{U}$, the finite iterates agree on that member, and the same holds
for a change of the representatives $\vec p_i$; so $H_j$ is a well-defined map.
It depends only on $f$: if also $f=g^{B}(-,\vec q)$, \L o\'s applied to
$\forall u\,h(u,\vec p)=g(u,\vec q)$ gives a member of $\mathcal{U}$ on which
$f_i$ and $g_i$ agree at every argument, hence so do their iterates of every
finite length. The maps have three properties.
\begin{enumerate}[label=(\alph*)]
\item $H_{j+1}=f\circ H_j$. Term functions of $B$ are computed coordinatewise,
and $\max(0,i+j)+1=\max(0,i+j+1)$ for all $i\ge -j$; this cofinite set of
coordinates is the same for every seed.
\item $H_j\circ f=f\circ H_j$, since finite iterates of $f_i$ commute with
$f_i$.
\item If $f(a)=a$, then $H_j(a)=a$: for $a=[a_i]$ the equation $f_i(a_i)=a_i$
holds on a member of $\mathcal{U}$, and there every finite iterate fixes $a_i$.
\end{enumerate}
For a fixed seed $a$ the family $(H_j(a))_{j\in\mathbb{Z}}$ is a two-sided orbit
of $f$ as in \cref{obs:lem:pi-orbit}; the difference is that one family of maps
serves all seeds at once.

\emph{The expansion.} Assign $\mathit{It}_f$ to every presented function $f$ as
in the proof of \cref{obs:lem:pi-extension}, with one change: if $f$ is not
old, put $\mathit{It}_f(\Suc^{k}0,a)=f^{k}(a)$ for $k\in\mathbb{N}$ and
$\mathit{It}_f(\Suc^{j}b,a)=H_j(a)$ for $j\in\mathbb{Z}$, for each copy of
$\mathbb{Z}$ with origin $b$ (\cref{obs:lem:pi-chains}), negative powers being
iterated predecessors. Define $I_h$ for $h\in\mathcal{D}$ as there. The
computation axioms and (Ext) hold by the argument of
\cref{obs:lem:pi-extension}: on a copy of $\mathbb{Z}$ property (a) replaces
$x_{j+1}=f(x_j)$, and (Ext) uses only that $\mathit{It}_f$ is assigned once for
each function.

\emph{(FP) and (Shift).} A symbol $I_h$ at the parameters $\vec p$ is
$\mathit{It}_f$ for $f=h^{B}(-,\vec p)$, and $h(a,\vec p)=f(a)$. It therefore
suffices to show, for every presented $f$ and all $n,a$,
\[
\mathit{It}_f\bigl(n,f(a)\bigr)=f\bigl(\mathit{It}_f(n,a)\bigr),\qquad
f(a)=a\ \Rightarrow\ \mathit{It}_f(n,a)=a .
\]
If $f$ is old, $f=g^{B}(-,\vec q)$ with $I_g$ in $L$ and
$\mathit{It}_f=I_g^{B}(-,-,\vec q)$, and the two statements are instances of
(Shift) and (FP) for $I_g$ at $\vec q$, which hold in $B$. If $f$ is not old,
on the standard ray $f^{k}(f(a))=f(f^{k}(a))$, and $f(a)=a$ gives $f^{k}(a)=a$;
at $n=\Suc^{j}b$ the two statements are (b) and (c).
\end{proof}

The choice of orbits in \cref{obs:lem:pi-extension}, made separately for each
seed, does not give (Shift): the orbits chosen at $a$ and at $f(a)$ are
unrelated. The ultrapower supplies the maps $H_j$ because the iteration length
$\max(0,i+j)$ depends on the coordinate only, not on the seed.

\begin{proposition}[The tower over the doubling seed]
\label{obs:prop:fp-tower}
Let $M_0$ be the carrier of \cref{obs:prop:pi-cancellative},
$M_0=\{(0,z):z\ge 0\}\cup\{(q,z):q\ge 1,\ z\in\mathbb{Z}\}$ with componentwise
addition, $0=(0,0)$ and $\Suc(q,z)=(q,z+1)$. Let $L_0=\{0,\Suc,J,K\}$, where
$J$ and $K$ are the iterator symbols of the descriptors $\Suc u$ and
$\Suc(\Suc u)$ without parameters, and put
\[
d(q,z)=(q,2z),\qquad J(n,a)=a+n,\qquad K(n,a)=a+d(n).
\]
Then $M_0$ is fixed-point coherent, it satisfies the sentences
\begin{equation}
\label{obs:eq:fp-monoid}
\begin{gathered}
\forall n\ J(n,0)=n,\qquad \forall a\ J(0,a)=a,\qquad
\forall a\,\forall b\ J(b,a)=J(a,b),\\
\forall a\,\forall b\,\forall c\ J(c,J(b,a))=J(J(c,b),a),\qquad
\forall a\,\forall b\,\forall c\ \bigl(J(c,a)=J(c,b)\to a=b\bigr),
\end{gathered}
\end{equation}
and $K(c,0)\ne J(c,c)$ at $c=(1,0)$. The tower of \cref{obs:prop:pi-tower}
over $M_0$, with \cref{obs:lem:fp-extension} in place of
\cref{obs:lem:pi-extension} at every stage, gives a fixed-point coherent
$L_\omega$-structure $\mathfrak{M}^{\mathrm{fp}}_\omega$ that extends $M_0$, in
which every term of $L_\omega$ with a designated state variable has an iterator
symbol, and which satisfies \eqref{obs:eq:fp-monoid} and $K(c,0)\ne J(c,c)$.
\end{proposition}

\begin{proof}
The successor conditions hold in $M_0$ as in the proof of
\cref{obs:prop:pi-cancellative}. The map $d$ keeps the first coordinate and
keeps the second coordinate nonnegative when the first is $0$, so it maps $M_0$
into $M_0$; moreover $d(0)=0$ and $d(\Suc n)=\Suc^{2}d(n)$. Since
$\Suc x=x+(0,1)$,
\[
J(0,a)=a,\quad J(\Suc n,a)=\Suc J(n,a),\qquad
K(0,a)=a,\quad K(\Suc n,a)=a+\Suc^{2}d(n)=\Suc^{2}K(n,a),
\]
which are the computation axioms for the bodies $\Suc u$ and $\Suc(\Suc u)$.
(Shift) holds: $J(n,\Suc a)=\Suc J(n,a)$ and $K(n,\Suc^{2}a)=\Suc^{2}K(n,a)$.
(FP) is vacuous, since $\Suc$ and $\Suc^{2}$ have no fixed point. The mixed
instances of (Ext) have a false antecedent, because $\Suc 0\ne\Suc^{2}0$; the
two self-instances have no parameters and are tautologies. The sentences
\eqref{obs:eq:fp-monoid} say that $J(n,0)=0+n=n$ and that $a+b=J(b,a)$ is the
addition of a submonoid of $\mathbb{Z}^{2}$, which is commutative, associative
and cancellative with unit $0$. Finally $K(c,0)=d(1,0)=(1,0)$ and
$J(c,c)=(2,0)$.

The proof of \cref{obs:prop:pi-tower} applies without change: every instance of
(FP) and (Shift) mentions finitely many symbols, so it is a sentence of some
$L_j$ that holds in $\mathfrak{M}_j$ and is preserved by the elementary chain.
The sentences \eqref{obs:eq:fp-monoid} and the inequation, with the parameter
$c$ of $M_0$, hold in $M_0$ and are preserved in the same way.
\end{proof}

Native terms are translated into $L_\omega$ as in the subsection on native
terms in the tower; the body $\Suc u$ translates to the descriptor of $J$, the
body $\Suc(\Suc u)$ to that of $K$, and the body $u$ receives its symbol at the
first stage. \Cref{obs:lem:pi-subst} and its proof hold for
$\mathfrak{M}^{\mathrm{fp}}_\omega$ unchanged, since they use only the
computation axioms and (Ext).

\begin{theorem}[Fixed points and seed shift do not give induction]
\label{obs:thm:fixed-points}
In the native term language put
\[
J(n,a)\defeq\It{n}{a}{u}{\Suc u},\qquad
K(n,a)\defeq\It{n}{a}{u}{\Suc(\Suc u)},\qquad
a+b\defeq J(b,a),\qquad
E(n)\defeq\bigl(K(n,0)=J(n,n)\bigr).
\]
\begin{enumerate}[label=(\roman*)]
\item Under the translation and with actual equality,
$\mathfrak{M}^{\mathrm{fp}}_\omega$ satisfies: $\Suc x=\Suc y\to x=y$,
$\Suc x\ne 0$, $x=0\vee\exists y\,x=\Suc y$ and $\Suc^{k}x\ne x$ for each
$k\ge 1$; both computation equations for every iterator term; the
extensionality axioms of \cref{obs:cor:pi-extensional}; for every native body
$h$, the universal closures of
\[
h\subst{a}{u}=a\to\It{n}{a}{u}{h}=a,\qquad
\It{n}{h\subst{a}{u}}{u}{h}=h\subst{\It{n}{a}{u}{h}}{u};
\]
the law $\It{n}{0}{u}{\Suc u}=n$; and the laws of a commutative cancellative
monoid with unit $0$ for $+$.
\item The computation equations, the second scheme of (i) for the body
$\Suc u$, and equality reasoning prove $E(0)$ and $E(n)\to E(\Suc n)$, without
induction.
\item $E(c)$ is false in $\mathfrak{M}^{\mathrm{fp}}_\omega$ for $c=(1,0)$.
\end{enumerate}
Consequently the classical first-order theory $T_{\mathrm{fix}}$ whose axioms
are the statements listed in (i) proves $E(0)$ and
$\forall n\,(E(n)\to E(\Suc n))$ but not $\forall n\,E(n)$. It does not prove
the induction instance
$E(0)\wedge\forall n\,(E(n)\to E(\Suc n))\to\forall n\,E(n)$, and neither does
any intuitionistic subtheory of $T_{\mathrm{fix}}$. Preservation of all fixed
seeds, which contains $\It{n}{a}{u}{u}=a$, seed shift for every update, the
canonical successor iterator and cancellative addition do not replace
induction for literal equations over computation, coverage and separation.
\end{theorem}

\begin{proof}
(i) The successor laws and coverage are successor conditions of
\cref{obs:def:pi-coherent}. The computation equations follow from the
computation axioms and \cref{obs:lem:pi-subst} as in part (ii) of the proof of
\cref{obs:thm:progress-independence}, and the extensionality axioms as in the
proof of \cref{obs:cor:pi-extensional}. For the two schemes fix an assignment
$\rho$ and let $f=(h^{*})^{\mathfrak{M}^{\mathrm{fp}}_\omega}(-,\rho(\vec p))$,
where $\vec p$ are the free variables of $h$ other than $u$. By
\cref{obs:lem:pi-subst} the value of $h\subst{a}{u}$ is $f(\bar a)$, where
$\bar a$ is the value of $a$, and the value of $\It{n}{a}{u}{h}$ is
$I_{h^{*}}(\bar n,\bar a,\rho(\vec p))$; the value of
$h\subst{\It{n}{a}{u}{h}}{u}$ is $f$ applied to it. The two schemes are
therefore the instances of (FP) and (Shift) for $I_{h^{*}}$ at
$\rho(\vec p)$. The terms $\It{n}{0}{u}{\Suc u}$ and $a+b$ translate to
$J(n,0)$ and $J(b,a)$, so the last two items are \eqref{obs:eq:fp-monoid}.

(ii) $E(0)$: the zero computations give $K(0,0)=0$ and $J(0,0)=0$, and
symmetry and transitivity give $K(0,0)=J(0,0)$. Step: the successor
computation for $J$ and the scheme for $\Suc u$ at the seed $n$ give
\[
J(\Suc n,\Suc n)=\Suc J(n,\Suc n)=\Suc\Suc J(n,n).
\]
The successor computation for $K$, the hypothesis $E(n)$ and congruence give
\[
K(\Suc n,0)=\Suc\Suc K(n,0)=\Suc\Suc J(n,n)=J(\Suc n,\Suc n).
\]
The seed-shift instance can be replaced by commutativity of $+$ and the
computation equation:
$J(n,\Suc a)=J(\Suc a,n)=\Suc J(a,n)=\Suc J(n,a)$.

(iii) The translation of $E(c)$ is $K(c,0)=J(c,c)$, which is false by
\cref{obs:prop:fp-tower}.

The classical structure $\mathfrak{M}^{\mathrm{fp}}_\omega$ is a model of
$T_{\mathrm{fix}}$ by (i), and a sound test of intuitionistic derivability
(\cref{obs:sec:kinds}). If $T_{\mathrm{fix}}$ or a subtheory proved the
induction instance, then with (ii) $T_{\mathrm{fix}}$ would prove
$\forall n\,E(n)$, which is false by (iii).
\end{proof}

\begin{corollary}[The native form]
\label{obs:cor:fp-native}
Let $\Nat^{-}_{\mathrm{fix}}$ be $\Nat^{-}$ with two further rule families, for
every native body $h$ and all terms $n$, $a$:
\[
\Gamma\proves\It{n}{h\subst{a}{u}}{u}{h}=h\subst{\It{n}{a}{u}{h}}{u}
\ \ \rulename{Shift},\qquad
\infrule{\Gamma\proves h\subst{a}{u}=a}{\Gamma\proves\It{n}{a}{u}{h}=a}
\ \ \rulename{FP}.
\]
With $E$ as in \cref{obs:thm:fixed-points}:
\begin{enumerate}[label=(\roman*)]
\item $\Nat^{-}_{\mathrm{fix}}$ derives $E(0)$ and the sequent
$E(z)\proves E(\Suc z)$.
\item Every sequent derivable in $\Nat^{-}_{\mathrm{fix}}$ is valid in
$\mathfrak{M}^{\mathrm{fp}}_\omega$ under the ordinary readings, with actual
equality.
\item $\Q{y}E(y)$ is false in $\mathfrak{M}^{\mathrm{fp}}_\omega$.
\end{enumerate}
Hence $\Nat^{-}_{\mathrm{fix}}$ does not derive $\Q{y}E(y)$, and \rulename{QI}
is not admissible over $\Nat^{-}_{\mathrm{fix}}$. In particular $\Nat^{-}$
extended by the leaves $\Gamma\proves\It{n}{a}{u}{u}=a$, which
$\Nat^{-}_{\mathrm{fix}}$ derives by \rulename{FP} from reflexivity
(\cref{base:lem:eq}(i)), does not derive $\Q{y}E(y)$.
\end{corollary}

\begin{proof}
(i) The equality rules of \cref{base:lem:eq} use only \rulename{Tr} and
\rulename{I0}, so they are available in $\Nat^{-}$. $E(0)$: two \rulename{I0}
leaves, symmetry and transitivity. For the step, in the context $E(z)$:
\rulename{IS} gives $K(\Suc z,0)=\Suc\Suc K(z,0)$; \rulename{Hyp} and
congruence in the template $\Suc\Suc v$ give
$\Suc\Suc K(z,0)=\Suc\Suc J(z,z)$; \rulename{IS} gives
$J(\Suc z,\Suc z)=\Suc J(z,\Suc z)$; the \rulename{Shift} leaf for the body
$\Suc u$, the counter $z$ and the seed $z$ gives
$J(z,\Suc z)=\Suc J(z,z)$, and congruence in the template $\Suc v$ gives
$\Suc J(z,\Suc z)=\Suc\Suc J(z,z)$; symmetry and transitivity give
$K(\Suc z,0)=J(\Suc z,\Suc z)$.

(ii) For the seven rules of $\Nat^{-}$ the argument of part (ii) of the proof of
\cref{obs:thm:progress-independence} applies verbatim; it uses only the
coherence of the structure, \cref{obs:lem:pi-subst} and coverage. The
\rulename{Shift} leaves are valid by \cref{obs:thm:fixed-points}(i). For
\rulename{FP}, if $\rho$ satisfies $\Gamma$, the premise gives
$h\subst{a}{u}=a$ at $\rho$, and the first scheme of
\cref{obs:thm:fixed-points}(i) gives $\It{n}{a}{u}{h}=a$ at $\rho$.

(iii) $\Q{y}E(y)$ reads $E(0)\to\forall y\,E(y)$; $E(0)$ is true and $E(c)$ is
false.
\end{proof}

\begin{remark}[Scope of the fixed-point model]
\label{obs:rem:fp-scope}
The step of \cref{obs:cor:fp-native}(i) uses the \rulename{Shift} leaf for the
body $\Suc u$, or, by the alternative in the proof of
\cref{obs:thm:fixed-points}, commutativity of $+$. Neither is shown derivable
in $\Nat^{-}$ extended by the identity scheme $\It{n}{a}{u}{u}=a$ alone. The
corollary shows that the conclusion $\Q{y}E(y)$ is underivable in that
extension, but not that the premise $E(z)\proves E(\Suc z)$ is derivable there,
so it does not by itself show that \rulename{QI} fails to be admissible over
$\Nat^{-}$ with the identity scheme; that is shown with another equation and
another model in \cref{obs:cor:qi-identity}. % ed.: the package reads its theorem as an answer to the question of adjoining the identity scheme; its step uses (Shift), which holds in the model but is not derived from the seven rules and the identity scheme, so the settlement is stated for the rule families FP and Shift, and the stronger reading is left open.
In $\mathfrak{M}^{\mathrm{fp}}_\omega$ the term $n+\It{n}{0}{u}{u}$ that
refutes the uniqueness rule for the successor update over the seed of
\cref{obs:prop:pi-cancellative} equals $n$ by (FP), so that refutation does not
apply; whether \rulename{SU}, the uniqueness rule for the update $\Suc u$
(\cref{kern:def:SU}), is valid in $\mathfrak{M}^{\mathrm{fp}}_\omega$ is not
known. The result is a separation of principles over the native scalar
interface under the ordinary reading. It is not a smaller axiomatization of
$\iPRA$ and not a lower bound on the size of a checker, and it does not exclude
a bi-interpretation through a different definable domain. The package checks
the seed algebra of \cref{obs:prop:fp-tower} on finitely many coordinates by a
script; the tower and the native interpretation rest on the written proof
alone, and no Lean or Rocq verification is claimed.
\end{remark}

The missing strength sits at the comparison of different updates. For unary
term templates $f$, $g$, $h$, which may have parameters, the \emph{fusion rule}
is
\[
\infrule{\Gamma\proves f(g(z))=h(f(z))}
{\Gamma\proves f\bigl(\It{t}{a}{u}{g(u)}\bigr)=\It{t}{f(a)}{v}{h(v)}}
\ \rulename{Fusion},
\]
where $z$ is not free in $\Gamma$, $t$, $a$ or the parameters, and the three
templates contain $z$ only through the displayed argument. It is read, like
\rulename{UR} (\cref{obs:sec:partial}), in intuitionistic logic with equality
and the two computation equations.

\begin{proposition}[Fusion and uniqueness]
\label{obs:prop:fusion}
Over intuitionistic logic with equality and the computation equations:
\begin{enumerate}[label=(\roman*)]
\item \rulename{UR} derives \rulename{Fusion};
\item \rulename{Fusion} and the scheme $\It{t}{0}{u}{\Suc u}=t$ derive
\rulename{UR};
\item \rulename{Fusion} derives both schemes of \cref{obs:thm:fixed-points}(i):
from $\Gamma\proves h(a)=a$ it derives $\Gamma\proves\It{t}{a}{u}{h(u)}=a$, and
it derives $h\bigl(\It{t}{a}{u}{h(u)}\bigr)=\It{t}{h(a)}{v}{h(v)}$.
\end{enumerate}
Hence, in the presence of $\It{t}{0}{u}{\Suc u}=t$, \rulename{Fusion} and
\rulename{UR} are interderivable.
\end{proposition}

\begin{proof}
(i) Assume the premise of \rulename{Fusion}. Let $w$ be a variable not free in
$\Gamma$, $t$, $a$ or the templates, and put
$r(n)\defeq f\bigl(\It{n}{a}{u}{g(u)}\bigr)$. The successor computation and
congruence give $r(\Suc w)=f\bigl(g(\It{w}{a}{u}{g(u)})\bigr)$. Substituting
$\It{w}{a}{u}{g(u)}$ for the variable $z$, which is not free in $\Gamma$, in the
premise gives
$f\bigl(g(\It{w}{a}{u}{g(u)})\bigr)=h\bigl(f(\It{w}{a}{u}{g(u)})\bigr)=h(r(w))$.
So $r(\Suc w)=h(r(w))$, and \rulename{UR} for the template $r$ gives
$r(t)=\It{t}{r(0)}{v}{h(v)}$. The zero computation gives $r(0)=f(a)$, and
congruence in the seed gives the conclusion of \rulename{Fusion}.

(ii) Take $g(u)\defeq\Suc u$ and $a\defeq 0$. The premise of \rulename{Fusion}
is then $f(\Suc z)=h(f(z))$, the premise of \rulename{UR}. Its conclusion is
$f\bigl(\It{t}{0}{u}{\Suc u}\bigr)=\It{t}{f(0)}{v}{h(v)}$, and the scheme and
congruence turn the left side into $f(t)$.

(iii) For fixed points take the constant template $f(w)\defeq a$, any $g$, for
instance $g(u)\defeq\Suc u$, and the seed $a$. The premise $a=h(a)$ follows
from $h(a)=a$ by symmetry, and the conclusion is $a=\It{t}{a}{v}{h(v)}$. For
seed shift take $f=g=h$: the premise $h(h(z))=h(h(z))$ is reflexivity, and the
conclusion is the displayed equation.
\end{proof}

With $\It{t}{0}{u}{\Suc u}=t$, which \rulename{UR} derives
(\cref{obs:lem:ur-basic}(iii)), the proposition and \cref{kern:thm:ur} show
that $\Teq$ with induction for literal equations replaced by \rulename{Fusion}
and this scheme has the same theorems as $\Teq$: \rulename{Fusion} follows in
$\Teq$ from \cref{kern:prop:uniq} and (i). Consequently no structure with actual
equality that satisfies the computation equations, coverage, separation, the
scheme $\It{t}{0}{u}{\Suc u}=t$ and every instance of \rulename{Fusion} refutes
an induction instance for a literal equation; an instance is satisfied if, at
each value of its parameters, the premise at all values of $z$ implies the
conclusion at all values of $t$ and $a$, and then \rulename{UR} is sound with
arbitrary contexts because $z$ is fresh. So (FP) and (Shift), which
\rulename{Fusion} derives, are strictly weaker over this background, and
\cref{obs:thm:fixed-points} is the witness. The model violates one mixed
instance: the doubling template $f(w)\defeq J(w,w)$ satisfies
$f(\Suc z)=\Suc\Suc f(z)$ by \cref{obs:thm:fixed-points}(ii), which is the
premise of \rulename{Fusion} for $g(u)\defeq\Suc u$, $h(v)\defeq\Suc(\Suc v)$
and $a\defeq 0$, and also the premise of \rulename{UR} for the update
$\Suc(\Suc u)$. With $f(0)=0$ and $\It{t}{0}{u}{\Suc u}=t$, both conclusions
read $J(t,t)=K(t,0)$, which fails at $c$. Adding this one instance excludes
the model; it is not shown to give induction. Neither \rulename{Fusion} nor the
scheme $\It{t}{0}{u}{\Suc u}=t$ can be omitted from part (ii) of
\cref{obs:prop:fusion}, even in the presence of induction against every fixed
standard numeral (\cref{obs:sec:barriers}).

\subsection{Numeral induction, numbering and fusion}
\label{obs:sec:barriers}

This subsection draws on \pkg{iteration\_barriers}. It
separates the two hypotheses of \cref{obs:prop:fusion}(ii) by two classical
structures with actual equality, coverage, separation and an injective
successor, both built over a nonstandard model of Peano arithmetic, and it
turns their automorphisms into obstructions to bi-interpretation with
$\iPRA$. The first structure refines the two-fiber model of
\cref{ctr:thm:zeroind}; the second has no counterpart earlier in the article.

\emph{The theories.} Let $\mathsf{B}$ be the ordinary theory over the native
terms with intuitionistic first-order logic with equality
(\cref{base:sec:teq}), the two computation equations, coverage and
separation; it is $\Teq$ without induction. Write $\overline{k}=\Suc^{k}0$.
$T_{\mathrm{num}}$ is $\mathsf{B}$ with $\Suc x=\Suc y\to x=y$ and, for every
term $t(n,\vec p)$ and every $k\in\mathbb{N}$ fixed in the metatheory, the
numeral induction rule
\[
\infrule{\Gamma\proves t(0,\vec p)=\overline{k}\qquad
\Gamma,\ t(z,\vec p)=\overline{k}\proves t(\Suc z,\vec p)=\overline{k}}
{\Gamma\proves t(s,\vec p)=\overline{k}}
\]
with $z$ fresh for $\Gamma$, $s$ and $\vec p$; there is no predicate for
standardness in the language. Consider the laws
\[
\It{n}{a}{u}{u}=a\quad(\mathsf{Id}),\qquad
\It{n}{\overline{k}}{u}{\Suc u}=\Suc^{k}n\quad(\mathsf{N}_k,\ k\in\mathbb{N}),
\]
and the extensionality axioms $(\mathsf{Ext})$ of \cref{obs:cor:pi-extensional}.
$\mathsf{N}_0$ is the scheme $\It{t}{0}{u}{\Suc u}=t$ of
\cref{obs:prop:fusion}(ii). Put
\[
T_{\mathrm{cal}}\defeq T_{\mathrm{num}}+\mathsf{Id}+\{\mathsf{N}_k:k\in\mathbb{N}\}+\mathsf{Ext},
\qquad
T_{\mathrm{fus}}\defeq T_{\mathrm{num}}+\mathsf{Id}+\mathsf{Ext}+\rulename{Fusion},
\]
with \rulename{Fusion} as displayed before \cref{obs:prop:fusion}. All these
axioms and rules are derivable in $\Teq$: $\mathsf{Id}$, $\mathsf{N}_k$ and
$\mathsf{Ext}$ by induction on the displayed equation, \rulename{Fusion} by
\cref{kern:prop:uniq} and \cref{obs:prop:fusion}(i), and injectivity from the
predecessor of \cref{base:thm:pred}. In $T_{\mathrm{fus}}$ the laws
$\mathsf{Id}$ and $\mathsf{Ext}$ are redundant: $\mathsf{Id}$ is
\cref{obs:prop:fusion}(iii) for the body $u$, whose premise $a=a$ is
reflexivity, and under the antecedent $\forall u\,(h=g)$ of an extensionality
axiom, instantiation at a fresh $z$ gives $h(z)=g(z)$, which is the premise of
\rulename{Fusion} for the identity template $f$ and the two updates $h$, $g$.
A classical structure \emph{validates} a theory if it satisfies its axioms and
every rule is valid in it in the sense stated before
\cref{disch:def:sc-model}; it then satisfies every theorem, and is a sound
test of intuitionistic derivability (\cref{obs:sec:kinds}).

\begin{theorem}[Numbering does not repair numeral induction]
\label{obs:thm:numbering-barrier}
Let $E(p,n)\defeq\bigl(\It{n}{\Suc p}{u}{\Suc u}=\Suc\It{n}{p}{u}{\Suc u}\bigr)$.
\begin{enumerate}[label=(\roman*)]
\item $\mathsf{B}$ proves $E(p,0)$ and $E(p,n)\to E(p,\Suc n)$, using only
the computation equations and equality reasoning.
\item There is a classical structure $\mathcal{D}$ with actual equality that
validates $T_{\mathrm{cal}}$ and in which $\forall n\,E(p,n)$ is false at
some value of $p$.
\item $\mathcal{D}$ has an automorphism $\tau$ of order two, of the full
staged language, that fixes every standard numeral and is not the identity.
\end{enumerate}
Hence $T_{\mathrm{cal}}$ does not prove $\forall p\,\forall n\,E(p,n)$, and
induction for literal equations is not admissible in $T_{\mathrm{cal}}$.
\end{theorem}

\begin{proof}
(i) $E(p,0)$: the zero computations give $\It{0}{\Suc p}{u}{\Suc u}=\Suc p$
and $\It{0}{p}{u}{\Suc u}=p$, and congruence gives $\Suc p=\Suc\It{0}{p}{u}{\Suc u}$.
Step: by the successor computation, the hypothesis $E(p,n)$ under $\Suc$, and
the successor computation again,
\[
\It{\Suc n}{\Suc p}{u}{\Suc u}=\Suc\It{n}{\Suc p}{u}{\Suc u}
=\Suc\Suc\It{n}{p}{u}{\Suc u}=\Suc\It{\Suc n}{p}{u}{\Suc u}.
\]

(ii) Let $M$, $D$, $\pi$, $0_D$, $\Suc_D$, $\mathscr{H}$, $\operatorname{It}_f$,
the nonstandard chains $C=(n_j)_{j\in\mathbb{Z}}$ and the fibers
$B_j=\pi^{-1}\bigl(\operatorname{It}_f(\pi(n_j),\pi(a))\bigr)$ be as in the
proof of \cref{ctr:thm:zeroind}; a \emph{path} for a triple $(h,a,C)$ is a
family $b_j\in B_j$ with $b_{j+1}=h(b_j)$. Only the choice of paths changes.
Let $\tau\colon D\to D$ fix $(m,0)$ for standard $m$ and exchange $(m,0)$ and
$(m,1)$ for nonstandard $m$. Then $\tau^{2}=\mathrm{id}$, $\pi\tau=\pi$,
$\tau(0_D)=0_D$ and $\tau\Suc_D=\Suc_D\tau$. If $h\in\mathscr{H}$ with
projected function $f$, then $\tau h\tau\in\mathscr{H}$ with the same $f$.
The label is constant along a nonstandard chain $C$, so $\tau C$ is the chain
over the same chain of $M$ with the other label, and $\tau C\ne C$.

\emph{Prescribed paths.} Fix a nonstandard $b\in M$ and put $a_0=(b,0)$,
$a_1=(b,1)$. On every nonstandard chain $C=(n_j)$ prescribe:
\begin{enumerate}[label=(\alph*)]
\item for $h=\mathrm{id}_D$ and every seed $a$, the constant path $b_j=a$;
\item for $h=\Suc_D$ and the seed $\overline{k}_D=\Suc_D^{k}0_D$, $k\in\mathbb{N}$,
the path $b_j=\Suc_D^{k}n_j$;
\item for $h=\Suc_D$ and the four seeds $a_0$, $\Suc_Da_0$, $a_1$, $\Suc_Da_1$,
the paths $b_j=(\pi(n_j)+b,0)$, $(\pi(n_j)+\Suc b,1)$, $(\pi(n_j)+b,1)$ and
$(\pi(n_j)+\Suc b,0)$ respectively, with the addition of $M$.
\end{enumerate}
These are paths. In (a) $\operatorname{It}_{\mathrm{id}}(\pi(n_j),\pi(a))=\pi(a)$
and the identity fixes $a$. In (b) and (c) the projected iteration of the
successor from $k$, $b$ or $\Suc b$ is $k+\pi(n_j)$, $b+\pi(n_j)$ or
$\Suc b+\pi(n_j)$; in (c) these first coordinates are nonstandard, so both
labels exist; and $\Suc_D$ raises the first coordinate by one and keeps the
label. The prescribed triples are pairwise distinct: the identity and the
successor are different functions, the seeds of (b) are standard, and the four
seeds of (c) are distinct and nonstandard. The involution
$(h,a,C)\mapsto(\tau h\tau,\tau a,\tau C)$ maps the set of prescribed triples
onto itself, and the prescription of the image is the image under $\tau$ of
the prescription: $\tau\,\mathrm{id}_D\,\tau=\mathrm{id}_D$ and the constant
path at $\tau a$ is $\tau$ of the constant path at $a$;
$\tau\Suc_D\tau=\Suc_D$, $\tau\overline{k}_D=\overline{k}_D$ and
$\Suc_D^{k}(\tau n_j)=\tau(\Suc_D^{k}n_j)$; and in (c) the involution
exchanges $a_0$ with $a_1$ and $\Suc_Da_0$ with $\Suc_Da_1$, whose prescribed
values have equal first coordinates and opposite labels.

\emph{The remaining paths.} The involution on triples has no fixed point,
since $\tau C\ne C$. Choose one triple $(h,a,C)$ from each two-element orbit
that is not prescribed, a path $(b_j)$ for it by \cref{ctr:lem:path}, and for
the partner $(\tau h\tau,\tau a,\tau C)$ the family that takes the value
$\tau b_j$ at $\tau n_j$. It is a path: $\tau h\tau(\tau b_j)=\tau b_{j+1}$, and
$\pi(\tau b_j)=\pi(b_j)=\operatorname{It}_f(\pi(\tau n_j),\pi(\tau a))$. Define
$\mathcal{I}(n,a,h)$ as $h^{q}(a)$ at a standard counter $\overline{q}_D$ and as
the chosen path at nonstandard counters. Then $\mathcal{I}$ satisfies the three
displayed equations of the proof of \cref{ctr:thm:zeroind}, and
\begin{equation}
\label{obs:eq:equivariant}
\mathcal{I}(\tau n,\tau a,\tau h\tau)=\tau\,\mathcal{I}(n,a,h)
\qquad(n,a\in D,\ h\in\mathscr{H}),
\end{equation}
at standard counters because $\tau\overline{q}_D=\overline{q}_D$ and
$(\tau h\tau)^{q}(\tau a)=\tau(h^{q}(a))$, and at nonstandard counters by the
construction.

\emph{Validity.} Terms are interpreted as in the proof of \cref{ctr:thm:zeroind};
\eqref{ctr:eq:termproj} holds by the same induction, every body function lies
in $\mathscr{H}$, the substitution lemma and both computation schemas hold, and
the argument given there for an equation $\overline{k}=t(n,\vec a)$, with an
arbitrary context, validates the numeral induction rule. Injectivity, coverage
and separation are shown there. The value of an iterator depends only on the
actual body function, so $\mathsf{Ext}$ holds. $\mathsf{Id}$ holds at standard
counters by finite iteration and at nonstandard ones by (a). $\mathsf{N}_k$
holds at $\overline{q}_D$ because $\Suc_D^{q}\overline{k}_D=\Suc_D^{k}\overline{q}_D$,
and at nonstandard counters by (b). Finally let $p\mapsto a_0$ and let $n$ be
nonstandard. By (c) the left side of $E(a_0,n)$ is $(\pi(n)+\Suc b,1)$, and
the right side is $\Suc_D(\pi(n)+b,0)=(\pi(n)+\Suc b,0)$.

(iii) By induction on terms, $\tau[\![t]\!]_\rho=[\![t]\!]_{\tau\circ\rho}$
for every assignment $\rho$. Variables, $0$ and $\Suc$ are immediate. For
$t=\It{n}{a}{u}{h}$ let $G$ and $G'$ be the body functions at $\rho$ and at
$\tau\circ\rho$. The induction hypothesis for $h$ at $\rho[u\mapsto v]$ gives
$G'(\tau v)=\tau G(v)$, so $G'=\tau G\tau$; the counter and the seed have the
values $\tau[\![n]\!]_\rho$ and $\tau[\![a]\!]_\rho$; and
\eqref{obs:eq:equivariant} gives the claim. So $\tau$ commutes with the
interpretation of every function symbol of the staged language. It fixes the
standard numerals and moves $a_0$.
\end{proof}

The proof of \cref{ctr:cor:noreflection} applies to $\mathcal{D}$ verbatim,
since it uses only \eqref{ctr:eq:termproj}, the one-element fiber over zero and
two distinct points with a common projection. Hence $T_{\mathrm{cal}}$ proves
for no term $e(x,y)$ the law $e(x,y)=0\leftrightarrow x=y$, and the reduction
of equation induction to zero-equation induction through such a term is not
available in $T_{\mathrm{cal}}$.

The second structure lives on pairs. Fix a nonstandard model $M$ of Peano
arithmetic expanded by symbols for all primitive recursive descriptions, write
$\mathrm{St}$ for its standard part, and put
\[
D_2\defeq\{(m,m):m\in\mathrm{St}\}\cup\{(m,m'):m,m'\in M\setminus\mathrm{St}\}
\subseteq M^{2},
\]
with $0=(0,0)$ and the coordinatewise successor. Native terms are first
interpreted on all of $M^{2}$: variables, $0$ and $\Suc$ as stated, and
$\It{n}{a}{u}{h}$ at an assignment $\rho$ as the internal iterate in $M$,
of length the first coordinate $n_1$ of $[\![n]\!]_\rho$, of the function
$v\mapsto[\![h]\!]_{\rho[u\mapsto v]}$, starting at $[\![a]\!]_\rho$. By
induction on the term every term defines, with its parameters, a function on
$M^{2}$ given by a primitive recursive description on the two coordinates, so
the internal iterate exists at every count in $M$. Write $x_{i,1}$, $x_{i,2}$
for the coordinates of a pair $x_i$.

\begin{lemma}[Provenance]
\label{obs:lem:provenance}
For every term $t(x_1,\dots,x_r)$ there are functions $\sigma$ and $d$ of
$x_{1,1},\dots,x_{r,1}$ alone, primitive recursive in $M$ with parameters,
with $\sigma\in\{0,\dots,r\}$ and $d\in M$, such that on all of $M^{2}$
\[
[\![t]\!]=\bigl(x_{\sigma,1}+d,\ x_{\sigma,2}+d\bigr),\qquad x_0\defeq(0,0).
\]
\end{lemma}

\begin{proof}
Induction on $t$. A variable $x_i$ has $\sigma=i$, $d=0$; $0$ has $\sigma=0$,
$d=0$; $\Suc$ adds $1$ to $d$. Let $t=\It{n}{a}{u}{h}$, and apply the induction
hypothesis to $a$ and to $h$, where $h$ has the additional input $u$, read as
the index $r+1$. Define internally, by primitive recursion on $j\in M$, a pair
$(\sigma_j,d_j)$: $(\sigma_0,d_0)$ is the provenance of $a$; given
$(\sigma_j,d_j)$, evaluate the provenance $(\sigma_h,d_h)$ of $h$ at the first
coordinates of the inputs and at the first coordinate $x_{\sigma_j,1}+d_j$ of
the state; if $\sigma_h=r+1$ put $(\sigma_{j+1},d_{j+1})=(\sigma_j,d_j+d_h)$,
and otherwise $(\sigma_h,d_h)$. All of this depends on first coordinates only.
Induction in $M$ on $j$ proves that the $j$-th internal iterate equals
$(x_{\sigma_j,1}+d_j,\ x_{\sigma_j,2}+d_j)$: in the first case $h$ adds $d_h$ to
both coordinates of the state, and in the second it returns
$x_{\sigma_h}+(d_h,d_h)$. Evaluate at $j=n_1$, the first coordinate of the
counter, which by the induction hypothesis for $n$ is again a function of
first coordinates.
\end{proof}

Three consequences. First, every term function maps tuples of $D_2$ into
$D_2$: adding a common $d$ to a standard diagonal pair gives a diagonal pair,
whose coordinates are both standard or both nonstandard, adding it to a pair of
nonstandard coordinates keeps both nonstandard, and $x_0$ is diagonal. Second,
for a body $h$ with parameters in $D_2$ and a seed $a\in D_2$, every internal
iterate $h^{j}(a)$, $j\in M$, lies in $D_2$, since it is the value of the
iterator at the counter $(j,j)\in D_2$. Third, by induction on $t$, the first
coordinate of $[\![t]\!]$ is the value in $M$ of $t$ at the first coordinates,
with the iterator of $M$; in the iterator case this is an induction in $M$ on
the count. Restricted to $D_2$ the interpretation is therefore a structure for
the staged language, with actual equality.

\begin{theorem}[Fusion does not give numbering]
\label{obs:thm:fusion-barrier}
The structure $D_2$ validates $T_{\mathrm{fus}}$, and $\It{n}{0}{u}{\Suc u}=n$
is false at every counter $(c,c')$ with $c\ne c'$ nonstandard. For every finite
tuple $\vec c$ of elements of $D_2$ there is an automorphism of $D_2$ of order
two, of the full staged language, that fixes $\vec c$ and is not the
identity.
\end{theorem}

\begin{proof}
\emph{Background.} The successor is injective and $0$ is not a successor. A
nonzero standard diagonal pair $(m,m)$ is the successor of $(m-1,m-1)$, and a
pair of nonstandard coordinates is the successor of the pair of their
predecessors, again nonstandard; so coverage holds. The computation equations
are the recursion equations of the internal iterate at the counts $0$ and
$n_1+1$. The interpretation is compositional on $M^{2}$, so the substitution
lemma holds by induction on terms: substituting $s$ for $x$ in a body does not
change the function on $M^{2}$ that it defines once $x$ takes the value of $s$.
$\mathsf{Id}$ holds because every internal iterate of the identity is the
identity.

\emph{Numeral induction.} The first projection maps $D_2$ onto $M$, through
the diagonal, and its fiber over a standard $k$ is $\{(k,k)\}=\{\overline{k}\}$.
By the third consequence of \cref{obs:lem:provenance}, the argument for an
equation $\overline{k}=t(n,\vec a)$ in the proof of \cref{ctr:thm:zeroind},
with the first projection in place of $\pi$, validates the numeral induction
rule with an arbitrary context.

\emph{Fusion.} Fix an assignment that satisfies $\Gamma$; it fixes the values
of the parameters of $f$, $g$, $h$ and of $t$ and $a$. Since the eigenvariable
$z$ is fresh, the premise gives $f(g(z))=h(f(z))$ for every value $z\in D_2$.
Write $a_j=g^{j}(a)$ for the internal iterates, $j\in M$. By the second consequence
of \cref{obs:lem:provenance} every $a_j$ lies in $D_2$, so
$f(g(a_j))=h(f(a_j))$ for every $j\in M$. With the primitive recursive
descriptions of $f$, $g$, $h$ and of $j\mapsto a_j$ substituted, this is a
statement of $M$ that mentions no predicate for $D_2$, and it is true in $M$.
Induction in $M$ on $j$ gives $f(a_j)=h^{j}(f(a))$: at $j=0$ both sides are
$f(a)$, and $f(a_{j+1})=f(g(a_j))=h(f(a_j))=h(h^{j}(f(a)))$. At $j$ the first
coordinate of the value of $t$, this is the conclusion
$f\bigl(\It{t}{a}{u}{g(u)}\bigr)=\It{t}{f(a)}{v}{h(v)}$. $\mathsf{Ext}$ is
derivable from \rulename{Fusion}, hence valid.

\emph{Numbering fails.} At $n=(c,c')$ the iterator $\It{n}{0}{u}{\Suc u}$ is the
$c$-fold internal successor of $(0,0)$, that is $(c,c)\ne(c,c')$.

\emph{Involutions.} For distinct positive standard $r$, $s$ define on $M^{2}$
\[
\tau_{r,s}(m,m')=
\begin{cases}
(m,m+s), & m'=m+r,\\
(m,m+r), & m'=m+s,\\
(m,m'), & \text{otherwise}.
\end{cases}
\]
The cases are disjoint, $\tau_{r,s}$ is an involution, it fixes the first
coordinate and every diagonal pair, and by cancellation in $M$ it commutes with
adding a common $d\in M$ to both coordinates. It maps $D_2$ onto $D_2$: the
standard elements of $D_2$ are diagonal and fixed, and a moved pair has a
nonstandard first coordinate, so both second coordinates are nonstandard. By
\cref{obs:lem:provenance}, applying $\tau_{r,s}$ to all arguments of a term
function changes neither $\sigma$ nor $d$, which depend on first coordinates,
so
$[\![t]\!](\tau_{r,s}\vec x)=\tau_{r,s}(x_\sigma)+(d,d)=\tau_{r,s}\bigl([\![t]\!](\vec x)\bigr)$,
also when $\sigma=0$. Hence the restriction to $D_2$ is an automorphism of the
staged language. Given $\vec c$, let $Q$ be the finite set of positive standard
$q$ such that some component of $\vec c$ has the form $(m,m+q)$, and choose
$r\ne s$ outside $Q$. Then $\tau_{r,s}$ fixes $\vec c$, and it exchanges
$(m,m+r)$ and $(m,m+s)$ for every nonstandard $m$.
\end{proof}

\begin{corollary}[Obstructions to bi-interpretation]
\label{obs:cor:barrier-biinterp}
\begin{enumerate}[label=(\roman*)]
\item No theory validated by $D_2$, in particular $T_{\mathrm{fus}}$, is
bi-interpretable with $\iPRA$ in the sense of \cref{obs:def:biinterp}, with
arbitrary definable domains, quotient equality, finite dimension and finite
parameter tuples, provided the interpretations and the return isomorphism exist
in $D_2$; this covers a parameter condition that the theory proves inhabited.
\item No theory validated by $\mathcal{D}$, in particular $T_{\mathrm{cal}}$,
is bi-interpretable with $\iPRA$ in that sense by a bi-interpretation whose
parameters, those of the first interpretation, of the return isomorphism and
representatives of those of the second interpretation, can be chosen in
$\mathcal{D}$ among the standard numerals; in particular by no
parameter-free bi-interpretation.
\end{enumerate}
\end{corollary}

\begin{proof}
As in the proof of \cref{obs:thm:symmetry}: a hypothetical bi-interpretation,
applied in the structure, yields a classical model of $\iPRA$; collect the
parameters $\vec c$ as stated. In (i) \cref{obs:thm:fusion-barrier}, and in
(ii) \cref{obs:thm:numbering-barrier}(iii), gives an automorphism of order two
that fixes $\vec c$ and is not the identity. Its image under the homomorphism
of \cref{obs:lem:aut} has order dividing two and preserves the parameter-free
definable order of \cref{obs:lem:rigid}, so it is the identity, and injectivity
of the homomorphism gives a contradiction.
\end{proof}

The automorphism $\tau$ of $\mathcal{D}$ moves every nonstandard element, so
(ii) says nothing about bi-interpretations that need a nonstandard parameter.
Neither part excludes equiconsistency or a one-way interpretation.

The witness of \cref{obs:thm:fusion-barrier} is parameter-free and its step
needs no auxiliary law, which settles the question left open in
\cref{obs:rem:fp-scope}. Let $\Nat^{-}_{\mathrm{id}}$ be $\Nat^{-}$ with the
leaves $\Gamma\proves\It{n}{a}{u}{u}=a$ for all terms $n$, $a$.

\begin{corollary}[Progress over the identity scheme]
\label{obs:cor:qi-identity}
Let $E(y)\defeq\bigl(\It{y}{0}{u}{\Suc u}=y\bigr)$.
\begin{enumerate}[label=(\roman*)]
\item $\Nat^{-}$ derives $E(0)$ and the sequent $E(z)\proves E(\Suc z)$.
\item Every sequent derivable in $\Nat^{-}_{\mathrm{id}}$, also when the rule
\rulename{Fusion} with native contexts and the rule \rulename{QI} restricted
to matrices $t(y)=\overline{k}$, $k\in\mathbb{N}$, are added, is valid in
$D_2$ under the ordinary readings, with actual equality.
\item $\Q{y}E(y)$ is false in $D_2$.
\end{enumerate}
Hence $\Nat^{-}_{\mathrm{id}}$ does not derive $\Q{y}E(y)$, and \rulename{QI}
is not admissible over $\Nat^{-}$ with the identity scheme, nor over that
extension.
\end{corollary}

\begin{proof}
(i) $E(0)$ is a zero computation leaf. In the context $E(z)$: the successor
computation gives $\It{\Suc z}{0}{u}{\Suc u}=\Suc\It{z}{0}{u}{\Suc u}$,
\rulename{Hyp} and congruence in the template $\Suc v$ give
$\Suc\It{z}{0}{u}{\Suc u}=\Suc z$, and transitivity gives $E(\Suc z)$. The
equality rules of \cref{base:lem:eq} use only \rulename{Tr} and \rulename{I0},
as in the proof of \cref{obs:cor:fp-native}.

(ii) For the seven rules of $\Nat^{-}$ the argument of part (ii) of the proof
of \cref{obs:thm:progress-independence} applies: it uses actual equality, the
computation equations, the substitution lemma and coverage, which hold in
$D_2$ by the proof of \cref{obs:thm:fusion-barrier}. The identity leaves hold
since $D_2$ satisfies $\mathsf{Id}$. \rulename{Fusion} is valid by
\cref{obs:thm:fusion-barrier}. For \rulename{QI} at a matrix $t(y)=\overline{k}$,
let $\rho$ satisfy $\Gamma$; the premise holds at every value of the fresh $z$,
so if $t(0)=\overline{k}$ holds at $\rho$ then numeral induction gives
$t(y)=\overline{k}$ at every value of $y$, which is the reading of $\Q{y}$.

(iii) $\Q{y}E(y)$ reads $E(0)\to\forall y\,E(y)$; $E(0)$ holds and $E$ fails at
every $(c,c')$ with $c\ne c'$ nonstandard.
\end{proof}

The structure $\mathcal{D}$ gives a second witness with a parameter:
$\Nat^{-}$ derives $E(p,z)\proves E(p,\Suc z)$ for the equation of
\cref{obs:thm:numbering-barrier}, by the derivation in part (i) of its proof,
and $\Nat^{-}$ with the identity leaves, the leaves
$\Gamma\proves\It{n}{\overline{k}}{u}{\Suc u}=\Suc^{k}n$ and \rulename{QI}
restricted to matrices $t(y)=\overline{k}$ is sound in $\mathcal{D}$ by the
same argument, while $\Q{y}E(p,y)$ fails at $p\mapsto a_0$.

\begin{remark}[Comparison and scope]
\label{obs:rem:barriers}
The equation $E(p,n)$ of \cref{obs:thm:numbering-barrier} is the instance of
(Shift) for the body $\Suc u$. It is also an instance of \rulename{SU}
(\cref{kern:def:SU}): for $f(z)\defeq\Suc\It{z}{p}{u}{\Suc u}$ the successor
computation and congruence give $f(\Suc z)=\Suc f(z)$, and $f(0)=\Suc p$, so
\rulename{SU} gives $\Suc\It{t}{p}{u}{\Suc u}=\It{t}{\Suc p}{u}{\Suc u}$.
Hence $\mathcal{D}$ refutes (Shift) and \rulename{SU}; it is no witness for
\cref{kern:q:su}, and it is complementary to the model of
\cref{obs:thm:fixed-points}, which satisfies (FP), (Shift) and $\mathsf{N}_0$
and is not shown to satisfy numeral induction. In \cref{ctr:thm:zeroind} the
identity iterator is the specified family $(\pi(a),0)$, so $\mathsf{Id}$ fails
there; \cref{obs:thm:numbering-barrier} keeps the base and adds $\mathsf{Id}$,
every $\mathsf{N}_k$ and the equivariance. The structure $\mathcal{M}$ of
\cref{disch:def:sc-model} also validates \rulename{Fusion} and refutes
$\mathsf{N}_0$, but its successor is not injective and coverage fails there;
$D_2$ has both, together with numeral induction. With
\cref{obs:prop:fusion,obs:thm:fixed-points} the picture over $\mathsf{B}$ is:
\rulename{Fusion} with $\mathsf{N}_0$ recovers induction for literal equations;
\rulename{Fusion}, $\mathsf{Id}$, $\mathsf{Ext}$ and numeral induction without
$\mathsf{N}_0$ do not; $\mathsf{N}_0$ with every $\mathsf{N}_k$, $\mathsf{Id}$,
$\mathsf{Ext}$ and numeral induction does not; and $\mathsf{N}_0$ with (FP),
(Shift), $\mathsf{Ext}$ and cancellative addition does not.

The package also states the positive equivalence: over $\mathsf{B}$, equation
induction proves \rulename{Fusion} and $\mathsf{N}_0$, and these two prove
\rulename{UR}, hence equation induction by \cref{kern:thm:ur}. This is
\cref{obs:prop:fusion}(i), (ii) with the discussion after it and is not
restated; the package derives \rulename{Fusion} from equation induction
directly, by induction on $f(\It{n}{a}{u}{h(u)})=\It{n}{f(a)}{u}{k(u)}$, in
place of the route through \rulename{UR}, and it notes that the converse half
of \cref{kern:thm:ur} is a written proof that is not part of the Lean
formalization. The constructions of the subsection are written proofs in
classical metamathematics with a nonstandard model of Peano arithmetic and,
for $\mathcal{D}$, the axiom of choice and \cref{ctr:lem:path}; none of these
principles enters the object theories. No Lean or Rocq verification is
claimed, and the package's test suites concern inherited checkers only.
\end{remark}

\subsection{No successor cycle without progress}

The countermodel has no successor cycle, and it cannot have one. Let $T^{-}$ be
intuitionistic logic with equality, the two computation equations and
coverage; there is no induction. Write $Z(x)\defeq\sel(x,1,0)$, so $Z(0)=1$ and
$Z(\Suc x)=0$, write $\overline{j}$ for $\Suc^{j}0$, and $h^{[r]}$ for the
$r$-fold composition of a unary term, which introduces no iterator. All
inductions below are on the integer that specifies a term.

\begin{lemma}[Finite clocks]
\label{obs:lem:clocks}
For every $m\ge 2$ there is a term $h_m(x)$, built from the leaves
$\overline{0},\dots,\overline{m-1}$ by the selector with arbitrary counter
terms, such that $T^{-}$ proves
\[
\bigvee_{j<m}h_m(x)=\overline{j},\qquad
h_m(\overline{j})=\overline{j+1}\ (j<m-1),\qquad
h_m(\overline{m-1})=0 .
\]
\end{lemma}

\begin{proof}
Let $h_2=Z$; coverage gives $Z(x)=0\vee Z(x)=1$. Given $h_m$, put
$T_m(x)\defeq\It{x}{0}{u}{h_m(u)}$. By coverage on $x$ and the two
computations, $T_m(x)$ is $0$ or a value of $h_m$, so it is one of
$\overline{0},\dots,\overline{m-1}$; nothing is asserted about the preceding
iterator value. Finite computation gives $T_m(\overline{j})=\overline{j}$ for
$j<m$ and $T_m(\overline{m})=0$. For $1\le r<m$ put
\[
d_{m,r}(z)\defeq Z\bigl(h_m^{[m-r]}(z)\bigr),\qquad
V_{m,0}(z)\defeq 0,\qquad
V_{m,r}(z)\defeq\sel\bigl(d_{m,r}(z),\,V_{m,r-1}(z),\,\overline{r+1}\bigr),
\]
and $V_m\defeq V_{m,m-1}$. Ground computation gives $d_{m,r}(\overline{s})=1$
for $s=r$ and $d_{m,r}(\overline{s})=0$ for $s<m$, $s\ne r$; hence
$V_m(0)=0$ and $V_m(\overline{r})=\overline{r+1}$ for $1\le r<m$. At an
arbitrary argument $V_m$ takes one of the values
$0,\overline{2},\dots,\overline{m}$, by coverage on the selector counters. Put
\[
h_{m+1}(x)\defeq\sel\bigl(x,\,1,\,V_m(T_m(x))\bigr).
\]
Its values are among $\overline{0},\dots,\overline{m}$. At $0$ its value is
$1$; at $\overline{j}$ with $1\le j<m$ it is
$V_m(\overline{j})=\overline{j+1}$; at $\overline{m}$ it is $V_m(0)=0$.
\end{proof}

For example $h_3(x)=\sel\bigl(x,1,\sel(Z(Z(T_2(x))),0,2)\bigr)$ with
$T_2(x)=\It{x}{0}{u}{Z(u)}$.

\begin{lemma}[Separation of numerals]
\label{obs:lem:numerals}
For $m\ge 2$ and distinct $i,j<m$, $T^{-}$ proves
$\overline{i}=\overline{j}\to 0=1$, without injectivity of the successor.
\end{lemma}

\begin{proof}
Apply $h_m^{[r]}$ with $r\equiv -i\pmod m$, $0\le r<m$. The result is
$0=\overline{s}$ with $1\le s<m$, and $Z$ gives $1=0$.
\end{proof}

\begin{theorem}[No finite successor cycle]
\label{obs:thm:no-cycle}
For every $k\ge 1$ and every term $c$, $T^{-}$ proves
$\Suc^{k}c=c\to 0=1$, and $\Nat^{-}$ derives the sequent
$\Suc^{k}c=c\proves 0=1$.
\end{theorem}

\begin{proof}
Let $m=k+1$ and $v=T_m(c)$. The hypothesis and $k$ successor computations give
$v=T_m(\Suc^{k}c)=h_m^{[k]}(v)$. By \cref{obs:lem:clocks}, $v=\overline{j}$ for
some $j<m$, and the equation becomes
$\overline{j}=\overline{(j+k)\bmod m}$. Since $k=m-1$ the two residues differ,
and \cref{obs:lem:numerals} gives $0=1$.

For $\Nat^{-}$ the finite disjunctions are removed. The cases rule for equation
templates, from $E(0)$ and $E(\Suc z)$ to $E(t)$, is derivable in $\Nat^{-}$:
\rulename{DI} by cases with an unused witness, \rulename{DE} and \rulename{QE}
(\Lean{NoInd.cases}). First, if $E(\ell)$ is derivable for each leaf $\ell$ of
a term built from finitely many leaves by the selector, then $E$ of the term is
derivable, by induction on the term with the cases rule at each selector
counter. Second, since $h_m$ is such a term over
$\overline{0},\dots,\overline{m-1}$, the cases rule at the counter of $T_m$
gives: if $E(\overline{j})$ is derivable for every $j<m$, then $E(T_m(c))$ is
derivable; the zero branch is $E(0)$ and the successor branch is
$E(h_m(T_m(z)))$. Third, put
$b_r(a)\defeq Z\bigl(h_m^{[(-r)\bmod m]}(a)\bigr)$ for $r<m$, so that
$b_r(\overline{j})=1$ exactly when $r=j$, and
\[
\operatorname{xor}(p,q)\defeq\sel(p,q,Z(q)),\qquad
\operatorname{or}(p,q)\defeq\sel(p,q,1),
\]
and let $D_m(a,b)$ be the left-nested $\operatorname{or}$, starting at $0$, of
the terms $\operatorname{xor}(b_r(a),b_r(b))$, $r<m$. The cases rule for the
template $\operatorname{xor}(Z(w),Z(w))=0$, instantiated at the arguments of
$Z$ in the terms $b_r(a)$, and the computation $\operatorname{or}(0,0)=0$ give
$D_m(a,a)=0$ for every $a$.
Ground computation gives
$D_m(\overline{j},h_m^{[k]}(\overline{j}))=1$ for $j<m$, so the second step
gives $D_m(v,h_m^{[k]}(v))=1$ at $v=T_m(c)$. Under the hypothesis
$v=h_m^{[k]}(v)$, and transport gives $D_m(v,v)=1$, hence $0=1$. No
disjunction, no characteristic of equality of arbitrary numbers and no
\rulename{QI} is used.
\end{proof}

The README records the odd case in Lean (\Lean{NoInd.no\_odd\_cycle},
\Lean{NoInd.no\_succ\_fixed}) and leaves the even case open, because it needs a
counter of another period whose behavior has to be established without
induction; the clocks $h_m$ are such counters. The theorem is a constraint on
countermodels and does not make \rulename{QI} redundant.

\begin{remark}[What is proved and what is not]
\label{obs:rem:pi-scope}
\Cref{obs:thm:progress-independence} is an independence result for one rule of
one fixed kernel, with the identity translation of formulas. It is proved by a
written model construction that uses a nonprincipal ultrafilter and the axiom
of choice. It is not formalized: the README lists the package
as filed and not formalized, and the Lean readings of \Lean{NativeModels} fix
the standard successor, so they cannot express the model. The result does not
bound the size of a checker, does not concern presentations in which
\rulename{QI} is fused with another rule or replaced by a weaker induction
interface, and does not show that the eight rules are a minimum: equivalent
presentations with fewer rule families exist. It does not decide whether the
uniqueness rule for the successor update is equivalent to induction for
literal equations, since the seed of \cref{obs:prop:pi-cancellative} refutes
that uniqueness instance; the refined model of \cref{obs:thm:fixed-points} is
not known to satisfy it (\cref{obs:rem:fp-scope}).
\end{remark}

\section{The progressive binder presentation}
\label{obs:sec:progressive}

This section draws on \pkg{research}. It shows that the calculus of
\cref{obs:def:progressive}, with \rulename{BoolCases}, is a complete
presentation of $\iPRA$, so that \cref{obs:thm:boolcases} isolates the one rule
that cannot be removed. Structural weakening, contraction, exchange,
substitution, renaming and assumption cut are admissible by induction on
derivations, with eigenparameters renamed first.

\begin{lemma}[Identity and equation induction, without cases]
\label{obs:lem:Hidentity}
The case-free calculus derives $t=t$, symmetry, transitivity and numerical
congruence; $\Gamma,A\proves A$ for every object formula $A$; and contextual
induction for literal equations: from $\Gamma\proves E(0)$ and
$\Gamma,E(z)\proves E(\Suc z)$ with $z$ fresh, $\Gamma\proves E(t)$.
\end{lemma}

\begin{proof}
The computation $\mathsf{R}(0,t;u,v.0)=t$, used as the major and as the minor
of \rulename{Tr} with the template $w=t$, gives $t=t$; the templates $w=s$,
$s=w$ and $f(s)=f(w)$ give the other equality rules. An available equation
$s=t$ transports $s=s$ to $s=t$. For an available $\mathsf{H}_{x,y}B$, induct
on the binder depth: \rulename{HE} supplies a fresh $a$ and
$\mathsf{P}_{y}B(a,y)$; reintroduce the formula by \rulename{HI} at the witness
$a$, where the local assumption $B(a,z)$ yields a proof of itself by the
induction hypothesis and \rulename{Step} gives $B(a,\Suc z)$. For induction,
\rulename{HI} with an unused witness gives $\mathsf{H}_{x,y}E(y)$,
\rulename{HE} supplies $\mathsf{P}_{y}E(y)$, and \rulename{Iter} gives $E(t)$.
\end{proof}

\begin{lemma}[Boolean cases and constructor cases]
\label{obs:lem:Hcases}
Put $P(n)=\mathsf{R}(n,0;u,v.u)$ and $C(n,a,b)=\mathsf{R}(n,a;u,v.b)$. The
case-free calculus derives
\[
C(\operatorname{sg}(x),0,x)=x,\qquad C(\operatorname{sg}(x),x,\Suc(Px))=x .
\]
Over the case-free calculus, \rulename{BoolCases} is equivalent to the rule
\[
\infrule{\Gamma,\ t=0\proves A\quad\Gamma,\ t=\Suc z\proves A}{\Gamma\proves A}\ \rulename{Cases},
\qquad z \text{ fresh for } \Gamma,t,A .
\]
\end{lemma}

\begin{proof}
Both identities follow by literal-equation induction; each base is a zero
computation and each successor step is a successor computation that does not
use its induction hypothesis. Hence $\operatorname{sg}(t)=0$ gives $t=0$ and
$\operatorname{sg}(t)=1$ gives $t=\Suc(Pt)$. For \rulename{Cases} from
\rulename{BoolCases}, substitute $P(t)$ for $z$ in the second premise.
Conversely \rulename{Cases} gives \rulename{BoolCases} by the equations of
$\operatorname{sg}$.
\end{proof}

The missing principle is thus the two-value range of $\operatorname{sg}$;
equation induction supplies the numerical identities.

\begin{lemma}[Binder congruence]
\label{obs:lem:Hcong}
If $A(x,y)$ and $B(x,y)$ are interderivable uniformly in $x,y$, with $x,y$ not
in the retained context, then $\mathsf{H}_{x,y}A$ and $\mathsf{H}_{x,y}B$ are
interderivable. No cases and no compound induction are used.
\end{lemma}

\begin{proof}
From $\mathsf{H}_{x,y}A$, \rulename{HE} gives $\mathsf{P}_{y}A(a,y)$. Introduce
$\mathsf{H}_{x,y}B$ at the witness $a$: convert the local $B(a,z)$ to $A(a,z)$,
apply \rulename{Step}, convert $A(a,\Suc z)$ to $B(a,\Suc z)$.
\end{proof}

\begin{definition}[Gates, selectors and macros]
\label{obs:def:Hmacros}
With bound variables renamed apart, define by recursion on formulas
\begin{align*}
G_v(s=t)&\defeq\bigl(C(v,0,s)=C(v,0,t)\bigr), &
G_v(\mathsf{H}_{x,y}A)&\defeq\mathsf{H}_{x,y}G_vA,\\
M_v(s=t,\,r=w)&\defeq\bigl(C(v,s,r)=C(v,t,w)\bigr), &
M_v(\mathsf{H}_{x,y}A,\mathsf{H}_{x,y}B)&\defeq\mathsf{H}_{x,y}M_v(A,B),\\
M_v(E,\mathsf{H}_{x,y}B)&\defeq\mathsf{H}_{x,y}M_v(G_yE,B), &
M_v(\mathsf{H}_{x,y}A,E)&\defeq\mathsf{H}_{x,y}M_v(A,G_yE),
\end{align*}
for equations $E$; the maximal binder depth decreases in each recursive call.
Let $\mathsf{T}=(0=0)$ and let $p,q,d$ be the terms with
\[
p(0)=0,\ p(\Suc n)=C(p(n),1,0);\quad
q(0)=0,\ q(\Suc n)=C(p(n),q(n),\Suc(q(n)));\quad
d(0)=0,\ d(\Suc n)=\Suc\Suc(d(n)),
\]
which are parity, half and doubling in $\mathbb{N}$. The logical macros are
\begin{align*}
A\Rightarrow_{\mathsf{H}}B&\defeq\mathsf{H}_{u,z}M_z(A,B), &
\exists_{\mathsf{H}}x\,A(x)&\defeq\mathsf{H}_{x,z}G_zA(x),\\
A\wedge_{\mathsf{H}}B&\defeq\mathsf{H}_{u,z}G_zM_{p(z)}(A,B), &
A\vee_{\mathsf{H}}B&\defeq\exists_{\mathsf{H}}x\,M_x(A,B),\\
\forall_{\mathsf{H}}x\,A(x)&\defeq\mathsf{H}_{u,z}M_{p(z)}\bigl(\mathsf{T},A(q(z))\bigr), &
\bot_{\mathsf{H}}&\defeq(0=\Suc 0),
\end{align*}
with all unused coordinates fresh. These are translations of fixed finite
syntax, not truth predicates.
\end{definition}

\begin{lemma}
\label{obs:lem:Hgates}
The calculus of \cref{obs:def:progressive} derives:
\begin{enumerate}[label=(\roman*)]
\item $G_0A$, and $G_{\Suc v}A$ is interderivable with $A$;
\item $A$ is interderivable with $\mathsf{H}_{u,y}G_yA$ for fresh $u,y$;
\item $M_0(A,B)$ is interderivable with $A$ and $M_{\Suc v}(A,B)$ with $B$;
\item $p(d(n))=0$, $p(\Suc d(n))=1$, $q(d(n))=n$ and $q(\Suc d(n))=n$.
\end{enumerate}
Items (i)--(iii) hold in the case-free calculus.
\end{lemma}

\begin{proof}
(i) By induction on the formula, with the computations of $C$ at equations and
\cref{obs:lem:Hcong} at binders; for $G_0$ at a binder, \rulename{HI} ignores
its local assumption and uses the induction hypothesis for the successor
matrix. (ii) Forward, \rulename{HI} ignores its assumption and uses
$G_{\Suc z}A$ from $A$. Backward, \rulename{HE} supplies the progress clause,
$G_0A$ is a theorem, and one \rulename{Step} gives $G_1A$, hence $A$. (iii) By
computation, (ii) and congruence. (iv) Induct on the single equation
$p(d(n))=0$: at a successor two parity updates send $0$ to $1$ and back. Then
$p(\Suc d(n))=1$ by one computation. Induct on the single equation
$q(d(n))=n$: the two half updates at $d(n)$ and $\Suc d(n)$ first keep the
value and then increment it, by the parity equations. Then
$q(\Suc d(n))=n$ by its defining update. Two single-equation inductions are
used in sequence; no induction on a conjunction or a compound formula.
\end{proof}

\begin{proposition}[All intuitionistic rules]
\label{obs:prop:Hlogic}
The calculus of \cref{obs:def:progressive} derives the introduction and
elimination rules of intuitionistic natural deduction for the macros of
\cref{obs:def:Hmacros}, with the usual eigenvariable conditions, explosion from
$\bot_{\mathsf{H}}$, and coverage
$(n=0)\vee_{\mathsf{H}}\exists_{\mathsf{H}}v\,(n=\Suc v)$. No excluded middle
for quantified formulas is used.
\end{proposition}

\begin{proof}
By \cref{obs:lem:Hcases}, \rulename{Cases} is available, and with
\cref{obs:lem:Hgates}(iii) it selects either endpoint of a selector at an
arbitrary control, also when the context mentions the control; in particular
proofs of $A$ and of $B$ give a proof of $M_t(A,B)$.

\emph{Implication.} Given $\Gamma,A\proves B$, apply \rulename{HI}; the local
assumption is $M_z(A,B)$, and \rulename{Cases} on $z$ gives $A$ and then $B$ at
$z=0$, and $B$ at $z=\Suc w$; in both branches the successor matrix is $B$. For
elimination, \rulename{HE} supplies the clause, $A$ is its zero matrix, and
\rulename{Step} at $0$ gives $B$.

\emph{Existential quantifier.} From $A(t)$, \rulename{HI} at the witness $t$
ignores the local gate. For elimination, \rulename{HE} gives a fresh $a$ and
the clause; the zero gate is a theorem and \rulename{Step} at $0$ gives
$A(a)$.

\emph{Conjunction.} If $A$ and $B$ are proved, \rulename{Cases} and the
selector endpoints prove every successor matrix. Conversely the zero gate is a
theorem, \rulename{Step} at $0$ gives the matrix at $1$, which is $B$ since
$p(1)=1$, and a second \rulename{Step} gives the matrix at $2$, which is $A$
since $p(2)=0$.

\emph{Disjunction.} The introductions use the witnesses $0$ and $1$. For
elimination, eliminate the existential at a fresh $a$ and apply
\rulename{Cases} to $a$.

\emph{Universal quantifier.} Let $\Gamma\proves A(x)$ with $x$ fresh. Then
$\mathsf{T}$ and $A(q(\Suc z))$ are proved, so \rulename{Cases} and the
selector endpoints prove the successor matrix, and \rulename{HI} ignores its
assumption. For specialization, \rulename{HE} gives the clause; the matrix at
$d(t)$ is a theorem by \cref{obs:lem:Hgates}(iv), and \rulename{Step} at
$d(t)$ gives the matrix at $\Suc d(t)$, which is $A(t)$. No compound progress
clause is iterated.

\emph{Falsity and coverage.} Under $0=1$, computation and transport give
$s=C(0,s,t)=C(1,s,t)=t$; induction on formulas then gives every object formula,
since at a binder \rulename{HI} ignores its assumption. \rulename{Cases} proves
coverage by its two introductions.
\end{proof}

\begin{theorem}[\pkg{research}]
\label{obs:thm:progressive}
Let $T_{\mathsf{H}}$ be intuitionistic first-order logic with equality in the
staged signature of $\mathsf{R}$, with the computation equations, induction for
literal equations, coverage and $0=1\to\bot$. The calculus of
\cref{obs:def:progressive} is sound in $T_{\mathsf{H}}$ under its ordinary
readings and derives the translations of all rules and axioms of
$T_{\mathsf{H}}$. For every matrix $B$,
\[
\mathsf{H}_{x,y}B(x,y)\ \dashv\vdash\
\exists_{\mathsf{H}}x\,\forall_{\mathsf{H}}y\,\bigl(B(x,y)\Rightarrow_{\mathsf{H}}B(x,\Suc y)\bigr).
\]
$T_{\mathsf{H}}$ and $\iPRA$ have the same theorems in the staged primitive
recursive language. The progressive calculus is equiconsistent and
bi-interpretable with $\iPRA$, with whole domains, actual equality and return
isomorphisms with graph $x=y$.
\end{theorem}

\begin{proof}
\emph{Soundness.} \rulename{HI}, \rulename{HE} and \rulename{Step} are logical
rules under the readings; \rulename{Iter} is equation induction under the
retained base; \rulename{BoolCases} is the range of $\operatorname{sg}$, which
follows from coverage.

\emph{Completeness.} \Cref{obs:prop:Hlogic} gives the logical rules and
coverage, \cref{obs:lem:Hidentity} gives equation induction, and separation
translates to the identity implication
$(0=1)\Rightarrow_{\mathsf{H}}(0=1)$.

\emph{The composite.} Forward, \rulename{HE} supplies a fresh $a$ and the
clause; at a fresh $y$ assume $B(a,y)$ and use \rulename{Step}; introduce the
implication, the universal and the existential macro and close \rulename{HE}.
Backward, eliminate the existential at a fresh $a$, specialize the universal
at the eigenparameter $z$ of \rulename{HI}, and eliminate the implication under
the assumption $B(a,z)$. The clause composite is the universal progress
formula. In $T_{\mathsf{H}}$ the gate, selector and clock equations prove the
intended meaning of every macro; for the universal macro, the even-to-odd
progress implications give every instance $A(n)$, and all instances of $A$ make
every progress implication true, with no induction on $A$. Induction on
formulas, with \cref{obs:lem:Hcong}, gives both formula composites.

\emph{Arithmetic.} $T_{\mathsf{H}}$ has the computation equations of ordinary
counter-aware primitive recursion, coverage, separation and equation
induction. The scalar iterator is the case of $\mathsf{R}$ whose update ignores
its counter, so the construction of a characteristic of actual equality and
the derivation of quantifier-free induction from equation induction
(\cref{base:thm:reflect}) apply to the terms of the larger language, and
$T_{\mathsf{H}}$ proves every quantifier-free induction instance.
Counter-aware recursion compiles every primitive recursive description
directly, in place of the scalar compiler. Numerical descriptions are
identified with themselves, and equation induction identifies two fixed
recursions with the same equations, so both domain isomorphisms have the graph
$x=y$.
\end{proof}

\begin{remark}
The presentation has one unrestricted introduction and pays for it with the
explicit \rulename{BoolCases} and with the distinction between \rulename{Step}
and \rulename{Iter}. The universal quantifier is recovered with one progress
step through the clock $p,q,d$. No smaller complete checker for this
presentation is established, and the presentation is not assigned the
measurements of the checkers of its package.
\end{remark}

\section{The counter-congruence form of the calculus}
\label{obs:sec:cc}

This section draws on \pkg{counter\_congruence}. General equality transport
carries an arbitrary formula or term context as rule data. Here it is recovered
from two specific operations: replacement of the complete left side of an
equation, and congruence in the counter of a recursor. The result is a
restriction of the primitive interface; it does not by itself give a shorter
checker. \Cref{obs:prop:root-vs-congruence} shows that root replacement alone
is not enough, and \cref{obs:thm:RL-independent} that counter congruence alone
is not enough.

\begin{definition}[The calculus $\mathcal{K}_{\mathrm{ctr}}$]
\label{obs:def:cc}
Terms are variables, $0$ and the successor-seeded counter-aware recursor
$\mathsf{R}_{\Suc}(n,a;u,v.h)$, with $\Suc a\defeq\mathsf{R}_{\Suc}(0,a;u,v.0)$
and $1\defeq\Suc 0$. The computation schemas are
\begin{align*}
\mathsf{R}_{\Suc}(0,a;u,v.h)&=\Suc a, \tag{R0}\\
\mathsf{R}_{\Suc}(\Suc n,a;u,v.h)&=h\subst{n}{u}\subst{\mathsf{R}_{\Suc}(n,a;u,v.h)}{v}. \tag{RS}
\end{align*}
Object formulas are $A::=(s=t)\mid\D{x}{y}A$. Contexts are finite lists of
object formulas and internal clauses $\Q{y}B$; a clause is available by
membership only, and an object formula by membership or by a proof, written
$\Gamma\mathrel{\triangleright}M$. The rules are \rulename{Hyp} for object
formulas, R0, RS,
\[
\infrule{\Gamma\mathrel{\triangleright}s=t\quad\Gamma\proves s=r}{\Gamma\proves t=r}\ \rulename{RL}
\qquad
\infrule{\Gamma\mathrel{\triangleright}s=t}{\Gamma\proves\mathsf{R}_{\Suc}(s,a;u,v.h)=\mathsf{R}_{\Suc}(t,a;u,v.h)}\ \rulename{CC}
\]
and \rulename{QE}, \rulename{DE}, \rulename{DI} by cases and
\rulename{DI-progress} as in $\Nat_{\mathrm{p}}$, with majors available in the
sense of $\triangleright$. In \rulename{CC} the initializer, the update and the
passive parameters are the same on both sides. There is no primitive
reflexivity, no general transport and no other congruence. The designated
contradiction is $0=1$.
\end{definition}

Weakening, exchange, contraction, renaming, substitution of a term for a
parameter and cut of an object assumption are admissible by induction on
derivations; substitution preserves the matching initializer and update of
\rulename{CC} and the matching root of \rulename{RL}. These are operations on
proof syntax and do not replace equal terms inside a formula.

\begin{lemma}[Equality algebra from root replacement]
\label{obs:lem:cc-eq}
Reflexivity, symmetry and transitivity are derivable from RS and
\rulename{RL}, without \rulename{CC}.
\end{lemma}

\begin{proof}
For a term $s$ put $c_s\defeq\mathsf{R}_{\Suc}(\Suc s,0;u,v.u)$; RS gives
$c_s=s$. With one copy as the major and one as the minor, \rulename{RL} gives
$s=s$. An available $s=t$ with the minor $s=s$ gives $t=s$. For $s=t$ and
$t=r$, symmetry gives $t=s$, and \rulename{RL} with this major and the minor
$t=r$ gives $s=r$. Paths of equations are composed at their roots; no proper
subterm changes.
\end{proof}

\begin{lemma}[Numerical cases]
\label{obs:lem:cc-cases}
For every object template $A(y)$, from $\Gamma\proves A(0)$ and
$\Gamma\proves A(\Suc z)$ with $z$ fresh, $\Gamma\proves A(t)$ is derivable
without either equality rule and without \rulename{DI-progress}.
\end{lemma}

\begin{proof}
Weaken the second premise by $A(0)$ and apply \rulename{DI} by cases with the
witness $0$ and an unused coordinate $x$, which gives $\D{x}{y}A(y)$.
\rulename{DE} supplies the clause $\Q{y}A(y)$, and \rulename{QE} with the first
premise gives $A(t)$.
\end{proof}

\begin{lemma}[Counter representation of a context]
\label{obs:lem:cc-rep}
Let $W(n)\defeq\mathsf{R}_{\Suc}(n,0;u,v.u)$, so that $W(0)=\Suc 0$ and
$W(\Suc n)=n$ by R0 and RS. For a numerical template $F(x)$ put
\[
H_F(n)\defeq\mathsf{R}_{\Suc}\bigl(n,F(0);u,v.\Suc(F(\Suc u))\bigr),\qquad
J_F(n)\defeq W(H_F(n)),
\]
with fresh bound variables. Then $\mathcal{K}_{\mathrm{ctr}}$ derives
$J_F(n)=F(n)$ for arbitrary $n$ in every context, by cases and without
\rulename{DI-progress}.
\end{lemma}

\begin{proof}
The update of $H_F$ does not contain $v$, so R0 and RS give literally
$H_F(0)=\Suc(F(0))$ and $H_F(\Suc z)=\Suc(F(\Suc z))$. Let $q$ be $0$ or
$\Suc z$. \rulename{CC} for the recursor $W$ gives
$W(H_F(q))=W(\Suc(F(q)))$, and RS gives $W(\Suc(F(q)))=F(q)$. Composition by
\cref{obs:lem:cc-eq} gives $J_F(q)=F(q)$. \Cref{obs:lem:cc-cases} for the
equation $J_F(y)=F(y)$ gives the claim. The only change of a proper subterm is
the one \rulename{CC} instance through $W$; replacing $H_F(q)$ by
$\Suc(F(q))$ under $W$ without it would assume the congruence that is being
derived.
\end{proof}

\begin{theorem}[All numerical congruence from counter congruence]
\label{obs:thm:cc-congruence}
For every numerical template $F(x)$, from $\Gamma\mathrel{\triangleright}s=t$
the calculus $\mathcal{K}_{\mathrm{ctr}}$ derives $\Gamma\proves F(s)=F(t)$.
Hence for templates $U,V$, from $\Gamma\mathrel{\triangleright}s=t$ and
$\Gamma\proves U(s)=V(s)$ it derives $\Gamma\proves U(t)=V(t)$.
\end{theorem}

\begin{proof}
\rulename{CC} with the initializer $F(0)$ and the update $\Suc(F(\Suc u))$
gives $H_F(s)=H_F(t)$; a second \rulename{CC} for $W$ gives $J_F(s)=J_F(t)$.
\Cref{obs:lem:cc-rep} and \cref{obs:lem:cc-eq} give
$F(s)=J_F(s)=J_F(t)=F(t)$. For the second claim,
$U(t)=U(s)=V(s)=V(t)$.
\end{proof}

\begin{proposition}[Constant overhead in inference nodes]
\label{obs:prop:cc-nodes}
Count one node for each primitive inference, hypothesis lookups included, and
count repeated subproofs repeatedly. If the supplied proof of $s=t$ has $p$
nodes, the derivation of $F(s)=F(t)$ has $p+56$ nodes, independently of the
syntactic depth of $F$.
\end{proposition}

\begin{proof}
Reflexivity has three nodes; symmetry adds four nodes to its premise and
transitivity five to the sum of its premises. Each of the two proofs of
$J_F(q)=F(q)$ has eight nodes: a computation, a \rulename{CC} lifting, a second
computation and the transitivity overhead. The cases construction adds
\rulename{DE}, \rulename{DI}, \rulename{QE} and one clause lookup, so
$J_F(t)=F(t)$ has twenty nodes. The two \rulename{CC} liftings of $s=t$ have
$p+2$ nodes, the reversed left representation $20+4$, the right one $20$, and
two transitivities give $(p+2)+24+20+5+5=p+56$.
\end{proof}

The count is not a bound on bytes or running time: term and formula labels
contain copies of $F$. It is not a constant bound for the lifting to compound
formulas below.

\begin{theorem}[Transport in object formulas]
\label{obs:thm:cc-transport}
For every object template $A(w)$, from $\Gamma\mathrel{\triangleright}s=t$ and
$\Gamma\proves A(s)$ the calculus $\mathcal{K}_{\mathrm{ctr}}$ derives
$\Gamma\proves A(t)$, with no progress instance in the expansion. An available
clause $\Q{y}B(s,y)$ together with an available $s=t$ has the force of
$\Q{y}B(t,y)$: from a proof of $B(t,0)$ one obtains $B(t,r)$ for every $r$.
\end{theorem}

\begin{proof}
Induction on the number of binders of $A$, uniformly in contexts, terms and
parameters. An equation is \cref{obs:thm:cc-congruence}. Let
$A(w)=\D{x}{y}B(w,x,y)$ with the bound variables renamed apart. Eliminate
$A(s)$ at a fresh $a$, which retains $\Q{y}B(s,a,y)$. Introduce
$\D{x}{y}B(t,x,y)$ by cases with the witness $a$ and a fresh $z$: the local
assumption $B(t,a,0)$ is transported backward along $t=s$ by the induction
hypothesis to $B(s,a,0)$; the clause at $\Suc z$ gives $B(s,a,\Suc z)$; the
induction hypothesis transports it forward to $B(t,a,\Suc z)$. Close the
introduction and the elimination. Both recursive uses concern the smaller
matrix $B$, and the only introductions added are case introductions. For the
clause, transport the base backward, specialize, and transport forward.
\end{proof}

\begin{corollary}[Identity-language equivalence]
\label{obs:cor:cc-equiv}
Let $\mathcal{K}_{\mathrm{tr}}$ be $\mathcal{K}_{\mathrm{ctr}}$ with
\rulename{RL} and \rulename{CC} replaced by unrestricted transport in object
formulas. For every context $\Gamma$ of object formulas and clauses and every
object formula $A$, $\Gamma\proves A$ in $\mathcal{K}_{\mathrm{ctr}}$ if and
only if $\Gamma\proves A$ in $\mathcal{K}_{\mathrm{tr}}$. The translations are
effective and preserve end sequents literally.
\end{corollary}

\begin{proof}
From $\mathcal{K}_{\mathrm{tr}}$, replace each transport by
\cref{obs:thm:cc-transport}. Conversely \rulename{RL} is transport with the
template $w=r$; reflexivity follows from a computation by the same instance;
and \rulename{CC} is obtained by transporting the left occurrence of $s$ in
$\mathsf{R}_{\Suc}(s,a;u,v.h)=\mathsf{R}_{\Suc}(s,a;u,v.h)$ to $t$ and applying
symmetry.
\end{proof}

Since the language and the meaning of every primitive are unchanged, the
bi-interpretation of $\mathcal{K}_{\mathrm{tr}}$ with $\iPRA$, with whole
domains, actual equality and identity return isomorphisms, transfers to
$\mathcal{K}_{\mathrm{ctr}}$. In the successor-seeded basis the numerical
bridge is short: with $C(n,a,b)\defeq W(\mathsf{R}_{\Suc}(n,a;u,v.\Suc b))$
and $P(n)\defeq C(n,0,W(n))$ one has $C(0,a,b)=a$, $C(\Suc n,a,b)=b$,
$P(0)=0$, $P(\Suc n)=n$, and for terms $g$ and $h$ the term
$W\bigl(\mathsf{R}_{\Suc}(n,g(\vec p);u,v.\Suc h(u,W(v),\vec p))\bigr)$
satisfies the equations of the primitive recursion with base $g$ and step $h$,
without predecessor search, pairing or induction. The package attributes this
compiler, the logical macros and the lifting of \cref{obs:thm:cc-transport} to
earlier drafts; its own contribution is \cref{obs:lem:cc-rep} and
\cref{obs:thm:cc-congruence}, with the computation schemas unchanged and the
counter as the only congruence position. An earlier draft on fixed-position
transport uses replacement in the initializer of a scalar iterator together
with an additional first-step computation schema.

\begin{theorem}[Root replacement is not derivable]
\label{obs:thm:RL-independent}
Let $\mathcal{K}^{-}$ be $\mathcal{K}_{\mathrm{ctr}}$ without \rulename{RL},
with no reflexivity or other equality rule added. Then $\mathcal{K}^{-}$ does
not derive $0=0$ from the empty context. In particular \rulename{RL} is not
admissible in $\mathcal{K}^{-}$.
\end{theorem}

\begin{proof}
Interpret the atom $s=t$ by a binary predicate, not assumed to be equality.
The domain is $\{0,1\}$, $0$ denotes $0$, and every recursion term denotes $1$
under every assignment, whatever its counter, initializer, update and
parameters; this is compatible with substitution and renaming. So $\Suc a$ is
constantly $1$. Let $s=t$ be true exactly when $s$ has the value $1$; the right
side is immaterial. Read $\Q{y}B$ and $\D{x}{y}B$ by their ordinary readings.

R0, RS and the conclusion of \rulename{CC} have a recursion term on the left
and are true. \rulename{Hyp}, availability, \rulename{QE} and \rulename{DE}
are sound as logical rules. For \rulename{DI} by cases let $b$ be the value of
the witness; a sound premise gives $B(b,0)\to B(b,1)$, because $\Suc z$ denotes
$1$, and on the two-element domain this is $B(b,0)\to\forall y\,B(b,y)$. For
\rulename{DI-progress}, with parameters and witness fixed, a term $U(y)$ has
one of the value functions $y\mapsto 0$, $y\mapsto 1$, $y\mapsto y$, so the
truth set of a literal equation $U(y)=V(y)$ is $\varnothing$, $\{0,1\}$ or
$\{1\}$; in each case truth at $0$ implies truth everywhere, and the closure
holds for every witness. So every rule of $\mathcal{K}^{-}$ is sound, and
$0=0$ is false. In $\mathcal{K}_{\mathrm{ctr}}$ two copies of
$\mathsf{R}_{\Suc}(\Suc 0,0;u,v.u)=0$ as major and minor of \rulename{RL} prove
$0=0$.
\end{proof}

The theorem concerns the specified inference system. The interpretation is not
a model of an ordinary theory whose equality carries its usual axioms, and the
theorem does not exclude a bi-interpretation with a different equality or
domain.

\section{Constructed sources and certificate interfaces}
\label{obs:sec:sources}

This section draws on \pkg{constructed\_computation} and
\pkg{uniform\_predecessor} for the computation interface, and on
\pkg{checker\_and\_barriers}, \pkg{compact\_kernel}, \pkg{local\_scopes} and
\pkg{research} for statements about their checkers. The object calculus is
$\Nat_{\mathrm{p}}$ with \rulename{Comp} throughout; the results change the
certificate format, not the theory.

\subsection{Successor sources constructed from predecessors}

Earlier certificate formats submit an iterator source $[c,a,h]$ and must check
that the counter $c$ has successor shape before using a successor computation;
deleting that check admits a false equation. The constructed-source formats
submit the predecessor.

\begin{definition}
\label{obs:def:packets}
A \emph{constructed computation packet} supplies a predecessor $p$, an
initializer $a$, an update $h$ and a mode. In base mode the checker generates
$\It{0}{a}{v}{h}=a$. In step mode it constructs the successor itself and
generates
\[
\It{\Suc p}{a}{v}{h}=h\subst{\It{p}{a}{v}{h}}{v}.
\]
In \pkg{constructed\_computation} the mode is an explicit flag in a previously
unused packet field, and the generated equation is opened at the fresh
parameter of the continuation. In \pkg{uniform\_predecessor} the mode is read
off the raw representation of $p$ (an empty representation selects the base
equation), the packet is marked by a null major, and the generated equation is
opened with the image zero. In both, the recursive image
$\It{p}{a}{v}{h}$ is validated as a numerical term only when opening the update
requests it.
\end{definition}

\begin{lemma}[Known-instance discharge]
\label{obs:lem:known-instance}
If $w\notin\FV(\Gamma,A)$, $\Gamma\proves H(0)$ and $\Gamma,H(w)\proves A$,
then $\Gamma\proves A$.
\end{lemma}

\begin{proof}
Rename the eigenparameters of the second derivation, substitute $0$ for $w$,
which leaves $\Gamma$ and $A$ unchanged, and cut with the first derivation.
\end{proof}

\begin{lemma}[Boxed zero]
\label{obs:lem:boxing}
Put $B_0(p)\defeq\It{0}{p}{w}{0}$. The base computation and the successor
computations at predecessors of the form $B_0(p)$ derive every instance of the
successor computation.
\end{lemma}

\begin{proof}
The base computation gives $B_0(p)=p$. Compute at $B_0(p)$ and transport along
$B_0(p)=p$ with the template
$\It{\Suc w}{a}{u}{h}=h\subst{\It{w}{a}{u}{h}}{u}$. The argument holds for
every element of the domain.
\end{proof}

\begin{theorem}[Constructed computation supply]
\label{obs:thm:constructed-supply}
\begin{enumerate}[label=(\roman*)]
\item In the format of \pkg{constructed\_computation}, every successfully
generated equation $H(w)$ has a zero instance $H(0)$ derivable from the
computation equations and transport. Hence its continuation is admissible by
\cref{obs:lem:known-instance}.
\item In the format of \pkg{uniform\_predecessor}, every successfully supplied
equation is a theorem uniformly in every newly admitted parameter; no
known-instance discharge is needed.
\item Both formats derive exactly the closed theorems of $\Nat_{\mathrm{p}}$
with \rulename{Comp}. In the first, the step at the predecessor $0$ is the
packet with $p=0$ in step mode. In the second, a literal $0$ selects the base
equation, and the step at $0$ is obtained by \cref{obs:lem:boxing}.
\end{enumerate}
\end{theorem}

\begin{proof}
Validation separates into two obligations: the source counter is a successor,
and a recursive image used in the result is a valid term. Construction
establishes the first, and demand-driven substitution the second.

(i) In base mode, after $w=0$ the initializer and its copy on the right agree
by structural zero completion, and the equation is the base computation. In
step mode, suppose first that opening the update requests the recursive image.
The image $[p,a,h]$ is then validated at numerical root; successful validation
is idempotent, so validating on demand gives the same substituted value as
validating eagerly. The additional dummy slot below the constructed successor
introduces no dependency, since root validation of the image rejects a
meaningful reference to it. The two occurrences of the predecessor agree and
the zero instance is the successor computation. If the update never requests
its image, it is independent of the state on every meaningful path, and after
zero completion the equation is $\It{\Suc p_0}{a_0}{v}{h_0}=h_0$ with $h_0$
independent of $v$, which is a successor computation without the unused image.
The old goal and context were checked before $w$ was allocated and do not
contain it.

(ii) With a null major the continuation opens the equation with the image zero,
and the newly admitted parameters stay rigid. For a predecessor with empty
representation the decoded counter is zero and the two occurrences of the
initializer agree, so the equation is a base computation. Otherwise the source
counter is the constructed successor, a successor whatever the value of the
predecessor. If the update uses its state, the image is checked at the old
parameter floor and the opening lemma of the inherited codec identifies the
right side with the update at the preceding iterator. If not, the update is
independent of the state, and the successor computation with a
state-independent update holds for every predecessor and initializer, including
new parameters.

(iii) Soundness is (i), (ii) and the unchanged logical branches. For
completeness, encode each instance of \rulename{Comp} by its packet and close
its use by \rulename{Hyp}; the other instructions are unchanged.
\end{proof}

The raw codec is inherited. Input is a JSON array that ends with the closed
goal; four-field packets are consumed from the end; membership in the context
closes a goal without a packet. One list grammar is read in two sorts, with
arities standing for truth or zero, closure or successor, equation, and paired
binder or iterator. Negative coordinates name allocated rigid parameters and
positive coordinates name protected binder slots. An opener substitutes a
checked image at exposed numerical coordinates. These conventions are syntax,
not arithmetic operations.

The uniform format shortens certificates: a step at a scalar parameter is one
record where the earlier flat format needs three (boxing, computation,
transport). For the same named derivations and the delivered encoders the
record counts are as follows; they are not lower bounds on shortest proofs.

\begin{center}
\begin{tabular}{lrr}
\toprule
Named derivation & Earlier format & Uniform predecessor\\
\midrule
Nonvacuous zero witness & 4 & 2\\
Atomic progress closure & 6 & 4\\
Arbitrary-term specialization & 8 & 6\\
Contextual induction with seed congruence & 48 & 40\\
\bottomrule
\end{tabular}
\end{center}

\subsection{Statements about the checkers of the packages}

\emph{Composition of the arity codec with predecessor packets}
(\pkg{checker\_and\_barriers}, \pkg{compact\_kernel}). The arity codec encodes
$s=t$ as $[s,t]$, $\D{x}{y}B$ as $[B]$ and $\Q{y}B$ as $[B,B,B]$, with the
depth increment $5-\mathrm{len}$ for every array, numerical root depth three,
and fixed truth and zero aliases at the other arities. The composed checker
uses the packet $[\mathtt{null},p,a,h]$ of \cref{obs:def:packets}. For a
truthy $p$ it corresponds to the arity packet $[[\,],[p,a,h],1,0]$ and for a
falsy $p$ to $[[\,],[0,a,h],p,0]$: the generated raw equation, the image
checks at the old floor, the next floor and the opened goal and hypothesis are
identical in the two executions, so the inherited supply proof applies. The
literal $0$ stays available as a witness of an introduction, and the step at
the predecessor $0$ is the three-node derivation of \cref{obs:lem:boxing}.

\begin{theorem}[Composed verifier]
\label{obs:thm:composed}
Every successful execution of each of the supplied composed checkers has a
decoded conclusion derivable in $\Nat_{\mathrm{p}}$, and every closed theorem
of $\Nat_{\mathrm{p}}$ has a finite accepted canonical certificate, given
sufficient resources.
\end{theorem}

\begin{proof}
Induction on the successful trace. Membership is \rulename{Hyp}, by raw
equality and the opening invariant: at a meaningful sorted root a successful
opening implements capture-free filling of the distinguished hole, by
induction on the tree, and the interpretation of a constant alias is unchanged
by traversal of its children. A proved equality or closure major with matching
endpoints is transport or \rulename{QE}; a mismatch leaves the goal and adds
truth. Triplication of a unary major is \rulename{DE} at a fresh parameter. The
positional argument check makes an introduction a genuine instance, with the
restriction of progress to literal equations. Arithmetic supply is
\cref{obs:thm:constructed-supply}. Parameters admitted inside a completed
major and unused by its end sequent are replaced by $0$. For completeness,
encode a derivation rule by rule, with the boxed step at the predecessor $0$.
\end{proof}

The invariant is conditional: opening is not idempotent on every raw tree
below ignored annotations. The two packages obtain the same 612-byte program.

\emph{Local eigenparameter scopes} (\pkg{local\_scopes}). The inherited checker
decreases one global integer at every allocation, so the proof of a major
permanently advances the allocator of the continuation. The local-scope checker
passes the current floor to each call by value.

\begin{theorem}[Local scope preserves the calculus]
\label{obs:thm:local-scope}
With the same sorted decoding and rule correspondence, the local-scope checker
is sound for $\Nat_{\mathrm{p}}$ with \rulename{Comp}, and every finite
canonical derivation has a local-scope certificate.
\end{theorem}

\begin{proof}
At the entry of a call with floor $k$ and context $g$ the scope contains the
labels $k,\dots,-1$, and $k=-|g|$, since each continuation appends one context
entry and decreases the floor by one. Induct on a successful trace. The entry
traversal checks the goal at the current scope before a name is allocated. A
major is checked at the same scope, so its conclusion contains no name
allocated inside its subproof, and context extension builds a new list, so its
local assumptions do not escape. The caller then allocates $k-1$; the
discharged side proof imposes no freshness requirement on this independent
binding. The openings, the transport, the closure use, the paired rules and the
generated arithmetic assumption are justified as before; the generated
assumption needs only a derivable zero instance, by
\cref{obs:lem:known-instance}. For completeness, rename the binders of a named
derivation, label the active eigenparameters at a branch with $d$ context
entries by members of $-1,\dots,-d$, serialize a major with the same
environment and discard its private extensions on return.
\end{proof}

The theorem states preservation of the calculus under scope-respecting
renaming, not equality of accepted byte strings: two sibling subproofs may
reuse a raw label, and an old packet may mention a name from a finished side
proof.

\emph{Iterative continuations} (\pkg{research}). The loop checker keeps
recursion for majors and runs continuations in a loop over an immutable
context.

\begin{lemma}[Continuation simulation]
\label{obs:lem:continuation}
On canonical certificates and with unbounded resources, the iterative checker
has the same successful derivations as the recursive flat checker, with the
same order of allocation of fresh names.
\end{lemma}

\begin{proof}
The invariant is that the current goal, the contents of the context, the
global freshness floor and the unconsumed instruction stack equal those of the
corresponding recursive call. Hypothesis closure is the same membership test.
The recursive call for a major uses the same preceding context, and its
allocations persist in both programs. Every branch constructs the same next
goal and added assumption; tuple concatenation has the contents of list
concatenation and changes no earlier context. Opening the goal and opening the
assumption have no effect on the floor, so their order is immaterial.
Induction on the trace gives both directions.
\end{proof}

Persistence of the context is necessary by \cref{obs:ex:leak}. The loop removes
the dependence of recursion depth on the number of continuation steps, not on
nested majors, syntax traversal or parsing.

\section{Partial routes}
\label{obs:sec:partial}

This section draws on the second research note of \pkg{scope\_kernel}, on
\pkg{local\_scopes} and on the appendix on stationarity of the second write-up
of \pkg{progress\_independence}. All three results are conditional: each
isolates the statement that is missing for a complete replacement of a rule.

\subsection{From uniqueness of iteration toward induction}

Work in intuitionistic logic with equality over the scalar iterator, with the
two computation equations. The \emph{uniqueness rule} is
\[
\infrule{\Gamma\proves f(\Suc z)=h(f(z))}{\Gamma\proves f(t)=\It{t}{f(0)}{u}{h(u)}}\ \rulename{UR},
\]
where $f$ and $h$ are fixed term templates, $z$ is fresh for $\Gamma$, $t$ and
the parameters, and $h(u)$ does not contain $z$. The premise has no discharged
induction hypothesis and there is no separate base premise. Uniqueness of
recursion as an arithmetic principle is classical for richer languages with
counter-aware recursion and product types; the question here is the scalar
update.

\begin{lemma}
\label{obs:lem:ur-basic}
\begin{enumerate}[label=(\roman*)]
\item Induction for literal equations derives \rulename{UR}, without coverage
or an equality test.
\item \rulename{UR} derives: from $f(0)=g(0)$, $f(\Suc z)=h(f(z))$ and
$g(\Suc z)=h(g(z))$ with $z$ fresh, $f(t)=g(t)$.
\item \rulename{UR} derives $\It{t}{0}{u}{\Suc u}=t$.
\end{enumerate}
\end{lemma}

\begin{proof}
(i) Induct on $f(n)=\It{n}{f(0)}{u}{h(u)}$: the base is the zero computation,
and $f(\Suc n)=h(f(n))=h(\It{n}{f(0)}{u}{h(u)})=\It{\Suc n}{f(0)}{u}{h(u)}$.
(ii) Apply \rulename{UR} to $f$ and to $g$ and use congruence in the
initializer. (iii) Take $f(z)=z$ and $h(u)=\Suc u$; the premise is
reflexivity.
\end{proof}

Part (i) is formalized for the kernel as \Lean{NPrf.uniq}.

\begin{lemma}[Retraction lemma]
\label{obs:lem:retraction}
Let $e(x)$ be a term template such that $\Gamma$ proves, with $x$ fresh,
\[
e(0)=0,\qquad e(e(x))=e(x),\qquad e(\Suc(e(x)))=\Suc(e(x)).
\]
Then the computation equations, equality reasoning and \rulename{UR} prove
$\Gamma\proves e(t)=t$. No induction is used.
\end{lemma}

\begin{proof}
Put $h(u)=e(\Suc(e(u)))$ and $r(n)=\It{n}{0}{u}{h(u)}$. Idempotence at
$\Suc(e(u))$ gives $e(h(u))=h(u)$, and idempotence under $\Suc$ and $e$ gives
$h(e(u))=h(u)$; the third hypothesis is $h(u)=\Suc(e(u))$. For
$f(n)=e(r(n))$: $f(0)=e(0)=0$ and
$f(\Suc n)=e(h(r(n)))=h(r(n))=h(e(r(n)))=h(f(n))$. \rulename{UR} and
replacement of the equal initializers give $e(r(n))=r(n)$; so the whole orbit
of $h$ lies in the image of the retraction. Then
$r(\Suc n)=h(r(n))=\Suc(e(r(n)))=\Suc(r(n))$ and $r(0)=0$, so \rulename{UR} for
$r$ with the update $\Suc u$ and \cref{obs:lem:ur-basic}(iii) give $r(n)=n$.
Hence $e(n)=e(r(n))=r(n)=n$.
\end{proof}

\begin{proposition}[Induction from a term-defined split]
\label{obs:prop:ur-split}
Let $E(n)$ be a formula with $\Gamma\proves E(0)$ and
$\Gamma,E(z)\proves E(\Suc z)$, $z$ fresh. Suppose there is a term $e(n)$ such
that, uniformly in a fresh $n$, $\Gamma\proves E(n)\vee\neg E(n)$,
$\Gamma,E(n)\proves e(n)=n$ and $\Gamma,\neg E(n)\proves e(n)=0$. Then
\rulename{UR} proves $\Gamma\proves E(t)$.
\end{proposition}

\begin{proof}
A case split proves $E(e(n))$: in the first branch by $e(n)=n$, in the second
by $e(n)=0$ and the base. The true-branch identity at $e(n)$ gives
$e(e(n))=e(n)$. The step at $e(n)$ gives $E(\Suc(e(n)))$, and the true-branch
identity gives $e(\Suc(e(n)))=\Suc(e(n))$. The base gives $e(0)=0$.
\Cref{obs:lem:retraction} gives $e(n)=n$, and transport of $E(e(n))$ gives
$E(n)$.
\end{proof}

\begin{corollary}[Induction for zero equations]
\label{obs:cor:ur-zero}
With coverage and separation, \rulename{UR} derives: from
$\Gamma\proves q(0)=0$ and $\Gamma,q(z)=0\proves q(\Suc z)=0$,
$\Gamma\proves q(t)=0$. Only coverage, separation, the selector computations,
equality reasoning and \rulename{UR} are used.
\end{corollary}

\begin{proof}
Coverage gives the case distinction for every term, and $\Suc x=0$ is
contradictory: the template $v\mapsto\sel(v,0,\Suc 0)$ gives $\Suc 0=0$. So
$q(n)=0$ is decidable. Put $e(n)=\sel(q(n),n,0)$. If $q(n)=0$ then $e(n)=n$;
otherwise coverage for $q(n)$ leaves the successor alternative, where
$e(n)=0$. Apply \cref{obs:prop:ur-split}.
\end{proof}

Coverage decides whether one number is zero; it does not give a term
$\delta(a,b)$ with $\delta(a,b)=0\leftrightarrow a=b$. If the arithmetic with
\rulename{UR} constructed such a term and proved the equivalence on its whole
domain, \cref{obs:cor:ur-zero} applied to $\delta(s(n),t(n))=0$ would recover
induction for literal equations, and with \cref{obs:lem:ur-basic}(i) the two
principles would be equivalent over the common background. The note does not
supply this: the available proofs of predecessor and of reflection of equality
(\cref{base:thm:pred,base:thm:reflect}) use literal-equation inductions with
non-constant right sides, and importing them would be circular. The note
records one step of a translation. With $Z(u)=\sel(u,1,0)$,
$N(n)=\It{n}{0}{u}{Z(u)}$, $a+b=\It{b}{a}{u}{\Suc u}$, $j(u)=Z(u)+2N(u)$,
$M(n)=\It{n}{0}{u}{j(u)}$ and $G(n)=\It{n}{0}{u}{u+M(u)+1}$, the comparison
$M(G(n))=N(n)$ can be set up as a comparison of two recurrences with the
common update $k(v)=\It{\Suc v}{v}{u}{j(u)}$, and finite Boolean computations
give $k(N(n))=Z(N(n))$. The README of the repository records that the package
\pkg{scalar\_uniqueness} derives predecessor, reflection of equality and
equation induction from \rulename{UR} and so closes this gap; that argument is
not part of the note. A native kernel with \rulename{UR} in place of the
progress rule would in addition have to derive its equality logic, case
distinction and cut without the removed rule.

\subsection{What fixed countdown uniqueness determines}

A draft on scalar countdown uses a bouncing predecessor $W$ and a countdown
$L$ with $L(0,a)=a$ and $L(\Suc n,a)=W(L(n,a))$, and the contextual rule
\[
\infrule{\Gamma\proves g(\Suc z,\vec p)=W(g(z,\vec p))}{\Gamma\proves g(t,\vec p)=L(t,g(0,\vec p))}\ (U_W),
\]
with $g$ a fixed template and $z$ fresh. It completes its theory with an
additional faithfulness sentence.

\begin{lemma}[\pkg{local\_scopes}]
\label{obs:lem:action}
If $\Gamma\proves f(0)=0$ and $\Gamma\proves f(\Suc z)=\Suc(f(z))$ with $z$
fresh, then $\Gamma\proves L(f(t),a)=L(t,a)$ for every $t$ and $a$, by $U_W$
alone. Consequently, if the theory also proves
\[
\forall x\,\forall y\,\bigl((\forall a\,L(x,a)=L(y,a))\to x=y\bigr),
\]
it derives $f(t)=t$ for every such $f$.
\end{lemma}

\begin{proof}
Take $z$ fresh for $a$ and $g(n)=L(f(n),a)$. Then
$g(\Suc z)=L(\Suc(f(z)),a)=W(L(f(z),a))=W(g(z))$ and $g(0)=L(0,a)=a$. Apply
$U_W$ and substitute $g(0)=a$. The second claim instantiates the displayed
sentence at $x=f(t)$, $y=t$.
\end{proof}

For example $L(A(t,0),a)=L(t,a)$ for an initialized successor iterator $A$.
Equality of all countdown observations is a statement about an action and is
not equality of the counters. No proof of the displayed faithfulness sentence
from $U_W$ is given, and no equivalence with the stronger single-observation
sentence $\forall a\,\forall b\,(L(b,\Suc\Suc a)=2\to a=b)$ of the draft is
claimed.

\subsection{Stationarity over an explicit primitive recursive interface}

The \emph{stationarity rule} is
\[
\infrule{\Gamma\proves f(\Suc z,\vec p)=f(z,\vec p)}{\Gamma\proves f(t,\vec p)=f(0,\vec p)}\ \rulename{Stat},
\]
for a term $f$ of the language below, with $z$ fresh for $\Gamma$, $t$ and
$\vec p$. It has no induction hypothesis, and its conclusion mentions no
iterator. Read in the native language over the computation equations, it is
equivalent to the instance of \rulename{UR} for the identity update,
$\Gamma\proves f(t)=\It{t}{f(0)}{u}{u}$, together with the scheme
$\It{t}{a}{u}{u}=a$; the scheme is \rulename{Stat} for
$f(z)=\It{z}{a}{u}{u}$, whose premise is a successor computation. The scheme
fails at $t=c$, $a=0$ in the countermodel of
\cref{obs:thm:progress-independence}, so \rulename{Stat} is not sound there.

The \emph{interface} is intuitionistic first-order logic with equality over
the terms of primitive recursive arithmetic, with the defining equations of
all primitive recursive function symbols of fixed description, coverage
$x=0\vee\exists y\,x=\Suc y$, separation $0\ne 1$, and, for each
quantifier-free formula $A(\vec x)$, a fixed primitive recursive term
$\chi_A$ with the truth-flag axioms
\[
\chi_A(\vec x)=0\leftrightarrow A(\vec x),\qquad
\chi_A(\vec x)=0\vee\chi_A(\vec x)=1 .
\]
All interface axioms are theorems of $\iPRA$. They are not claimed to be
derivable from the scalar computation equations without induction.

\begin{proposition}[Stationarity recovers quantifier-free induction]
\label{obs:prop:stationary}
Over the interface, \rulename{Stat} and induction for quantifier-free formulas
derive each other, with arbitrary contexts. Hence the interface with
\rulename{Stat} has the same theorems as $\iPRA$ in its definitional expansion
by the symbols $\chi_A$, with the identity numerical domain and actual
equality.
\end{proposition}

\begin{proof}
Let $A$ be quantifier-free with $\Gamma\proves A(0)$ and
$\Gamma,A(z)\proves A(\Suc z)$, $z$ fresh; the parameters of $A$ are kept and
suppressed. Write $e(n)\defeq\chi_A(n)$ and take primitive recursive symbols
with
\begin{gather*}
\kappa(0)=0,\quad \kappa(\Suc x)=1,\qquad C(0,a,b)=a,\quad C(\Suc x,a,b)=b,\\
\beta(0)=0,\quad \beta(\Suc n)=C\bigl(e(\Suc n),\,\kappa(\beta(n)),\,1\bigr).
\end{gather*}
Coverage at $v$ and the equations of the clamp $\kappa$ give
$\kappa(v)=0\vee\kappa(v)=1$.
For a fresh $n$, a case split on $n$ by coverage, with no induction, proves
\[
\beta(n)=0\vee\beta(n)=1,\qquad \beta(n)=0\to A(n).
\tag{$*$}
\]
At $n=0$, $\beta(0)=0$ and the base gives $A(0)$. At $n=\Suc m$ split on the
flag $e(\Suc m)$. If $e(\Suc m)=0$, then $\beta(\Suc m)=\kappa(\beta(m))$, which is
$0$ or $1$, and $A(\Suc m)$ holds by the truth-flag axiom. If
$e(\Suc m)=1$, then $\beta(\Suc m)=1$, and $\beta(\Suc m)=0$ gives $0=1$,
hence $A(\Suc m)$ by separation. Only the outermost recursion equation of $\beta$ is used.

The premise of \rulename{Stat}, $\beta(\Suc n)=\beta(n)$ for fresh $n$,
follows by cases on the first part of $(*)$. If $\beta(n)=1$, then
$\kappa(\beta(n))=1$, and
both flag cases give $\beta(\Suc n)=1$. If $\beta(n)=0$, the second part of
$(*)$ gives $A(n)$, the step gives $A(\Suc n)$, so $e(\Suc n)=0$ and
$\beta(\Suc n)=\kappa(0)=0$. \rulename{Stat} gives $\beta(t)=\beta(0)=0$, and $(*)$
at $t$ gives $A(t)$.

Conversely, induction on the equation $f(n,\vec p)=f(0,\vec p)$ derives
\rulename{Stat}: the base is reflexivity, and the step composes the premise
$f(\Suc z,\vec p)=f(z,\vec p)$ with the hypothesis
$f(z,\vec p)=f(0,\vec p)$. Since the interface axioms are $\iPRA$ theorems, the
two theories prove the same formulas.
\end{proof}

The proposition is an alternative induction principle, not a smaller checker.
Installing the interface inside the scalar kernel without \rulename{QI} would
need predecessor and reflection of equality, and their proofs
(\cref{base:thm:pred,base:thm:reflect}) use \rulename{QI}, so citing them would
be circular. The rule \rulename{UR} of \pkg{scalar\_uniqueness} has an
arbitrary update and does not by itself justify the restriction to the
identity update.

\section{Checker sizes}
\label{obs:sec:sizes}

\Cref{obs:tab:sizes} lists the complete standalone programs of the packages of
this chapter with the sizes that the packages state, under the protocol of
\cref{base:sec:measure}. The sizes of \pkg{scope\_kernel},
\pkg{progress\_independence} and \pkg{ipra\_fixed\_points\_and\_checker}, whose
241-token compact checker is a trace-equivalent reimplementation of an inherited
calculus, belong to their checker work and are listed with it in other
chapters, the last in \cref{prov:ch:cert}.
\pkg{iteration\_barriers} delivers no new program; it
re-measures five unchanged programs of \pkg{whole\_json\_refinement} and
\pkg{shared\_leaf} under CPython 3.12.14, with the compact proof interface of
\cref{prov:ch:cert} as input contract (a promised exact-integer and list
alphabet and a canonical closed guarded public goal, not a validator of
arbitrary JSON values). Its values agree with \cref{cert:sec:sizes}. It
reports the envelope $248/268/469$ for whole-input programs and $242/262/469$
when JSON-line streams are admitted, each attained by three different
programs. These envelopes are stated for the programs the package inspects;
\cref{cert:sec:sizes} also records a stream program of
\pkg{ipra\_fixed\_points\_and\_checker} with $241$ tokens and later programs
with smaller envelopes. The observed minima at the
snapshots of the packages are: 277 tokens, 311 nodes and 672 bytes, attained by
different programs, for \pkg{adaptive\_seed}, \pkg{research},
\pkg{local\_scopes} and \pkg{uniform\_predecessor}; 268, 300 and 616 for
\pkg{checker\_and\_barriers} and \pkg{compact\_kernel}; 281, 315 and 685 for
\pkg{constructed\_computation} and \pkg{counter\_congruence}. These minima are
attained sizes, hence upper bounds on an unknown optimum; they are not lower
bounds.

\begingroup\small
\begin{longtable}{@{}p{0.33\linewidth}rrrp{0.40\linewidth}@{}}
\caption{Checker sizes as stated by the packages. The control column gives
tokens / nodes / bytes of the matched control.}\label{obs:tab:sizes}\\
\toprule
Program & Tokens & Nodes & Bytes & Matched control\\
\midrule
\endfirsthead
\toprule
Program & Tokens & Nodes & Bytes & Matched control\\
\midrule
\endhead
\bottomrule
\endfoot
\multicolumn{5}{@{}l}{\pkg{adaptive\_seed}}\\
rejected local-scope prototype & 279 & 322 & 685 & 277 / 317 / 684\\
\midrule
\multicolumn{5}{@{}l}{\pkg{research}}\\
iterative checker & 276 & 305 & 749 & 278 / 314 / 679 (flat computation)\\
byte profile & 280 & 319 & 652 & same\\
explanatory checker & 303 & 339 & 1401 & ---\\
\midrule
\multicolumn{5}{@{}l}{\pkg{local\_scopes}}\\
local-scope checker & 284 & 323 & 668 & 281 / 322 / 672 (zero-product control)\\
readable checker & 320 & 363 & 1426 & ---\\
\midrule
\multicolumn{5}{@{}l}{\pkg{checker\_and\_barriers}}\\
explicit list dispatch & 266 & 310 & 612 & 265 / 312 / 616 (old arity, explicit lists, same driver)\\
compact list predicate & 266 & 312 & 604 & 265 / 314 / 608 (old arity, compact lists, same driver)\\
Boolean-expression profile & 267 & 315 & 593 & 274 / 318 / 601 (old affine, same driver)\\
guarded traversal, tuple context & 269 & 300 & 690 & see below\\
\midrule
\multicolumn{5}{@{}l}{\pkg{compact\_kernel}}\\
selected composition & 266 & 310 & 612 & 269 / 306 / 628 (matched old-flat control)\\
shorter-byte composition & 267 & 315 & 597 & ---\\
smaller-AST composition & 267 & 298 & 639 & ---\\
immutable-context loop & 266 & 296 & 689 & ---\\
\midrule
\multicolumn{5}{@{}l}{\pkg{uniform\_predecessor}}\\
selected checker & 269 & 307 & 629 & 269 / 306 / 628 (flat computation, matched driver)\\
predecessor packet, two verdicts & 278 & 315 & 680 & 278 / 314 / 679 (flat computation)\\
ancillary byte profile & 269 & 311 & 616 & see below\\
ancillary AST profile & 268 & 300 & 649 & see below\\
readable checker & 312 & 359 & 1531 & ---\\
\midrule
\multicolumn{5}{@{}l}{\pkg{constructed\_computation}}\\
selected profile & 279 & 314 & 682 & 283 / 315 / 683 (old generator, same initialization profile)\\
explicit assignments & 277 & 317 & 684 & 281 / 318 / 685 (generated assumptions, principal)\\
AST profile & 279 & 314 & 709 & 283 / 315 / 710 (generated assumptions, AST profile)\\
\midrule
\multicolumn{5}{@{}l}{\pkg{counter\_congruence}}\\
readable checker of $\mathcal{K}_{\mathrm{ctr}}$ & 1617 & 1822 & 6834 & ---\\
module import only & 281 & 320 & 680 & 281 / 318 / 685 (frozen principal baseline)\\
module import, exact class & 280 & 321 & 680 & same\\
module import, exact type & 281 & 322 & 676 & same\\
module import, zero-product test & 281 & 322 & 672 & same\\
inverted hypothesis guard & 283 & 314 & 751 & 283 / 315 / 710 (frozen AST baseline)\\
compact, separate \rulename{RL}, \rulename{CC}, \rulename{QE} & 343 & 395 & 881 & 296 / 333 / 704 (general transport)\\
compact, \rulename{CC} fused into local facts & 341 & 388 & 892 & same\\
\midrule
\multicolumn{5}{@{}l}{\pkg{iteration\_barriers} (unchanged inherited programs, re-measured)}\\
whole input, token profile (\pkg{whole\_json\_refinement}) & 248 & 277 & 615 & --- (no new program)\\
whole input, AST profile (\pkg{whole\_json\_refinement}) & 254 & 268 & 613 & --- (no new program)\\
whole input, byte profile (\pkg{shared\_leaf}) & 272 & 306 & 469 & --- (no new program)\\
JSON-line stream, token profile (\pkg{shared\_leaf}) & 242 & 268 & 565 & --- (no new program)\\
JSON-line stream, AST profile (\pkg{shared\_leaf}) & 246 & 262 & 556 & --- (no new program)\\
\end{longtable}
\endgroup

For \pkg{checker\_and\_barriers} the further same-driver controls are the old
arity token profile, 264 / 303 / 641, and the old research loop,
267 / 297 / 698; a screen of 105 old profiles under the same rejection driver
gives the matched minima 264 tokens, 297 nodes and 601 bytes, of which only the
byte minimum is improved, by the 593-byte profile. For
\pkg{uniform\_predecessor} the matched-driver controls are: token leader
268 / 309 / 633, AST leader 271 / 303 / 654, byte leader 272 / 314 / 621, flat
computation 269 / 306 / 628, predecessor joint profile 271 / 306 / 637; the
selected checker is one node and one byte larger than the matched flat checker,
and its position below the unchanged minima comes from rejection by exit status
in place of a printed verdict. In \pkg{constructed\_computation} the arithmetic
generator alone changes the explicit program from 281 / 318 / 685 to
277 / 317 / 684. The 612-byte programs of \pkg{checker\_and\_barriers} and
\pkg{compact\_kernel} are the same file.

The implementation devices are inherited source idioms and are not arithmetic
principles. In short: the list test \texttt{a * 0 == []} on parsed JSON values;
a short-circuit expression for guarded elimination; exception-based openers,
which reject nonempty strings only through the recursion limit of the runtime
and are therefore not compared with structurally terminating parsers; folding
of initialization into a default argument; module import in place of a name
import; and rejection by nonzero exit status. All programs import only the
standard JSON library, and none imports a producer or a theorem library.

\section{Further questions}
\label{obs:sec:questions}

The following statements of the packages are not results. First, what was not
achieved.

\begin{itemize}
\item No package of this chapter obtains a complete axiomatization that is
simpler than the common core under the measurement protocol. Where a checker
improves a coordinate of the observed minima, the gain is an implementation
change to an inherited theory. No delivered program is below all three
coordinates in \pkg{research}, \pkg{local\_scopes},
\pkg{checker\_and\_barriers}, \pkg{compact\_kernel} or
\pkg{counter\_congruence}; \pkg{constructed\_computation} is below its three
minima by two tokens, one node and three bytes; \pkg{uniform\_predecessor} is
below the unchanged minima only through the rejection convention.
\item No obstruction of the chapter is a lower bound on the size of a checker,
and none excludes equiconsistency or all one-way interpretations. Apart from
\cref{obs:prop:monotone}, \cref{obs:lem:ur-basic}(i), the strictness
statement cited after \cref{obs:prop:forced-progress} and the odd case of
\cref{obs:thm:no-cycle}, the README of the repository lists none of them as
formalized in Lean. In particular the countermodel of
\cref{obs:thm:progress-independence} and the finite clocks of
\cref{obs:lem:clocks} are written proofs only.
\item \pkg{progress\_independence}: the result is the independence of one rule
of the fixed kernel. No smaller arithmetic basis follows from it, and the
package does not decide whether uniqueness of iteration for the successor
update is equivalent to induction for literal equations; its strengthened
countermodel refutes that uniqueness instance and so gives no answer. The second
write-up likewise claims no smaller axiomatization, no new checker-size record
and no global minimality of the eight rules, and \cref{obs:prop:stationary}
gives no smaller native checker, since the interface it assumes is not derived
in the scalar kernel without \rulename{QI}.
\item \pkg{ipra\_fixed\_points\_and\_checker}: \cref{obs:thm:fixed-points} is a
separation of principles and deletes no arithmetic axiom; it is not a smaller
axiomatization, not a lower bound on checker size, and does not exclude a
bi-interpretation through a different definable domain. The package does not
show that \rulename{QI} fails to be admissible over the seven rules with the
identity scheme alone (\cref{obs:rem:fp-scope}); this is settled by
\cref{obs:cor:qi-identity} of \pkg{iteration\_barriers}.
It does not decide the uniqueness rule for the successor update in its model.
\item \pkg{iteration\_barriers}: no axiomatization below the
observed minima, no new checker and no minimality statement is obtained; the
positive equivalence of \rulename{Fusion} with $\It{n}{0}{u}{\Suc u}=n$ and
equation induction is \cref{obs:prop:fusion} and rests on the written proof of
\cref{kern:thm:ur}. \Cref{obs:cor:barrier-biinterp}(ii) does not address
bi-interpretations with nonstandard parameters, and neither part excludes
equiconsistency or one-way interpretations. \Cref{obs:thm:numbering-barrier}
concerns the identity law for the designated iterator, not the stronger rule
that every term satisfying a constant-step recurrence is constant, and it
settles no restricted uniqueness scheme. The package also states, without a
model or a derivation, that scope validation of inserted terms in the compact
openers cannot be deleted, that packing a checker into a decoded string or a
universal interpreter only relocates the trusted program, and that
\rulename{QI} for arbitrary formulas changes the target theory; these are
assessments of implementation proposals, not theorems. The countermodels are
not formalized in Lean or Rocq. Its finite check of the
seed algebra is supplementary; no Lean or Rocq verification of the
construction is claimed.
\item No complete checker is given for the progressive calculus of
\cref{obs:sec:progressive}. No implementation of $\mathcal{K}_{\mathrm{ctr}}$
is below the inherited programs; the compact attempts of the package have no
universal audit of their raw aliases.
\item \pkg{adaptive\_seed}: the deletion of derived majors is rejected only in
the form of \cref{obs:prop:majors}. For a redesigned calculus without public
closure goals, logical normalization does not justify deleting all derived
majors, because equation induction can produce an equality that is then used
for transport into a compound formula, and at an arbitrary parameter that
induction is not a reducible successor case. This is an open proof obligation,
not an incompleteness theorem.
\item \pkg{compact\_kernel}: two directions were examined without result.
Replacing computation supply by numerical definitional conversion requires a
reducer that respects dummy slots, the state binder, neutral counters and rigid
parameters, and gave no smaller checker. Binders in the certificate grammar,
with eigenparameters allocated by the opener, require a new proof-opening
discipline, and no complete smaller construction was obtained.
\item \pkg{uniform\_predecessor}: binary-tree domains and arithmetic whose
behavior depends on the syntax of a body were examined; no reduction of the
complete checking mechanism and no valid smaller presentation was found.
\item \pkg{scope\_kernel}: the route from \rulename{UR} to induction is
unfinished in the note (\cref{obs:sec:partial}); the note claims no
bi-interpretation and no smaller checker for a calculus with \rulename{UR}.
\item \pkg{validated\_leaves}: \cref{obs:thm:adaptive} is a written proof; no
Lean or Rocq verification of it is claimed. It is a term-definability
obstruction and does not exclude interpretations of $\iPRA$ whose operations
are defined by formulas or finite-case graphs. The package deletes no
arithmetic axiom, and its checkers do not go below the compact minima.
\end{itemize}

\begin{question}[Relational escapes]
\Cref{obs:thm:quotient} and the theorems of \cref{obs:sec:majorants} exclude
term realizers only, and \cref{obs:thm:no-biinterp} requires an
existential-free axiomatization and finite control. Is there a seedless basis,
with positive existence axioms or with an additional primitive, that is
bi-interpretable with $\iPRA$ through operations defined by graphs, and whose
complete checker, with the defining equations and scope rules of the addition
counted, is smaller? Under which witness-extraction property do graph
interpretations reduce to the term case?
\end{question}

\begin{question}[Leaving the almost unary class]
An independent initializer is one way to combine two unbounded inputs. Is there
another operation, such as pairing or addition, that restores completeness of a
seedless basis more cheaply, with both interpretation composites proved
internally? \Cref{obs:prop:selector-repair} gives numerical completeness for an
unbounded selector and no more.
\end{question}

\begin{question}[Extent of finite control]
Does finite control, and with it \cref{obs:thm:no-biinterp}, survive the
adjunction of bounded numerical primitives other than an equality
characteristic? Is every unary primitive recursive function denoted by a
tied-seed term? Neither is decided by \cref{obs:thm:fc}.
\end{question}

\begin{question}[Restricted access to the state]
\Cref{obs:thm:adaptive} answers the question of a common barrier for
\cref{obs:thm:single-body,obs:thm:circuit}: adaptive queries to one body in
seed and update, register vectors and counter-dependent locations do not leave
the almost unary class. It is broken at once by unrestricted numerical
retention of an old register (\cref{obs:rem:boundary}(iv)). Which restricted
retention interfaces suffice to define addition and a selector, support their
internal proofs, and give both interpretation round trips? Defining the
functions in the standard model is not enough (\pkg{validated\_leaves}).
\end{question}

\begin{question}[Weaker case principles]
Is there a primitive rule weaker in syntax than \rulename{BoolCases} that
supplies the mixed-parameter elimination refuted by the separator of
\cref{obs:thm:boolcases}, supports both formula round trips and admits a
smaller checker? Is there a restriction of numerical cases that avoids the
product model of \cref{obs:thm:direct-target} and still derives cases for all
native formulas? A proposal must say where the global case distinction enters;
predecessor identities do not supply it.
\end{question}

\begin{question}[Positive iteration]
Is there a small operation, cheaper than the split into base and step, that
destroys the zero-profile invariant of \cref{obs:thm:zeroprofile} and recovers
a selector and the logical bootstrap? Can initialized iteration be introduced
relationally from positive iteration without assuming the selector-based logic
that it is meant to reconstruct, and how does induction transfer through the
graph definitions when new updates are nested?
\end{question}

\begin{question}[Majors and cuts]
Can atomic cuts be compiled away without induction on compound formulas, so
that a focused calculus with fewer proved majors is complete? A normalization
theorem would have to handle equation-induction proofs used for transport in
arbitrary formulas, or an explicit generator of local induction clauses would
have to be counted in the checker.
\end{question}

\begin{question}[Uniqueness and countdown]
Does a complete checker for a calculus with \rulename{UR} improve the size
coordinates? Removing an induction hypothesis does not decide this, since the
conclusion of \rulename{UR} constructs several term instances and an iterator.
Does $U_W$ alone prove the faithfulness sentence of \cref{obs:lem:action}, so
that the additional observation axiom of the countdown draft can be removed?
\end{question}

\begin{question}[The countermodel to progress]
Can \cref{obs:lem:pi-extension} and the tower of \cref{obs:prop:pi-tower} be
formalized in the Lean development, with arbitrary carriers and an interpreted
successor in place of readings over the natural numbers? The same is asked for
the orbit-map extension of \cref{obs:lem:fp-extension}
(\pkg{ipra\_fixed\_points\_and\_checker}). Which parts of the
extension lemma survive when further laws of iteration, such as fusion or
uniqueness, are imposed on the coherent structures? Computation and
extensionality in the update are compatible with the countermodel; uniqueness
for the identity update excludes its particular witness, and excluding one
model does not give arithmetic strength. This part is answered in part by
\cref{obs:sec:fixed-points}: the extension survives preservation of all fixed
points and seed shift for every update (\cref{obs:lem:fp-extension}), and the
seed can carry the canonical successor iterator and commutative cancellative
addition (\cref{obs:prop:fp-tower}); it cannot survive general fusion together
with the canonical successor iterator, because the two give \rulename{UR} and
hence equation induction (\cref{obs:prop:fusion} and the discussion after it).
General fusion alone is compatible with a failure of equation induction: the
pair structure of \cref{obs:thm:fusion-barrier}, which is not a tower, refutes
the canonical successor iterator instead. Whether the tower itself can be
extended so as to keep all instances of \rulename{Fusion} is not decided. Is
there a counterpart of the model
that satisfies the uniqueness rule for the successor update and still refutes
an induction for a literal equation, which would separate the two principles?
The model of \cref{obs:thm:fixed-points} satisfies $\It{n}{0}{u}{\Suc u}=n$,
which is the instance of that rule at the template $z$, but it is not shown to
satisfy the rule, so the question stays open.
Is there a presentation that combines computational equality with an induction
interface weaker than \rulename{QI} without moving the cost into syntax
validation?
\end{question}

\begin{question}[After the countermodel to progress]
The second write-up of \pkg{progress\_independence} asks three further
questions. Is there an explicitly specified, for instance computable, algebra
that supports all native substitutions and refutes $\It{t}{0}{u}{u}=0$ at some
$t$? It would have to solve the problem of two-sided orbits for every update
term, including updates that use earlier nonstandard iterators.
\pkg{ipra\_fixed\_points\_and\_checker} asks the same for
\cref{obs:thm:fixed-points}: is there a direct model, without a nonprincipal
ultrafilter and an elementary tower, that keeps all passive-parameter and
native-substitution laws together with (FP) and (Shift)? The second question,
whether some induction instance for a literal equation remains underivable
after the scheme $\It{n}{a}{u}{u}=a$ is added to the seven retained rules, is
answered by \cref{obs:cor:fp-native} in the following exact form: the
conclusion $\Q{y}E(y)$ of the instance for
$E(n)\equiv\bigl(\It{n}{0}{u}{\Suc(\Suc u)}=\It{n}{n}{u}{\Suc u}\bigr)$ is
underivable even after the rule \rulename{FP}, which contains the scheme, and
the leaves \rulename{Shift} are added, and over these additions its premise
$E(z)\proves E(\Suc z)$ is derivable, so \rulename{QI} is not admissible there.
The stronger form, whether \rulename{QI} is admissible over $\Nat^{-}$ with the
identity scheme, is answered negatively by \cref{obs:cor:qi-identity}
(\pkg{iteration\_barriers}): for
$E(y)\equiv\bigl(\It{y}{0}{u}{\Suc u}=y\bigr)$ the base and the step are
derivable in $\Nat^{-}$ alone, and $\Q{y}E(y)$ is not derivable even after the
identity leaves, \rulename{Fusion} and \rulename{QI} for matrices
$t(y)=\overline{k}$ are added. After the identity scheme and every
$\It{n}{\overline{k}}{u}{\Suc u}=\Suc^{k}n$ are added instead, the model of
\cref{obs:thm:numbering-barrier} gives an underivable instance whose step is
derivable in $\Nat^{-}$, with a parameter in the equation. Both answers use
nonstandard models, not explicit ones, so the first question of this paragraph
stays open for them as well. Any such scheme would have to be counted in a
complete checker. Can predecessor and reflection of equality be derived
from \rulename{Stat} of \cref{obs:prop:stationary} in the bare scalar language,
without \rulename{QI}, so that the interface of that proposition is supplied
inside the kernel? Equivalently in spirit, does \rulename{Stat} recover the
strength missing from (FP) over the scalar interface? The model of
\cref{obs:thm:fixed-points} satisfies the identity scheme, which is the instance
of \rulename{Stat} at $f(z)\defeq\It{z}{a}{u}{u}$, but it is not shown to
satisfy \rulename{Stat}, so it does not settle this
(\pkg{ipra\_fixed\_points\_and\_checker}).
\end{question}

\begin{question}[Restricted fusion]
Which proper fragment of \rulename{Fusion} recovers induction for literal
equations over computation, coverage and separation? By
\cref{obs:prop:fusion} the whole rule does so with the canonical successor
iterator, and (FP) and (Shift) do not (\cref{obs:thm:fixed-points}). The
doubling instance excludes the model of \cref{obs:thm:fixed-points}, but
excluding one witness is not completeness.
(\pkg{ipra\_fixed\_points\_and\_checker}.) The scheme
$\It{t}{0}{u}{\Suc u}=t$ cannot be dropped from \cref{obs:prop:fusion}(ii):
the whole rule \rulename{Fusion}, with the identity law, extensionality,
successor injectivity and induction against every fixed numeral, does not
recover induction (\cref{obs:thm:fusion-barrier}); conversely the scheme with
every $\It{n}{\overline{k}}{u}{\Suc u}=\Suc^{k}n$, the identity law,
extensionality and numeral induction does not either
(\cref{obs:thm:numbering-barrier}). The question therefore concerns fragments
of \rulename{Fusion} taken together with the scheme
(\pkg{iteration\_barriers}).
\end{question}

\begin{question}[Parameters for the numbering theory]
Is $T_{\mathrm{cal}}$ of \cref{obs:sec:barriers} bi-interpretable with $\iPRA$
by a bi-interpretation that uses nonstandard parameters? The involution of
\cref{obs:thm:numbering-barrier} moves every nonstandard element, so
\cref{obs:cor:barrier-biinterp}(ii) does not decide this; a model of
$T_{\mathrm{cal}}$ with involutions fixing arbitrary finite tuples, as for
$T_{\mathrm{fus}}$, would. Do $T_{\mathrm{cal}}$ and $T_{\mathrm{fus}}$
interpret $\iPRA$ at all, and are they equiconsistent with it? Both are
subtheories of $\Teq$, so they are interpreted in $\iPRA$; the converse is not
addressed (\pkg{iteration\_barriers}).
\end{question}

\begin{question}[Counter congruence and the reversed binder]
Can \rulename{RL}, \rulename{CC} and numerical computation be encoded through
one checked representation that deletes a trusted operation, so that the
restricted interface of \cref{obs:sec:cc} removes code? Can a staged
introduction with delayed witness choice repair the reversed binder of
\cref{obs:prop:reversed}, with the reverse formula composite proved for nested
formulas?
\end{question}

\begin{question}[Formalization]
The packages name as formalization targets: the bi-interpretation of the native
calculus with its ordinary realization, with formula composites, freshness and
arbitrary contexts; the scalar predecessor and the reflection of equality,
without reintroducing quantifier-free induction before the characteristic is
proved correct; the derivation transformer for $J_F=F$ of
\cref{obs:lem:cc-rep}, without use of the ambient equality substitution; the
sorted codec with its opener, freshness conditions and input contract; and a
proved decoder that connects the arity checker with the Lean semantics of the
native core. None is supplied.
\end{question}

\begin{question}[Metric]
A coordinatewise envelope is not one program and not a lower bound. Would a
fixed implementation grammar, a fixed semantic contract, a fixed treatment of
the host runtime and a policy on aliases make strict-improvement claims more
informative? Triplicated closures make explicitly serialized certificates grow
as $3^r$ in closure depth, and the protocol charges checker source only; how
should certificate size and checking time enter the comparison? The empirical
target left by \pkg{checker\_and\_barriers} is one structurally terminating
complete program below 593 bytes, 264 tokens and 297 nodes under matched
drivers.
\end{question}
